\documentclass[11pt]{article}
\usepackage[a4paper,margin=25mm]{geometry}
\usepackage{amsmath,amssymb,amsthm,mathtools}
\usepackage{mathrsfs}
\usepackage[T1]{fontenc}
\usepackage{hyperref}
\emergencystretch=2em
\newtheorem{theorem}{Theorem}[section]
\newtheorem{proposition}[theorem]{Proposition}
\newtheorem{lemma}[theorem]{Lemma}
\newtheorem{corollary}[theorem]{Corollary}
\theoremstyle{definition}\newtheorem{definition}[theorem]{Definition}
\theoremstyle{remark}\newtheorem{remark}[theorem]{Remark}
\DeclareMathOperator{\Div}{Div}
\DeclareMathOperator{\lcm}{lcm}
\newcommand{\Z}{\mathbb Z}\newcommand{\Q}{\mathbb Q}
\newcommand{\N}{\mathbb N}
\title{Exact bounded congruence gluing and arithmetic successors\\
for Erd\H{o}s--Straus shells}
\author{The Clankers}
\date{8 September 2026}
\begin{document}
\maketitle
\begin{center}
\small Arithmetic contributor: Reddit user \texttt{u/UmbrellaCorp\_HR}.\\
Expanded edition: 72 statements with full proofs.\\
Project DOI: \href{https://doi.org/10.5281/zenodo.20401937}{10.5281/zenodo.20401937}.
\end{center}
\noindent This edition adds fourteen statements to the 58-statement snapshot
\cite{PriorEdition}. The canonical mathematical source files are unchanged;
this reading edition adds contribution credit and citations at the relevant
source interfaces. Exact computation supplements the proofs at stated finite
bounds. References distinguish a contribution's mathematical source from
broader historical context. The full announced Navier--Stokes construction
remains outside this article's independent verification claim.
\par\medskip


We study positive integer solutions of $4/p=1/x+1/y+1/z$ through
their complete bounded divisor-shell coordinates. The starting arithmetic,
positivity-code and prime-production constructions were contributed by Reddit
user \texttt{u/UmbrellaCorp\_HR} and developed further in this collaboration
\cite{UmbrellaContribution}. Their relation to the established Egyptian-fraction
parametrizations of Christian Elsholtz and Terence Tao is recorded explicitly;
the elementary reconstructions needed in this article are proved here
\cite{ElsholtzTao}. The cross-shell calculations classify ratio-preserving
transport and two elementary positive unimodular shears. The gluing
construction retains the finite exponent bounds throughout its integral lifts.
The divisor-category connection uses Alain Connes and Caterina Consani's
embedded ordered groups and their exact reciprocal-divisor locus
\cite{CCPericyclic}. A later part returns exact covering and inverse-operator
calculations to the auxiliary torus of the announced Navier--Stokes manuscript
\cite{OpenAINS}. No universal arithmetic occupancy, empty prime shell, or
complete fluid blowup proof is asserted by these calculations.

\section{Original shells, coordinates, and integer reconstruction}

Fix \(h\in\Z_{>0}\) such that \(p=12h+1\) is prime.  First choose
\[
 a\in A_p:=\{3h+1,\ldots,9h\},
\]
and then define \(R_a=4a-p\), \(S_a=pa\).  In particular
\[
 3\le R_a\le2p-3,\quad R_a\equiv3\pmod4,\quad 0<a<p.
\]
The order of these definitions is part of the construction.  The
identity \(4a=p+R_a\) will be used without changing that order.

\begin{lemma}\label{lem:units}
One has \(\gcd(R_a,S_a)=1\).  The finite set
\[
 \mathcal W_a=
 \{(m,n)\in\Z_{>0}^2:m,n\mid S_a,\ \gcd(m,n)=1,\
                         R_a\mid m+n\}
\]
is in bijection with all ordered positive integer solutions whose first
denominator is \(a\).  With \(k=(m+n)/R_a\), the bijection is
\begin{equation}\label{eq:reconstruct}
 (m,n)\longmapsto(a,y,z)
       =\bigl(a,(S_a/n)k,(S_a/m)k\bigr).
\end{equation}
Its inverse is the unique reduced positive pair \(m,n\) in
\[
 \frac mn=\frac{R_a y-S_a}{S_a}=\frac yz .
\]
Complementation interchanges \(m,n\) and interchanges \(y,z\).
\end{lemma}
\begin{proof}
Since \(0<a<p\), primality gives \(\gcd(a,p)=1\), hence
\(\gcd(a,4a-p)=1\).  A common divisor of \(p,4a-p\) would divide
\(4a\), which is prime to \(p\).  This proves the unit assertion.
For \((m,n)\in\mathcal W_a\), the integer \(k\) is positive and
\(S_a/n,S_a/m\) are positive integers.  Thus
\[
 \frac1y+\frac1z=\frac{n+m}{S_a k}=\frac{R_a}{S_a}
                   =\frac4p-\frac1a.
\]
Conversely the original equality implies
\[
 R_a yz=S_a(y+z),\qquad
 d:=R_a y-S_a=\frac{S_a y}{z}>0,\qquad
 (R_a y-S_a)(R_a z-S_a)=S_a^2.
\]
Thus \(d\) is a positive divisor of \(S_a^2\) and
\(d\equiv-S_a\pmod{R_a}\).  If \(v_\ell(S_a)=b_\ell\) and
\(v_\ell(d)=c_\ell\), then \(0\le c_\ell\le2b_\ell\).
The reduced ratio \(d/S_a\) has numerator and denominator exponents
\[
 v_\ell(m)=\max(c_\ell-b_\ell,0),\qquad
 v_\ell(n)=\max(b_\ell-c_\ell,0).
\]
Both are at most \(b_\ell\), and at most one is nonzero.  Hence \(m,n\)
divide \(S_a\) and are coprime.  Multiplication of
\(d\equiv-S_a\) by the unit \(n/S_a\) gives \(m+n\equiv0\).
The reverse formula \(d=S_a m/n\) is integral and divides \(S_a^2\)
by those same exponent inequalities.  Substituting it in
\(y=(S_a+d)/R_a\) and the complementary formula for \(z\)
gives \eqref{eq:reconstruct}.  These constructions are inverse by
uniqueness of the reduced fraction, with no reordered denominator.
\end{proof}

Define the full supply of primitive rational ratios by
\[
 \mathcal P(N):=\{m/n:m,n\in\Div(N),\ \gcd(m,n)=1\}.
\]
Every ratio retains its unique positive reduced numerator and
denominator.  If \(N=\prod_{\ell\mid N}\ell^{b_\ell}\), then
\begin{equation}\label{eq:signed-box}
 \prod_{\ell\mid N}([-b_\ell,b_\ell]\cap\Z)
 \xrightarrow{\;\sim\;}\mathcal P(N),\qquad
 \boldsymbol\beta\longmapsto\prod_{\ell\mid N}\ell^{\beta_\ell}
\end{equation}
is a bijection.  The inverse is \(v_\ell(m)-v_\ell(n)\);
primitivity gives the positive and negative parts just computed.
This is a bijection of finite sets, not the replacement of the set
by the subgroup of \(\Q_{>0}^{\times}\) that it generates.

For direct comparison with the two supplied channels, put
\[
 E_a=\{u\mid a^2:R_a\mid4u+1\},\qquad
 M_a=\{u\mid a^2:R_a\mid u+a\}.
\]
All divisors in this text are positive.  To recover their precise
multiplicities, write \(d\mid p^2a^2\) uniquely as \(d=p^i v\),
where \(i=0,1,2\), \(v\mid a^2\).  When \(i=2\), the congruence
\(d\equiv-pa\) is \(R_a\mid pv+a\), and multiplication by
\(4p^{-1}\) modulo \(R_a\), using \(4a=p\), makes it
\(R_a\mid4v+1\).  Complementation exchanges this case with \(i=0\).
When \(i=1\), it is \(R_a\mid v+a\), and complementation is
\(v\mapsto a^2/v\).  Its only possible fixed point \(v=a\) would
give \(R_a\mid2a\); by Lemma~\ref{lem:units} this would give
\(R_a\mid2\), impossible.  Thus the number of unordered passing
pairs is exactly \(|E_a|+|M_a|/2\).

The two direct exponent maps on
\[
 B_a=\prod_{\ell\mid a}\{0,\ldots,2v_\ell(a)\},\qquad
 U(\boldsymbol e)=\prod_{\ell\mid a}\ell^{e_\ell}
\]
therefore have respective targets
\[
 U(\boldsymbol e)\equiv -4^{-1}\pmod{R_a},
 \qquad U(\boldsymbol e)\equiv-a\pmod{R_a}.
\]
Their signed version is \(U(\boldsymbol e)/a\).  It has targets
\(-p^{-1}\) and \(-1\), respectively.  The maps and inverses are
the translations \(e_\ell\leftrightarrow e_\ell-v_\ell(a)\);
they retain every exponent bound.

\subsection{The full primewise biconditional and the two unary layers}

Let \(q_{R_a,p}\) count unordered pairs underlying
\(\mathcal W_a\), and retain the exact full aggregate
\[
 Q_p=\sum_{a=3h+1}^{9h}q_{R_a,p},\qquad
 q_{R_a,p}=|E_a|+\tfrac12|M_a|.
\]
These definitions do not identify \(Q_p\) with a restricted
selector.  There is a direct proof of the full biconditional
\[
 \operatorname{ES}(p)\quad\Longleftrightarrow\quad Q_p>0.
\]
For necessity choose a permutation of an actual integer
solution with first denominator \(a\) minimal, retaining
the permutation to recover its original order.  Positivity
gives \(4/p>1/a\), and minimality gives
\(4/p=1/a+1/y+1/z\le3/a\).  Hence
\(p/4<a\le3p/4\), whose integer values for \(p=12h+1\)
are exactly \(3h+1,\ldots,9h\).  Lemma~\ref{lem:units}
then gives a passing pair and thus a positive summand.
Conversely a positive summand has a passing pair, and
\eqref{eq:reconstruct} supplies an original positive
integer witness.  This proves both directions while
leaving every summand in the original full range.

The source's unary cutoff \(A\le2\) has a different
precise domain.  Reduce \(u/a=B/A\) with positive
coprime \(A,B\).  The signed box proves \(A,B\mid a\),
hence \(AB\mid a\).  Writing \(s=a/(AB)\) gives
\[
 a=ABs,\qquad u=B^2s,\qquad
 A=\frac{a}{\gcd(a,u)},\quad
 B=\frac{u}{\gcd(a,u)},\quad
 s=\frac{\gcd(a,u)^2}{u}.
\]
The first two descriptions prove the integrality of
the last formula and are inverse by the unique reduced
ratio.  Multiplying the exterior congruence by the
unit \(A\), and using \(4a\equiv p\bmod R_a\),
gives exactly \(R_a\mid A+pB\).  The middle congruence
gives exactly \(R_a\mid A+B\), since \(a,A\) are
units modulo \(R_a\).  Restricting to \(A=1,2\)
therefore keeps exactly those canonical pairs and
their stated marks, not all divisors of \(a^2\).

At \(p=37,a=12,R_a=11\), the full divisors of \(144\)
are
\[
 1,2,3,4,6,8,9,12,16,18,24,36,48,72,144 .
\]
The exterior target \(8\bmod11\) is attained only by
\(u=8\); the middle target \(10\bmod11\) is attained
by none.  Here \(u/a=2/3\), so \((A,B,s)=(3,2,2)\).
The full shell has the witness
\((12,42,1036)\) and no \(A\le2\) witness.
For the reciprocal check the common denominator 3108
gives \(259+74+3=336=4\cdot84\).
This disproves per-shell completeness of the cutoff.
It does not disprove the source's global assertion
that some shell admits that cutoff: at the same prime
the shell \(a=10,R_a=3,u=5\) has \(A=2,B=1,s=5\)
and \(3\mid4\cdot5+1\).  Its reconstructed triple is
\((10,2405,130)\), obtained from
\(d=37^2\cdot5=6845\), \(S_a=370\).
Both \(y=(370+6845)/3=2405\) and
\(z=(370+20)/3=130\) are positive integers and satisfy
the original identity by Lemma~\ref{lem:units}.

\subsection{An exact factor sieve in the first two original shells}

Keep \(p=12h+1\) prime, \(h\ge1\), and put
\(a_3=3h+1\), \(a_7=3h+2\). These are the original shells
with residuals \(R_{a_3}=3\) and \(R_{a_7}=7\).
They both lie in \(A_p\), and both are at most \(6h=(p-1)/2\).
For either retained integer \(a\), write its actual prime factorization
\(a=\prod_{\ell\mid a}\ell^{v_\ell}\). The exponent domain throughout
is exactly
\[
 \mathcal B_a=\prod_{\ell\mid a}\{0,\ldots,2v_\ell\},\qquad
 u(e)=\prod_{\ell\mid a}\ell^{e_\ell}.
\]
Unique prime factorization proves this is a bijection onto \(\Div(a^2)\),
with inverse \(e_\ell=v_\ell(u)\). Every prime, exponent cap and
integer \(a\) remains attached to the coordinates below.

\begin{theorem}[The complete residual-three factor sieve]
\label{thm:r3-factor-sieve}
At \(a=a_3\), define
\[
 P_1=\prod_{\substack{\ell\mid a\\\ell\equiv1\ (3)}}(2v_\ell+1),
 \qquad
 P_2=\prod_{\substack{\ell\mid a\\\ell\equiv2\ (3)}}(2v_\ell+1).
\]
Empty products are one. The full exterior and middle sets coincide,
and their complete exponent fibres are
\[
 E_a=M_a=
 \left\{u(e):e\in\mathcal B_a,\quad
           \sum_{\ell\equiv2\ (3)}e_\ell\equiv1\pmod2\right\}.
\]
In particular
\[
 |E_a|=|M_a|=\frac{P_1(P_2-1)}2,
 \qquad |E_a|+\frac{|M_a|}{2}=\frac{3P_1(P_2-1)}4.
\]
The shell has a witness if and only if \(a\) has a prime divisor
\(\ell\equiv2\pmod3\). Every such factor gives both channels
with primitive layer \(A=1\), via
\[
 u=a\ell,\quad A=1,\quad B=\ell,\quad D=a/\ell,
 \qquad C_0=(1+p\ell)/3,\quad C_1=(1+\ell)/3.
\]
Here the original tag \(\varepsilon=0\) uses \(C_0\), and
\(\varepsilon=1\) uses \(C_1\); their respective ordered triples are
\[
 (a,C_0D,p\ell C_0D),\qquad
 (a,pC_1D,p\ell C_1D).
\]
\end{theorem}
\begin{proof}
The identity \(a=3h+1\) gives \(a\equiv1\pmod3\), so no prime
factor is three. The exterior condition \(3\mid4u+1\) and the
middle condition \(3\mid u+a\) both say \(u\equiv2\pmod3\).
Every factor \(\ell\equiv1\pmod3\) contributes one to the
residue product, and every \(\ell\equiv2\pmod3\) contributes
\((-1)^{e_\ell}\). This proves both inclusions in the displayed
fibre equality, with the stated valuation inverse.

There are \(P_2\) exponent choices on the second residue class.
The number of even sums minus the number of odd sums equals
\[
 \prod_{\substack{\ell\mid a\\\ell\equiv2\ (3)}}
       \left(\sum_{e=0}^{2v_\ell}(-1)^e\right)=1,
\]
because every inner alternating sum is one. Thus the odd sums
number \((P_2-1)/2\); the first residue class independently
contributes all \(P_1\) choices. This proves the cardinality.
The complement \(u\mapsto a^2/u\) preserves the middle set:
its residue is \(1/2=2\) modulo three. Its only positive fixed
point would be \(u=a\), which has residue one and is absent.
Hence \(|M_a|\) is even, and the displayed original weighted
shell count is the stated integer.

If there is no second-class factor, every exponent sum is zero
and the fibre is empty. If \(\ell\equiv2\pmod3\) divides \(a\),
the integer \(u=a\ell\) has exponent \(v_\ell+1\le2v_\ell\)
at that prime and exponent \(v_t\) at every other prime \(t\).
It belongs to \(\Div(a^2)\) and has residue two. Its coordinates
satisfy \(a=ABD\), \(u=B^2D\), and \(\gcd(A,B)=1\) exactly.
Both \(C_0,C_1\) are positive integers since \(p\equiv1\pmod3\)
and \(\ell\equiv2\pmod3\). For each tag one has
\(3C_\varepsilon=1+p^{1-\varepsilon}\ell\) and
\(4a=p+3\). Multiplying gives
\[
 4AB C_\varepsilon D
   =A+p^{1-\varepsilon}B+pC_\varepsilon.
\]
Putting the three indicated reciprocals over \(pAB C_\varepsilon D\)
therefore proves their sum is \(4/p\). This proves the complete
arithmetic return with the original ordering and first denominator.
\end{proof}

\begin{theorem}[The complete residual-seven factor sieve]
\label{thm:r7-factor-sieve}
At \(a=a_7\), retain
\[
 n_3=\sum_{\substack{\ell\mid a\\\ell\equiv3\ (7)}}v_\ell,
 \qquad
 n_{24}=\sum_{\substack{\ell\mid a\\\ell\equiv2,4\ (7)}}v_\ell.
\]
Then \(E_a\cup M_a\ne\varnothing\) if and only if at least one
of the following three explicitly computed predicates holds:
\[
 \begin{gathered}
 \text{a prime divisor of }a\text{ is }5\text{ or }6\pmod7;\\
 n_3\ge3;\qquad n_3\ge1\text{ and }n_{24}\ge1.
 \end{gathered}
\]
The following choices give a witness in every passing case:
\[
 \begin{array}{c|c|c}
 \text{retained factors or divisor}&u&\text{channel}\\ \hline
 \ell\mid a,\ \ell\equiv5\ (7)&\ell&E_a\\
 \ell\mid a,\ \ell\equiv6\ (7)&a\ell&M_a\\
 d\mid a,\ d\text{ a product of three }3\ (7)\text{ prime factors}
        &ad&M_a\\
 \ell,t\mid a,\ \ell\equiv3\ (7),\ t\equiv2\ (7)
        &a\ell t&M_a\\
 \ell,t\mid a,\ \ell\equiv3\ (7),\ t\equiv4\ (7)
        &a\ell/t&M_a
 \end{array}
\]
The three factors defining \(d\) are counted with their original
prime multiplicities. No factor is used more often than it divides \(a\).

There is also an exact description of every witness, including those
not selected by this table. Set
\[
 \lambda(1)=0,\quad\lambda(2)=2,\quad\lambda(3)=1,\quad
 \lambda(4)=4,\quad\lambda(5)=5,\quad\lambda(6)=3.
\]
For each original prime \(\ell\mid a\), put
\(\lambda_\ell=\lambda(\operatorname{rem}_7\ell)\), and put
\(b=\sum_{\ell\mid a}\lambda_\ell v_\ell\) in \(\Z/6\Z\).
The full fibres, with the original valuation inverse, are
\[
 E_a=\{u(e):e\in\mathcal B_a,\ \sum_\ell\lambda_\ell e_\ell=5\ (6)\},
\]
\[
 M_a=\{u(e):e\in\mathcal B_a,\ \sum_\ell\lambda_\ell e_\ell=b+3\ (6)\}.
\]
In the finite group ring \(\Z[T]/(T^6-1)\), form
\[
 F_a(T)=\prod_{\ell\mid a}\sum_{e=0}^{2v_\ell}T^{\lambda_\ell e}
        =\sum_{r=0}^5 c_rT^r.
\]
Then \(|E_a|=c_5\), \(|M_a|=c_{\operatorname{rem}_6(b+3)}\).
Every coefficient counts its whole displayed exponent fibre; no
exponent representative replaces the finite integer budget.
\end{theorem}
\begin{proof}
The original unit calculation gives \(\gcd(a,7)=1\).
Explicitly, if seven divided \(a\), then \(p=4a-7\) would be
divisible by seven, contradicting primality and \(p\ge13\).
The powers \(3^0,\ldots,3^5\) modulo seven are
\(1,3,2,6,4,5\); the next power is one, and these six residues
are distinct. They therefore enumerate every nonzero residue
with the listed inverse \(\lambda\).
Since \(4^{-1}=2\pmod7\), the exterior target is
\(-4^{-1}=5=3^5\). The middle equation is \(u/a=-1=3^3\)
in the unit group. Taking these exact finite-group coordinates
proves the two displayed fibres. Expanding the finite product
\(F_a\) chooses one original exponent at each prime; collecting
terms with equal residues modulo six proves each count and its
entire underlying fibre.

Each table entry is a positive divisor of \(a^2\). For \(u=\ell\),
the exponent is one at \(\ell\) and zero elsewhere. For \(u=a\ell\)
the exponent at \(\ell\) increases from \(v_\ell\) to
\(v_\ell+1\le2v_\ell\). If \(n_3\ge3\), select exactly three
factors from the multiset having \(v_\ell\) copies of each
\(3\pmod7\) prime. Their product \(d\) divides \(a\), so
each exponent in \(ad\) is between \(v_\ell\) and \(2v_\ell\).
For the last two rows \(\ell\ne t\), since their residues differ.
For \(a\ell t\), both corresponding exponents increase by one
within their budgets. For \(a\ell/t\), the \(\ell\)-exponent
increases by one and the \(t\)-exponent decreases by one; the
latter is nonnegative since \(t\mid a\). All other exponents
remain exactly those of \(a\).

The first entry has residue five, hence belongs to \(E_a\).
The ratios \(u/a\) of the other entries have residues, in order,
\[
 6,\qquad 3^3=6,\qquad3\cdot2=6,\qquad3/4=3\cdot2=6
 \quad\pmod7.
\]
Thus they belong to \(M_a\). This proves sufficiency with full
integer and exponent checks.

For necessity suppose all three predicates fail. There is no prime
factor with residue five or six. If \(n_3=0\), every prime factor
has residue in \(H=\{1,2,4\}\). Direct multiplication modulo
seven shows \(H\) is the subgroup generated by two, of order
three. Every \(u\mid a^2\) and \(a\) itself then lie in \(H\).
The exterior target five is outside \(H\), and the middle target
\(-a\) is outside \(H\): otherwise its quotient by \(a\)
would put \(-1=6\) in \(H\). Both channels are empty.

The only remaining failed-predicate case has \(n_3=1\) or two
and \(n_{24}=0\). Every prime factor now has residue one or
three. Write \(e=\sum_{\ell\equiv3\ (7)}e_\ell\) for the
sum of the actual divisor exponents. Its range is
\(0\le e\le2n_3\le4\), and \(u\equiv3^e\pmod7\).
No such exponent is five modulo six, so the exterior target
is absent. For the middle condition,
\(u/a\equiv3^{e-n_3}\), with \(-2\le e-n_3\le2\).
None of these integers is three modulo six, so the middle
target is absent. This proves necessity without enlarging
the exponent intervals to their generated group.
\end{proof}

\begin{corollary}[The exact two-shell rejection set and integer reconstruction]
\label{cor:two-shell-sieve-return}
For every original prime \(p=12h+1\), the union of the two complete
shell loci is empty exactly when both of the following hold:
\[
 \begin{gathered}
 \text{every prime divisor of }3h+1\text{ is }1\pmod3;\\
 \text{no prime divisor of }3h+2\text{ is }5\text{ or }6\pmod7,\qquad
 n_3=0\ \text{or}\ (1\le n_3\le2\text{ and }n_{24}=0).
 \end{gathered}
\]
The counts and complete exponent fibres in the preceding theorems
give a sieve for these actual two shells for every such prime,
without a numerical cutoff on \(p\). A surviving prime is one
for which these two shells both fail; no assertion about the other
original shells, or the number of surviving primes, is implied.

For every retained passing pair \((\varepsilon,u)\), with
\(\varepsilon=0\) for the exterior channel and one for the middle,
set
\[
 g=\gcd(a,u),\quad A=a/g,\quad B=u/g,\quad D=g^2/u,\quad
 C=(A+p^{1-\varepsilon}B)/R_a.
\]
These are positive integers and give the exact original ordered witness
\[
 (a,y,z)=(ABD,p^\varepsilon ACD,pBCD),\qquad
 R_a=4a-p,\quad S_a=pa.
\]
The map from the retained tagged exponent vector to
\((p,a,R_a,S_a,\varepsilon,u,A,B,C,D,y,z)\) has singleton fibres
and inverse the original prime valuations of \(u\), retaining
\(p,a,\varepsilon\). All constructed witnesses have
\(a\le(p-1)/2\); membership in a smaller unary layer is imposed
on the computed \(A\), not assumed for every table branch.
\end{corollary}
\begin{proof}
Taking the conjunction of the two proved empty-fibre conditions
gives exactly the displayed rejection set, in both directions.
Each full locus remains available through its exponent fibre,
so the rejection calculation has not discarded any unselected
witness or any original prime factor.

To prove the stated reconstruction directly, fix a prime \(\ell\)
and write \(v=v_\ell(a)\), \(e=v_\ell(u)\), with \(0\le e\le2v\).
Then \(v_\ell(g)=\min(v,e)\), and the exponent of \(g^2/u\)
is \(2\min(v,e)-e\), which is \(e\ge0\) if \(e\le v\),
and \(2v-e\ge0\) if \(e\ge v\). Thus \(D\) is a positive
integer. The exact definitions give \(\gcd(A,B)=1\),
\(a=ABD\), \(u=B^2D\); their inverses are the displayed
formulas for \(g,A,B,D\).

All of \(A,B,D\) are units modulo \(R_a\), because their
prime factors divide \(a\), which is coprime to \(R_a\).
For the exterior channel multiply \(4u+1\equiv0\) by \(A\):
\[
 A(4B^2D+1)=4aB+A\equiv pB+A\pmod{R_a}.
\]
Thus \(R_a\mid A+pB\). For the middle channel,
\(u+a=BD(A+B)\), and the unit \(BD\) gives
\(R_a\mid A+B\). These prove that \(C\) is a positive integer
in each case, with no missing divisibility condition.
Now \((4a-p)C=A+p^{1-\varepsilon}B\), hence
\[
 4ABCD=A+p^{1-\varepsilon}B+pC.
\]
The three reciprocals, over the original common denominator
\(pABCD\), have numerators \(pC,p^{1-\varepsilon}B,A\).
Their sum is therefore \(4/p\), with precisely the asserted
ordering and unchanged \(a,R_a,S_a\). Conversely the graph
retains \(u\), so its original exponent vector and tag are
recovered exactly and uniquely. Finally
\(3h+1\le6h\) and \(3h+2\le6h\) for \(h\ge1\), proving
the retained first-half bound. The integer \(A=a/g\) remains
an explicit coordinate for any further layer condition.
\end{proof}

\paragraph{An occupied residual-seven shell outside the first two layers.}
The prime \(p=373=12\cdot31+1\) has
\(a_7=95=5\cdot19\). Primality is proved by the nonzero remainders
\(1,1,3,2,10,9,16,12\) on division by, respectively,
\(2,3,5,7,11,13,17,19\), all the primes at most
\(\sqrt{373}<20\). Both factors of \(a_7\) are five modulo seven.
Every original divisor is uniquely \(u=5^i19^j\),
\(0\le i,j\le2\), and its residue is \(5^{i+j}\).
For \(i+j=0,1,2,3,4\), these residues are
\(1,5,4,6,2\). The exterior target five occurs exactly at
\((i,j)=(1,0),(0,1)\), so \(E_{95}=\{5,19\}\).
The middle target is \(-95\equiv3\pmod7\), which does not occur;
thus \(M_{95}=\varnothing\).
Their complete primitive coordinates are, respectively,
\[
 (A,B,C,D,\varepsilon)=(19,1,56,5,0),\qquad
 (5,1,54,19,0),
\]
with ordered integer witnesses
\[
 (95,5320,104440),\qquad(95,5130,382698).
\]
Indeed \(7\cdot56=19+373\), \(7\cdot54=5+373\),
and each has \(ABD=95\); the preceding reciprocal calculation
therefore proves both identities with the original prime.
The full fibres show that every hit in this shell has \(A=19\)
or five. In particular no \(A\le2\) code in this shell is a hit.
This determines precisely the failure of the same-shell implication
from occupancy to the first two layers; all other original shells
remain in the full primewise system.

\subsection{The exact all-shell no-hit obstruction}
\label{sec:all-shell-no-hit}

Fix a prime \(p=12h+1\), \(h\geq1\), and retain the original shell interval
\[
 A_p=\{3h+1,\ldots,9h\}.
\]
For each \(a\in A_p\) put \(R_a=4a-p\). The unit lemma already
proved above gives \(\gcd(a,R_a)=\gcd(p,R_a)=1\). Write the actual
prime factorization, with no suppressed factors,
\[
 a=\prod_{\ell\mid a}\ell^{v_\ell(a)},\qquad
 {\cal B}_a=\prod_{\ell\mid a}\{0,1,\ldots,2v_\ell(a)\}.
\]
The map
\[
 \nu_a:{\cal B}_a\longrightarrow\operatorname{Div}(a^2),\qquad
 (e_\ell)_{\ell\mid a}\longmapsto U_a(e):=\prod_{\ell\mid a}\ell^{e_\ell}
 \tag{1}
\]
is bijective. Its inverse is the full valuation vector
\(e_\ell=v_\ell(u)\); in particular, no exponent interval is
replaced by the subgroup it generates.

Let \(U(R_a)=(\mathbb Z/R_a\mathbb Z)^\times\), written
multiplicatively, and let \(\mathbb N[U(R_a)]\) be its integral
semiring of finite formal sums. Since every \(\ell\mid a\) is a
unit modulo \(R_a\), define the finite unit-group packet
\[
 F_{a,R_a}:=\prod_{\ell\mid a}
   \left(\sum_{e=0}^{2v_\ell(a)}[\ell]^e\right)
 =\sum_{g\in U(R_a)}c_{a,R_a}(g)[g].
 \tag{2}
\]
The coordinate reduction map is the typed map
\[
 \pi_{a,R_a}:\operatorname{Div}(a^2)\longrightarrow U(R_a),
 \qquad u\longmapsto [u]_{R_a},
\]
and the exponent-to-unit map is $\pi_{a,R_a}\circ\nu_a$.
Expanding (2) records one and only one original exponent at every
prime. Thus, for every \(g\in U(R_a)\),
\[
 c_{a,R_a}(g)=
 \#\{e\in{\cal B}_a:U_a(e)\equiv g\pmod {R_a}\}. \tag{3}
\]
The two original channel sets are
\[
 E_a=\{u\mid a^2:R_a\mid4u+1\},\qquad
 M_a=\{u\mid a^2:R_a\mid u+a\}. \tag{4}
\]
Their exact coefficient form is
\[
 |E_a|=c_{a,R_a}(-4^{-1}),\qquad
 |M_a|=c_{a,R_a}(-a). \tag{5}
\]
Indeed, \(4\) and \(a\) are units modulo \(R_a\), so (4) is exactly
the pair of residue equations \(u\equiv-4^{-1}\) and
\(u\equiv-a\), and (3) counts the complete valuation fibres.

For later reconstruction, define the direct residual maps
\[
 \rho^E_a,\rho^M_a:{\cal B}_a\longrightarrow\mathbb Z/R_a\mathbb Z,
 \qquad
 \rho^E_a(e)=4U_a(e)+1\pmod {R_a},\qquad
 \rho^M_a(e)=U_a(e)+a\pmod {R_a}. \tag{6}
\]
The fibres of the zero element of these maps are, by (1),
\[
 (\rho^E_a)^{-1}(0)=\nu_a^{-1}(E_a),\qquad
 (\rho^M_a)^{-1}(0)=\nu_a^{-1}(M_a), \tag{7}
\]
and the inverse on either fibre is the coordinatewise valuation
map \(u\mapsto(v_\ell(u))_{\ell\mid a}\). Equations (1)--(7) are
coordinate-level maps between the exponent box, the divisor set,
and the finite unit group; they retain every prime and every
valuation bound.

\begin{theorem}[Shellwise no-hit criterion and integer return]
\label{thm:shellwise-no-hit}
For \(a\in A_p\), the shell has no original witness if and only if
\[
 c_{a,R_a}(-4^{-1})=0
 \quad\hbox{and}\quad
 c_{a,R_a}(-a)=0. \tag{8}
\]
The complete set of ordered witnesses with first denominator \(a\) is
\[
 \mathcal V_{a,p}^{\rm ord}
 =\{(a,y,z):y,z\in\mathbb Z_{>0},\quad 1/a+1/y+1/z=4/p\}.
\]
Its exact tagged coordinate domain is the disjoint union
\[
 \mathcal T_{a,p}
 =\{(0,u,\sigma):u\in E_a,\ \sigma\in\{0,1\}\}
 \ \coprod\ \{(1,u,0):u\in M_a\}.
\]
For a point \((\varepsilon,u,\sigma)\) in this domain, let
\[
 \begin{aligned}
 g&=\gcd(a,u),& A&=a/g,\\
 B&=u/g,&D&=g^2/u,\qquad
 C=\frac{A+p^{1-\varepsilon}B}{R_a}.
 \end{aligned}\tag{9}
\]
Then \(A,B,D,C\) are positive integers, \(a=ABD\), \(u=B^2D\),
and \(\gcd(A,B)=1\). Preserve the original chart values
\[
 y_0=p^\varepsilon ACD,\qquad z_0=pBCD.
\]
The map
\[
 \begin{gathered}
 \Phi_{a,p}:\mathcal T_{a,p}\longrightarrow\mathcal V_{a,p}^{\rm ord},\\
 \Phi_{a,p}(\varepsilon,u,\sigma)=
 \begin{cases}
  (a,y_0,z_0),&\sigma=0,\\
  (a,z_0,y_0),&\sigma=1
 \end{cases}
 \end{gathered} \tag{10}
\]
is a bijection. For its full inverse, retain
\(S_a=pa\), \(d_y=R_a y-S_a\), \(d_z=R_a z-S_a\), and put
\(i=v_p(d_y)\). Then \(d_y d_z=S_a^2\), \(i\in\{0,1,2\}\), and
\[
 \Phi_{a,p}^{-1}(a,y,z)=
 \begin{cases}
  (0,a^2/d_y,0),&i=0,\\
  (1,pa^2/d_y,0),&i=1,\\
  (0,d_y/p^2,1),&i=2.
 \end{cases} \tag{10a}
\]
Equations (9) and the valuation inverse in (1) then recover every
chart coordinate and every original exponent. The ordered cardinality is
\(2|E_a|+|M_a|\).

\begin{proof}
For \(u\mid a^2\), write \(v=v_\ell(a)\), \(e=v_\ell(u)\) at
each original prime \(\ell\mid a\). Then \(0\le e\le2v\), and
the exponents of \(A,B,D\) in (9) are respectively
\[
 v-\min(v,e),\qquad e-\min(v,e),\qquad 2\min(v,e)-e.
\]
All are nonnegative, the first two cannot both be positive, and
their sums for \(ABD\) and \(B^2D\) are \(v\) and \(e\).
This proves integrality, both product identities, and coprimality.
Conversely these chart coordinates recover \(u=B^2D\) and
\(g=BD\), hence every original exponent. The gate in (4) is
equivalent to \(4u+p^\varepsilon\equiv0\pmod{R_a}\). Indeed for
\(\varepsilon=1\) use \(p\equiv4a\) and the unit \(4\).
Keeping all factors, its chart form follows from
\[
 A(4u+p^\varepsilon)
 \equiv pB+p^\varepsilon A
 =p^\varepsilon(A+p^{1-\varepsilon}B)\pmod{R_a}.
\]
The multipliers \(A\) and \(p\) are units modulo \(R_a\), so the
equivalence holds in both directions and \(C\) is a positive
integer. Substitution of the unswapped chart, with common
denominator \(pABCD\), gives the exact integer identity
\[
 \frac{pABCD}{a}+\frac{pABCD}{p^\varepsilon ACD}
 +\frac{pABCD}{pBCD}
 =pC+p^{1-\varepsilon}B+A
 =4ABCD,
\]
because \((4ABD-p)C=R_aC=A+p^{1-\varepsilon}B\).
Thus (10) is a positive integer witness in either retained order.
The original residuals in the unswapped chart are exactly
\[
 \begin{aligned}
 R_a y_0-S_a
  &=p^\varepsilon A D(A+p^{1-\varepsilon}B)-pABD\\
  &=p^\varepsilon A^2D=p^\varepsilon a^2/u,\\
 R_a z_0-S_a
  &=pBD(A+p^{1-\varepsilon}B)-pABD\\
  &=p^{2-\varepsilon}B^2D=p^{2-\varepsilon}u.
 \end{aligned} \tag{10b}
\]
Swapping \(y_0,z_0\) swaps these two residuals.

For completeness, the typed residual domain is
\[
 \mathcal D_{a,p}
 =\{d\in\operatorname{Div}(S_a^2):R_a\mid d+S_a\}.
\]
The map \(\mathcal V_{a,p}^{\rm ord}\to\mathcal D_{a,p}\) sends
\((a,y,z)\) to \(d_y=R_a y-S_a\). The original equation gives
\(R_a yz=S_a(y+z)\), so \(d_y=S_a y/z>0\) and
\(d_y d_z=S_a^2\). Its inverse is
\[
 d\longmapsto
 \left(a,\frac{S_a+d}{R_a},\frac{S_a+S_a^2/d}{R_a}\right).
 \tag{10c}
\]
The first quotient is integral by the gate. Since every divisor
of \(S_a^2\) is a unit modulo \(R_a\), multiplication of the second
numerator by \(d\) gives \(S_a(d+S_a)\equiv0\); therefore the second
quotient is also integral. Both are positive, and expanding the
residual product gives \(R_a yz=S_a(y+z)\), hence the required
reciprocal equation. The formulas recover both residuals and both
denominators, so both inverse compositions are the identity.

As \(p\nmid a\), every \(d\in\mathcal D_{a,p}\) has a unique form
\(d=p^i v\) with \(i\in\{0,1,2\}\), \(v\mid a^2\).
We now compute all three inverse branches without deleting one of them.
For \(i=0\), put \(u=a^2/d\). Multiplication of
\(d+pa\equiv0\) by the unit \(u\), followed by division by the
unit \(a\), gives \(a+pu\equiv0\), which is exactly
\(4u+1\equiv0\) after using \(p\equiv4a\).
Thus \(u\in E_a\), and (10b) recovers \(d_y=a^2/u=d\)
with \((\varepsilon,\sigma)=(0,0)\).
For \(i=2\), put \(u=d/p^2=v\). Division of the gate by the unit
\(p\) gives \(pu+a\equiv0\), again equivalent to the exterior
gate. The swapped chart has \(d_y=p^2u=d\), with
\((\varepsilon,\sigma)=(0,1)\).
For \(i=1\), write \(d=pv\) and put \(u=a^2/v=pa^2/d\).
The gate is \(v+a\equiv0\). Multiplication by the unit \(u\)
and division by the unit \(a\) gives \(a+u\equiv0\).
Thus \(u\in M_a\), and (10b) recovers \(d_y=pa^2/u=d\)
with \((\varepsilon,\sigma)=(1,0)\).
Each modular multiplication and division used a unit, so every
gate equivalence has its converse. These branches prove that
(10a) maps every ordered witness into the asserted tagged domain
and that (10) recovers the original \(d_y\), hence the original
ordered triple by (10c). Conversely (10b) has \(p\)-adic exponent
respectively \(0,2,1\) in the exterior unswapped, exterior swapped,
and middle branches. Formula (10a) therefore recovers the original
\(\varepsilon,u,\sigma\) in each case. Every fibre is a singleton.
The domain has cardinality \(2|E_a|+|M_a|\). It is empty exactly
when both coefficients in (5) vanish, proving (8).
\end{proof}
\end{theorem}

The exact map corresponding to exchange of the last two original
denominators is the involution
\(\iota_{a,p}:\mathcal T_{a,p}\to\mathcal T_{a,p}\) given by
\[
 (0,u,\sigma)\longmapsto(0,u,1-\sigma),\qquad
 (1,u,0)\longmapsto(1,a^2/u,0). \tag{10d}
\]
For the middle branch \(R_a\mid u+a\) implies
\(R_a\mid a^2/u+a\), by multiplication by the unit
\([a]_{R_a}[u]_{R_a}^{-1}\in U(R_a)\).
Its chart changes \((A,B,C,D)\) to \((B,A,C,D)\): indeed
\(a^2/u=A^2D\), its gcd with \(a=ABD\) is \(AD\), and
\(C=(A+B)/R_a\) is unchanged. This proves the exchange in
(10d) directly in the original coordinates. There is no extra
middle swap bit, because it would duplicate these same ordered
witnesses. Define the unordered codomain and the two order-forgetting
maps explicitly by
\[
 \mathcal V_{a,p}^{\rm unord}
 =\{(a,\{y,z\}):(a,y,z)\in\mathcal V_{a,p}^{\rm ord}\},
\]
\[
 \begin{aligned}
 q_{\mathcal V}:\mathcal V_{a,p}^{\rm ord}&\longrightarrow
    \mathcal V_{a,p}^{\rm unord},&
 (a,y,z)&\longmapsto(a,\{y,z\}),\\
 q_{\mathcal T}:\mathcal T_{a,p}&\longrightarrow
    \mathcal V_{a,p}^{\rm unord},&
 q_{\mathcal T}&=q_{\mathcal V}\circ\Phi_{a,p}.
 \end{aligned}
\]
The exact fibres of \(q_{\mathcal T}\) are
\[
 \{(0,u,0),(0,u,1)\},\qquad
 \{(1,u,0),(1,a^2/u,0)\}.
\]
The corresponding \(q_{\mathcal V}\)-fibres are
\(\{(a,y,z),(a,z,y)\}\), by the proved bijection \(\Phi_{a,p}\).
These fibres all have two elements: a fixed middle point would have
\(u=a\) and \(R_a\mid2a\), forcing \(R_a\mid2\), impossible.
Equivalently a witness with \(y=z\) would have \(d_y=S_a\),
forcing \(R_a\mid2S_a\), again impossible by the unit lemma.

Define the shell count and global obstruction without identifying
different shell labels:
\[
 q_{a,p}=|E_a|+\frac{|M_a|}{2},\qquad
 Q_p=\sum_{a\in A_p}q_{a,p},\qquad
 {\mathcal O}_p=\prod_{a\in A_p}
 {\bf1}\!\left[c_{a,R_a}(-4^{-1})=0\right]
 {\bf1}\!\left[c_{a,R_a}(-a)=0\right]. \tag{11}
\]
The middle involution \(u\mapsto a^2/u\) has no fixed point on
\(M_a\): a fixed point would be \(u=a\), forcing
\(R_a\mid2a\), hence \(R_a\mid2\), impossible since \(R_a\ge3\).
Thus \(|M_a|\) is even. The proved two-element fibres show that
\(q_{a,p}\) is exactly the number of unordered pairs underlying
\(\mathcal V_{a,p}^{\rm ord}\), and
\(2q_{a,p}=|\mathcal V_{a,p}^{\rm ord}|\).

\begin{theorem}[Necessary-and-sufficient global obstruction]
\label{thm:global-no-hit}
For every prime \(p=12h+1\),
\[
 {\mathcal O}_p=1
 \Longleftrightarrow
 (\forall a\in A_p)\ E_a=M_a=\varnothing
 \Longleftrightarrow Q_p=0
 \Longleftrightarrow
 \text{\(p\) has no original Erd\H{o}s--Straus witness}. \tag{12}
\]
Consequently a hypothetical counterexample must satisfy the finite,
computable condition
\[
 \forall a\in\{3h+1,\ldots,9h\}:\quad
 [\,-4^{-1}\,]\notin\operatorname{supp}(F_{a,R_a})
 \quad\text{and}\quad
 [\,-a\,]\notin\operatorname{supp}(F_{a,R_a}), \tag{13}
\]
with every support computed from the unabridged boxes \({\cal B}_a\).

\begin{proof}
Each factor in (11) is \(0\) or \(1\), and equals \(1\) exactly
when the corresponding complete fibre is empty by (5)--(7).
Therefore the product is \(1\) exactly when every shell has both
empty channels. The summands \(q_{a,p}\) are nonnegative and
vanish exactly under the same condition, so their finite sum
vanishes exactly then. The full primewise biconditional already
proved above is \(\operatorname{ES}(p)\Longleftrightarrow Q_p>0\);
taking its negation yields the last equivalence. Since (2) is a
finite product, (13) is a finite prime-factor algorithm and not an
occupancy or density assertion.
\end{proof}
\end{theorem}

\begin{corollary}[Residual-three and residual-seven tests]
\label{cor:global-r3-r7}
For \(a=3h+1\), the coefficient in (5) is nonzero exactly when
\(a\) has a prime factor \(\ell\equiv2\pmod3\); this is the
complete residual-three sieve. For \(a=3h+2\), the two
coefficients in (5) are given by the six-class packet and the
three explicit predicates \(n_3,n_{24}\) in
Theorems~\ref{thm:r3-factor-sieve}--\ref{thm:r7-factor-sieve}.
Thus a global no-hit \(p\) must pass both exact local rejection tests,
together with (13) for every remaining shell. These conditions
are necessary and sufficient and retain the complete factor and
shell data.
\end{corollary}

\paragraph{A complete original-order fixture.}
At \(p=13\), \(h=1\), \(a=4\), one has \(R_a=3\), \(S_a=52\),
\(a^2=16\), and \(E_a=M_a=\{2,8\}\). These are exactly the
divisors in \(\{1,2,4,8,16\}\) congruent to \(2\pmod3\).
All six tagged points and their original residuals and denominators are
\[
\begin{array}{c|r|r|r|r}
 (\varepsilon,u,\sigma)&d_y&d_z&y&z\\\hline
 (0,2,0)&8&338&20&130\\
 (0,2,1)&338&8&130&20\\
 (0,8,0)&2&1352&18&468\\
 (0,8,1)&1352&2&468&18\\
 (1,2,0)&104&26&52&26\\
 (1,8,0)&26&104&26&52
\end{array}
\]
Here \(S_a^2=2704=2^4\cdot13^2\). Its passing divisors are
exactly \(2,8,26,104,338,1352\): among
\(2^e13^i\), \(0\le e\le4\), \(0\le i\le2\), the gate
\(d_y+52\equiv0\pmod3\) holds exactly for \(e=1,3\).
Formula (10c) gives precisely the six rows above. The shell has
three unordered pairs and six ordered witnesses; restricting the
chart to \(\sigma=0\) alone would omit the two exterior rows
with \(v_{13}(d_y)=2\).

\paragraph{A checked full-fibre warning.}
At $p=373=12\cdot31+1$ and $a=95=5\cdot19$, one has
$R_a=7$ and
$F_{95,7}=(1+[5]+[4])^2$ in $\mathbb N[U(7)]$ (the two
factors are separately retained, even though they have the same
residue).  The complete exterior fibre is
$E_{95}=\{5,19\}$ and the middle fibre is empty, since the
exponent sums $i+j$ with $0\le i,j\le2$ produce residues
$1,5,4,6,2$ and never $-95\equiv3\pmod7$.  The two retained
exponent vectors reconstruct $(A,B,C,D,\varepsilon)=(19,1,56,5,0)$
and $(5,1,54,19,0)$, so every hit in this shell has $A>2$.
This is a direct finite-fibre check and does not replace the full
all-shell criterion by an $A\le2$ cutoff.

\paragraph{Checked finite algorithm.}
For each input prime \(p=12h+1\), factor every \(a\in A_p\), enumerate
the literal exponent box \({\cal B}_a\), accumulate residues in
\(U(R_a)\), and test the two target coefficients. The checker
\texttt{check\_nohit.py} implements (1)--(13), compares each
coefficient with direct divisor enumeration, verifies (9)--(10d),
compares the complete ordered tagged image with direct enumeration
of every passing divisor of the original \(S_a^2\), checks both
inverse compositions and every order-forgetting fibre, and records
the first shell having a hit. It does not claim that
every shell is occupied, that a prescribed residue is hit by a
prime factor, or that infinitely many survivors exist.


\section{Local congruences with all integral lifts retained}

The following calculation applies to either preceding box, including
the signed box for \(S_a\).  Fix positive integers \(q_1,\ldots,q_s\)
and a finite integer box
\[
 B=\prod_{j=1}^s([L_j,U_j]\cap\Z),\qquad L_j\le U_j,
\]
where negative endpoints are allowed.  Take positive moduli
\(D_1,\ldots,D_t\) prime to \(q_1\cdots q_s\).  If a modulus is one,
its residue group has one element and all its congruences are
vacuous.  For \(D\) among these moduli or their least common
multiples, define
\[
 \phi_D:\Z^s\longrightarrow(\Z/D\Z)^\times,\qquad
 \phi_D(\boldsymbol e)=\prod_j(q_j\bmod D)^{e_j},\qquad
 K_D=\ker\phi_D .
\]
The image \(H_D\) is used only to describe integral fibres.  The
actual arithmetic domain remains \(B\).

\begin{theorem}[Integral exponent gluing and its bounded fibre]\label{thm:glue}
Let \(D,E\) be such moduli and \(L=\lcm(D,E)\).
There is an exact sequence of abelian groups
\begin{equation}\label{eq:lattice-exact}
 0\longrightarrow\Z^s/(K_D\cap K_E)
 \xrightarrow{\ i\ }\Z^s/K_D\oplus\Z^s/K_E
 \xrightarrow{\ \delta\ }\Z^s/(K_D+K_E)\longrightarrow0 ,
\end{equation}
where \(i([\boldsymbol e])=([\boldsymbol e],[\boldsymbol e])\)
and \(\delta([\boldsymbol x],[\boldsymbol y])
=[\boldsymbol x-\boldsymbol y]\).  Also \(K_L=K_D\cap K_E\).

For local targets \(r_D\in H_D,r_E\in H_E\), choose integral
lifts \(\boldsymbol x,\boldsymbol y\), respectively.  Their
well-defined obstruction is
\[
 o(r_D,r_E)=[\boldsymbol x-\boldsymbol y]\in
                        \Z^s/(K_D+K_E).
\]
If it is nonzero there is no integral common exponent lift.
If it is zero, express \(\boldsymbol x-\boldsymbol y
=\boldsymbol k_D+\boldsymbol k_E\) with
\(\boldsymbol k_D\in K_D,\boldsymbol k_E\in K_E\), and put
\[
 \boldsymbol e_0=\boldsymbol x-\boldsymbol k_D
                  =\boldsymbol y+\boldsymbol k_E.
\]
Then the complete set of common integral lifts, and the complete
set of arithmetic lifts, are respectively
\begin{equation}\label{eq:bounded-fibre}
 \boldsymbol e_0+K_L,\qquad
 B\cap(\boldsymbol e_0+K_L).
\end{equation}
In particular the vanishing obstruction is not a claim that the
second set is nonempty.
\end{theorem}
\begin{proof}
Changing a representative in the first map by \(K_D\cap K_E\)
does not change either coordinate; its kernel is zero because
vanishing of both coordinates puts it in this intersection.
Changing representatives in \(\delta\) changes their difference
by \(K_D+K_E\), so \(\delta\) is well defined.  It is onto, since
\(([\boldsymbol z],[0])\) maps to any class \([\boldsymbol z]\).
If \(\delta([\boldsymbol x],[\boldsymbol y])=0\), take the displayed
decomposition and \(\boldsymbol e_0\).  Its two residue classes
are exactly those of \(\boldsymbol x,\boldsymbol y\); hence it
is a preimage under \(i\).  Conversely every diagonal class has
difference zero.  This proves exactness, the asserted obstruction,
and the explicit lift.  Two common lifts differ by both kernels,
and addition of any member of their intersection preserves both
targets, proving the first formula in \eqref{eq:bounded-fibre}.
Requiring precisely the original inequalities gives the second.
Finally a unit is one modulo \(L\) exactly when it is one modulo
both \(D,E\): integer divisibility by both is divisibility by
their least common multiple.  Apply this to the numerator minus
denominator of the positive rational product if any exponent is
negative; that denominator is prime to all the moduli.  This
proves \(K_L=K_D\cap K_E\).
\end{proof}

\begin{proposition}[Noncoprime modulus compatibility and its inverse]\label{prop:crt}
Put \(g=\gcd(D,E)\), \(D=gD_0,E=gE_0\).
The reduction map gives the ring isomorphism
\[
 \Z/L\Z\xrightarrow{\sim}
 \Z/D\Z\times_{\Z/g\Z}\Z/E\Z.
\]
For representatives \(r,s\) of a compatible pair, set
\(c=(s-r)/g\), choose \(v\) with \(D_0v\equiv1\pmod{E_0}\),
and set
\[
 w=r+D(vc\bmod E_0).
\]
Then \(w\bmod L\) is its unique inverse.  The same map restricts
to a bijection on unit groups.  For a pair of unit targets
arising from local exponent maps, this ring-level compatibility
is necessary; after it, the additional integral obstruction
and bounded fibre are those of Theorem~\ref{thm:glue}.
\end{proposition}
\begin{proof}
Compatibility means \(g\mid s-r\), so \(c\) is integral.
The coprimality \(\gcd(D_0,E_0)=1\) gives the inverse \(v\) by
Bezout.  The formula is \(r\) modulo \(D\); modulo \(E\),
\(Dvc=gD_0vc\equiv gc=s-r\), as required.
Two inverse representatives differ by both moduli, hence by \(L\).
Addition and multiplication commute with reduction, proving
the ring statement.  An integer is prime to \(L\) exactly when
it is prime to both \(D,E\), proving the unit statement.
An exponent lift supplies this common residue automatically,
but a common residue need not lie in \(H_L\), and membership
in \(H_L\) need not supply a representative in \(B\).
Theorem~\ref{thm:glue} computes those two further questions.
\end{proof}

These constructions are effective without an unspecified lifting
assumption.  To see this directly, put
\[
 T_j=\operatorname{ord}_{L}(q_j),\qquad
 F=\prod_j\{0,\ldots,T_j-1\},
\]
where the order is one for modulus one.  These finite orders exist:
multiplication by a unit permutes the finite residue group, and
two equal powers can be cancelled.  Divide each exponent by \(T_j\)
with its unique remainder in the indicated interval.  The map
\(\phi_L\) is unchanged.  Consequently its complete kernel is
\[
 K_L=\bigcup_{\substack{\boldsymbol f\in F\\\phi_L(\boldsymbol f)=1}}
       \left(\boldsymbol f+\prod_j T_j\Z\right).
\]
For any target replace \(1\) by that target in this formula.
Intersection with \(B\) is calculated coordinate by coordinate
using the original intervals, with each surviving vector retained.
This enumerates exact integral fibres and their bounded portions.
It may be expensive; it is finite and does not posit their occupancy.
Repeated gluing of \(t\) moduli uses the intersection of the first
\(i\) kernels as the next \(K_D\); the new obstruction is computed
against this whole intersection.  Thus it does not substitute
pairwise compatibility for a simultaneous lift.

There is also a precise additive cohomological presentation.
The two-term complex
\[
 C^0=\Z^s/K_D\oplus\Z^s/K_E
       \xrightarrow{\delta}
 C^1=\Z^s/(K_D+K_E)
\]
has \(H^0=i(\Z^s/K_L)\) and \(H^1=0\), by
\eqref{eq:lattice-exact}.  A specified pair of local targets has
a gluing lift precisely when its difference \(\delta\) is zero.
The finite-box condition is the pullback of \(B\) along that
integral fibre.  In particular the vanishing of this complex's
\(H^1\) does not make every specified pair a cocycle and does not
make the bounded pullback nonempty. Their relation to cyclic
cohomology is the integral coordinate construction proved immediately below.

\paragraph{Cyclic-source interface.}
The degree-zero quotient and fibre-trace viewpoint is compared with
Connes--Consani \cite[Proposition 3.1 and equation (3.2)]{CCTrace}.
The following integral augmentation calculation is proved directly, including
its torsion and finite arithmetic return. It is not imported as a finite-box
lifting theorem from that source.
% Integral bridge for the actual exponent-gluing groups; wrapper-free.
The following construction gives the exact relation to degree-zero cyclic
homology and cyclic cohomology, including the integral torsion and the
original bounded exponent fibres. All group algebras here are over \(\Z\).

\begin{lemma}[Integral augmentation coordinates and their complete fibres]
\label{lem:integral-augmentation}
Let \(G\) be an abelian group written additively. Put
\(A_G=\Z[G]=\bigoplus_{g\in G}\Z u_g\), with
\(u_gu_h=u_{g+h}\) and unit \(u_0\), and put
\(I_G=\ker(\epsilon_G:A_G\to\Z)\), where
\(\epsilon_G(\sum_g n_gu_g)=\sum_g n_g\).
Every sum here has finite support. The maps
\[
 \kappa_G:G\longrightarrow I_G/I_G^2,
 \qquad g\longmapsto (u_g-u_0)+I_G^2,
\]
\[
 \theta_G:I_G/I_G^2\longrightarrow G,
 \qquad \sum_g n_gu_g+I_G^2\longmapsto\sum_g n_gg
\]
are mutually inverse homomorphisms. In degree zero the cyclic-chain
boundary is \(b(a\otimes b)=ab-ba\), so
\(HC_0^{\Z}(A_G)=A_G/[A_G,A_G]=A_G\). The homomorphism
\[
 \Theta_G:HC_0^{\Z}(A_G)\longrightarrow G,
 \qquad \sum_g n_gu_g\longmapsto\sum_g n_gg
\]
has kernel \(\Z u_0+I_G^2\) and fibre over \(g\) equal to
\((u_g-u_0)+\Z u_0+I_G^2\). Its restriction to \(I_G\) has
fibre \((u_g-u_0)+I_G^2\).
For every homomorphism \(f:G\to H\), the unital algebra map
\(f_*:A_G\to A_H\), \(u_g\mapsto u_{f(g)}\), preserves the
augmentation ideals and their squares and satisfies
\(\overline{f_*}\kappa_G=\kappa_H f\).
\end{lemma}
\begin{proof}
The displayed multiplication makes \(A_G\) a commutative unital ring.
An element of augmentation zero has the exact expression
\(\sum_g n_g(u_g-u_0)\). Consequently \(I_G^2\) is generated
as an additive group by the products
\[
 (u_g-u_0)(u_h-u_0)=u_{g+h}-u_g-u_h+u_0.
\]
The coefficient map \(\sum n_gu_g\mapsto\sum n_gg\) sends each
such product to \((g+h)-g-h+0=0\), and thus defines \(\theta_G\).
The same relation proves
\(\kappa_G(g+h)=\kappa_G(g)+\kappa_G(h)\).
Direct substitution gives \(\theta_G\kappa_G(g)=g\).
For \(w=\sum n_gu_g\in I_G\), additivity of \(\kappa_G\) gives
\[
 \kappa_G\theta_G(w+I_G^2)
 =\sum_g n_g(u_g-u_0)+I_G^2=w+I_G^2,
\]
where the final equality uses precisely \(\sum n_g=0\).
This proves the inverse and the restricted fibres. For arbitrary \(w\),
write \(w=\epsilon_G(w)u_0+(w-\epsilon_G(w)u_0)\). The second
summand lies in \(I_G\), while \(\Theta_G(u_0)=0\); hence the
kernel and every full fibre have the asserted form. Commutativity makes
the degree-zero boundary zero, proving the stated homology formula
directly. Finally the basis formula for \(f_*\) respects products,
unit, and augmentation. It sends \(I_G^2\) into \(I_H^2\),
and sends \(u_g-u_0\) to \(u_{f(g)}-u_0\). This proves naturality
on each coordinate, hence on every finite sum.
\end{proof}

\begin{theorem}[Cyclic characters with exact inverse and separation]
\label{thm:cyclic-character-bridge}
Let \(G\) be a finite abelian group and \(T=\Q/\Z\). Use the
cyclic cochain complex obtained by applying
\(\operatorname{Hom}_{\Z}(-,T)\) to the integral cyclic chain complex.
Its degree-zero cohomology consists of the additive traces
\(\tau:A_G\to T\), meaning \(\tau(ab)=\tau(ba)\).
Define the following subgroup of this degree-zero cohomology:
\[
 \mathcal T_G=\{\tau\in HC^0(A_G;T):
              \tau(u_0)=0,\quad \tau(I_G^2)=0\}.
\]
There is an isomorphism with specified inverse
\[
 \operatorname{Hom}(G,T)\longrightarrow\mathcal T_G,
 \quad \chi\longmapsto
 \left[\tau_\chi\!\left(\sum_g n_gu_g\right)=\sum_g n_g\chi(g)\right],
 \quad \chi(g)=\tau_\chi(u_g-u_0).
\]
Each fibre of this isomorphism is a singleton. These traces separate
the augmentation coordinates: \(g=0\) if and only if
\(\tau(u_g-u_0)=0\) for every \(\tau\in\mathcal T_G\).
If \(H\leq G\), every character of \(H\) extends to \(G\),
and the full fibre of restriction over a character \(\chi_H\) is
\(\chi_0+\operatorname{Hom}(G/H,T)\), where \(\chi_0\) is any
one extension and the second summand is pulled back along \(G\to G/H\).
\end{theorem}
\begin{proof}
The only cocycle condition in degree zero is vanishing on
\(ab-ba\), and there is no incoming cochain boundary in degree zero.
Since \(A_G\) is commutative every additive functional is a trace.
For a character \(\chi\), the formula for \(\tau_\chi\) is
additive, sends \(u_0\) to zero, and sends every generator
\(u_{g+h}-u_g-u_h+u_0\) of \(I_G^2\) to zero.
Conversely, for \(\tau\in\mathcal T_G\), set
\(\chi(g)=\tau(u_g-u_0)\). Applying \(\tau\) to that generator
proves \(\chi(g+h)=\chi(g)+\chi(h)\). Also \(\tau(u_0)=0\),
so additivity gives \(\tau(\sum n_gu_g)=\sum n_g\chi(g)\).
These substitutions prove both inverses and all singleton fibres.

Here is an explicit extension argument retaining all choices. Given
\(H\leq G\), \(\chi:H\to T\), and \(x\in G\), let \(m\)
be the least positive integer with \(mx\in H\); it exists because
\(G\) is finite. Choose a rational representative \(r\) of
\(\chi(mx)\). The solutions of \(m\beta=\chi(mx)\) in \(T\)
are exactly
\(\beta=(r+j)/m+\Z\), \(j=0,\ldots,m-1\): multiplying
verifies them, and subtracting two solutions shows their difference
is one of the \(m\) classes annihilated by \(m\).
For each such \(\beta\), define
\[
 \chi'(h+kx)=\chi(h)+k\beta\qquad(h\in H,\ k\in\Z).
\]
If \(h+kx=h'+k'x\), division of \(k-k'\) by \(m\) and
minimality of \(m\) show \(k-k'=\ell m\) for some integer
\(\ell\); then \(h-h'=-\ell mx\). The difference of the
two proposed values is \(-\ell\chi(mx)+\ell m\beta=0\).
Thus the formula is well defined and additive. It extends \(\chi\),
and every extension to \(H+\Z x\) is obtained uniquely by its
value \(\beta\) on \(x\). Adjoining an element outside the current
subgroup strictly increases its cardinality, so finitely many such
steps extend \(\chi\) to \(G\). The difference of two extensions
vanishes on \(H\), and factors uniquely through \(G/H\); adding
any such factor produces another extension. This proves the full fibre.

For nonzero \(g\), let its order be \(n>1\). The formula
\(kg\mapsto k/n+\Z\) defines a character on \(\langle g\rangle\):
two indices represent the same group element exactly when their
difference is divisible by \(n\). Extend this character by the
preceding construction. The resulting trace takes
\(u_g-u_0\) to \(1/n+\Z\ne0\). For \(g=0\), every such
evaluation is zero. This proves separation without discarding
the finite-order coordinates.
\end{proof}

\begin{theorem}[The original lattice obstruction in cyclic coordinates]
\label{thm:cyclic-lattice-gluing}
Keep the original \(q_1,\ldots,q_s,D,E,L,K_D,K_E,K_L\) and
the exact exponent box \(B\) of \eqref{eq:lattice-exact}. Set
\[
 G_L=\Z^s/K_L,\qquad
 G_0=(\Z^s/K_D)\oplus(\Z^s/K_E),\qquad
 O=\Z^s/(K_D+K_E).
\]
All three groups are finite. Under the coordinates of
Lemma~\ref{lem:integral-augmentation}, \eqref{eq:lattice-exact}
becomes the exact sequence
\[
 0\longrightarrow I_{G_L}/I_{G_L}^2
 \xrightarrow{\overline{i_*}} I_{G_0}/I_{G_0}^2
 \xrightarrow{\overline{\delta_*}} I_O/I_O^2
 \longrightarrow0,
\]
with the explicit generator maps
\[
 [u_{[e]_{K_L}}-u_0]\longmapsto
 [u_{([e]_{K_D},[e]_{K_E})}-u_0],\qquad
 [u_{([x]_{K_D},[y]_{K_E})}-u_0]\longmapsto
 [u_{[x-y]_{K_D+K_E}}-u_0].
\]
The corresponding cyclic-character maps form the exact sequence
\[
 0\longrightarrow\mathcal T_O
 \xrightarrow{\delta^*}\mathcal T_{G_0}
 \xrightarrow{i^*}\mathcal T_{G_L}\longrightarrow0,
 \qquad \delta^*(\tau)=\tau\circ\delta_*,\quad
 i^*(\tau)=\tau\circ i_*.
\]
For a specified local target pair
\(v=([x]_{K_D},[y]_{K_E})\), write
\(o=[x-y]_{K_D+K_E}\). Its integral gluing obstruction vanishes
if and only if every \(\tau\in\mathcal T_O\) has
\(\tau(u_o-u_0)=0\). If it vanishes, choose
\(x-y=k_D+k_E\) with \(k_D\in K_D,k_E\in K_E\), and set
\(e_0=x-k_D=y+k_E\). The complete original integer fibre and
its exact bounded part are, respectively,
\[
 e_0+K_L,\qquad B\cap(e_0+K_L).
\]
They use precisely the original coordinates and intervals; the bounded
part is enumerated by the fundamental-rectangle formula preceding
this construction. No nonemptiness of that intersection is assumed.
\end{theorem}
\begin{proof}
The map \([e]_{K_D}\mapsto\prod_j q_j^{e_j}\bmod D\) embeds
\(\Z^s/K_D\) in the finite unit group modulo \(D\): its kernel
is zero by the definition of \(K_D\). The same argument works
for \(E\) and \(L\). The group \(O\) is a quotient of
\(\Z^s/K_D\), hence is finite as well.
Lemma~\ref{lem:integral-augmentation} gives a bijective group
coordinate map \(\kappa\) at each term. Its naturality gives
the two displayed generator formulas and the commuting squares
with \(i\) and \(\delta\). Thus the middle kernel is exactly
\(\kappa_{G_0}(\ker\delta)=\kappa_{G_0}(\operatorname{im}i)\);
injectivity at the first term and surjectivity at the last follow
by applying the specified \(\theta\) inverses. Explicitly, the fibre
of \(\overline{\delta_*}\) over \(\kappa_O(o)\) is
\[
 \{\kappa_{G_0}(v_0+i(g)):g\in G_L\}
\]
for any \(v_0\) with \(\delta(v_0)=o\), and its parameter \(g\)
is unique because \(i\) is injective. Such a \(v_0\) exists:
if \(o=[t]_{K_D+K_E}\), take
\(v_0=([t]_{K_D},[0]_{K_E})\). A fibre of
\(\overline{i_*}\) over its image is a singleton, with inverse
\(\kappa_{G_L}i^{-1}\theta_{G_0}\); outside that image it is empty.
These are statements about the augmentation quotients. On them the
composite is zero, since \(u_0-u_0=0\), as the formulas show.

Under Theorem~\ref{thm:cyclic-character-bridge}, the dual maps are
\(\chi\mapsto\chi\circ\delta\) and
\(\psi\mapsto\psi\circ i\). Since \(\delta\) is surjective,
the first is injective. A character \(\psi\) on \(G_0\)
vanishes on \(i(G_L)=\ker\delta\) exactly when
\(\chi(o)=\psi(v)\), for any \(v\) with \(\delta(v)=o\),
defines a well-defined character on \(O\). This proves the middle
kernel and supplies the unique inverse on that image. Every character
on \(G_L\) extends from the subgroup \(i(G_L)\) to \(G_0\)
by the proved extension construction, so \(i^*\) is surjective.
Its full fibre over any specified trace is
\(\tau_0+\delta^*\mathcal T_O\), with \(\tau_0\) any one
extension; this follows from the computed kernel. The fibre of
\(\delta^*\) is a singleton on its image and empty elsewhere.

Character separation applied to \(O\) proves the criterion for
\(o=0\). If it holds, the displayed decomposition of \(x-y\)
gives \(e_0\equiv x\pmod{K_D}\) and
\(e_0\equiv y\pmod{K_E}\). A vector \(e\) has these same
two classes exactly when
\(e-e_0\in K_D\cap K_E=K_L\). This proves the entire fibre
and its intersection with \(B\), independently of the decomposition.
Each retained exponent vector reconstructs its original product
\(\prod_jq_j^{e_j}\); the original shell parameters
\(p,h,a,R_a,S_a\), when present, remain fixed throughout.
Every quotient used here has its specified inverse or complete fibre.
\end{proof}


\begin{proposition}[Nonzero integral exponent obstruction in an actual shell]
\label{prop:nonzero-obstruction}
At \(p=241,a=64,R_a=15\), the exterior and middle targets each
have bounded local lifts modulo three and five, but neither
has an integral exponent lift using the permitted prime \(2\),
even before its finite exponent budget is imposed.
\end{proposition}
\begin{proof}
Trial division by \(2,3,5,7,11,13\), the primes at most
\(\sqrt{241}<16\), proves primality.  Here \(h=20\),
\(a=64\in\{61,\ldots,180\}\), \(R_a=4\cdot64-241=15\),
and \(a^2=2^{12}\).  Both targets equal \(11\bmod15\):
\(-4^{-1}=11\bmod15\), and \(-64=11\bmod15\).
Writing \(u=2^j\), the target modulo three is \(2\), hence
\(j\equiv1\bmod2\); the target modulo five is \(1\), hence
\(j\equiv0\bmod4\).  These order statements follow from
\(2^2=1\bmod3\) and the four distinct powers
\(1,2,4,3\bmod5\).
Each congruence has a lift in \(\{0,\ldots,12\}\),
for example \(j=1\) and \(j=0\), respectively.  Their
kernel lattices are \(K_3=2\Z\) and \(K_5=4\Z\).
The obstruction is
\[
 [1-0]\in\Z/(2\Z+4\Z)=\Z/2\Z,
\]
which is nonzero.  Thus no integral exponent lift exists;
the actual bounded fibres are also empty.  Ordinary
residue CRT still has the integral fibre \(11+15\Z\).
This distinguishes the exponent obstruction from
noncoprime residue compatibility, which is vacuous here.
\end{proof}

\begin{proposition}[Zero integral obstruction with an empty bounded fibre]
\label{prop:local-counterexample}
At the actual full shell \(p=37,a=18,R_a=35,S_a=666\), both
channels have divisor hits at each prime-power factor \(5,7\)
of \(R_a\), but both full channels are empty.  For the exterior
target the integral exponent obstruction in
Theorem~\ref{thm:glue} is zero, while its bounded fibre is empty.
\end{proposition}
\begin{proof}
Here \(a^2=2^2 3^4=324\).  The complete divisor list, formed by
all fifteen choices \(u=2^i3^j\), \(0\le i\le2,0\le j\le4\), is
\[
 1,2,3,4,6,9,12,18,27,36,54,81,108,162,324.
\]
Direct substitution of this entire list gives
\[
 \begin{array}{c|l|l}
 D&\{u:D\mid4u+1\}&\{u:D\mid u+18\}\\\hline
 5&1,6,36,81&2,12,27,162\\
 7&12,54&3,108
 \end{array}.
\]
In either channel the two listed sets are disjoint.  Divisibility
by \(35\) requires both tests on the same divisor, so both full
channels are empty.  Their complete origin classification in
Section~1 proves \(\mathcal W_{18}=\varnothing\).

For the precise integral lift, the element \(2\) has order four
modulo five, \(3=2^3\bmod5\), the element \(3\) has order six
modulo seven, and \(2=3^2\bmod7\); their powers can be verified
by multiplication in the respective groups.  Thus
\[
 K_5=\{(i,j):i+3j\equiv0\pmod4\},\qquad
 K_7=\{(i,j):2i+j\equiv0\pmod6\}.
\]
The targets are \(1\bmod5\) and \(5=3^5\bmod7\).
Their exact bounded exponent fibres are
\[
 \{(0,0),(1,1),(2,2),(0,4)\},\qquad
 \{(2,1),(1,3)\}.
\]
The common integral lift \(\boldsymbol e_0=(5,1)\) satisfies
both congruences:
\(5+3=8\) is divisible by four and \(2\cdot5+1=11\equiv5\bmod6\).
It therefore witnesses the vanishing obstruction.  The whole
integral fibre is
\((5,1)+(K_5\cap K_7)\); its intersection with
\(\{0,1,2\}\times\{0,\ldots,4\}\) is empty by the displayed
complete bounded fibres.  At the integer level,
\[
 U(5,1)=96,\qquad4\cdot96+1=385=11\cdot35.
\]
The two prime factors are permitted, but the exponent \(5\)
of \(2\) exceeds its budget \(2\).  The ordinary CRT inverse
is \(u\equiv26\bmod35\), with full fibre \(26+35\Z\).
It includes \(96\), and no element of the actual divisor list.

The primitive orchard description gives the same failed
gluing.  The local pairs \((2,3)\) at modulus five and
\((1,6)\) at modulus seven are positive, primitive, and divide
\(666=2\cdot3^2\cdot37\).  Coordinatewise CRT has the positive
primitive lift \((22,13)\bmod35\), with full integral fibre
\((22,13)+35\Z^2\).  Its displayed numerator contains \(11\)
and its denominator \(13\), both unavailable in \(S_a\).
No other representative satisfies all the primitive divisor
bounds: if one did, it would belong to the full set
\(\mathcal W_{18}\), already proved empty by the complete
E/M calculation.  This proves the failure at the common
finite box, not a failure of the CRT map itself.
\end{proof}

For clarity, the exact finite gluing also has a direct
set-theoretic fibre-product form.  If \(B\) is the original
box and \(F_i\subset B\) is the passing set for one prime-power
factor, give each \(F_i\) its inclusion in \(B\).  Then
\[
 \bigcap_i F_i\ \xrightarrow{\sim}\
 F_1\times_B\cdots\times_B F_t,\qquad
 e\longmapsto(e,\ldots,e).
\]
The inverse takes any component, since all components have
equal image in \(B\); both compositions are identities.
The product over a single unlabelled point would instead
permit different exponent choices, exactly the error in
the preceding example.

\section{Bounded divisibility and simultaneous congruences}

Here is a version that allows nonunits, noncoprime moduli, several
divisor bounds, and mandatory factors.  It will be applied to the
original E/M expressions, not to an unspecified lifting predicate.
Take positive integers \(N_1,\ldots,N_v\), positive integers
\(f_1,\ldots,f_w\) (the list may be empty), and integers
\(c_i,b_i\), positive moduli \(M_i\) for \(1\le i\le t\)
(this list may also be empty).  The exact problem is to find all
\(u\in\Z_{>0}\) satisfying
\begin{equation}\label{eq:affine-system}
 u\mid N_j\ (1\le j\le v),\qquad
 f_j\mid u\ (1\le j\le w),\qquad
 c_i u\equiv b_i\pmod{M_i}\ (1\le i\le t).
\end{equation}
Assume \(v\ge1\), so this positive arithmetic problem is finite.
Zero coefficients \(c_i\) are allowed.

\begin{theorem}[Constructive bounded divisibility CRT]\label{thm:divisor-crt}
The following algorithm returns exactly every solution of
\eqref{eq:affine-system}, with its original positive integer
lift and all exponent coordinates.

First put \(N=\gcd(N_1,\ldots,N_v)\) and
\(f=\lcm(f_1,\ldots,f_w)\), taking \(f=1\) for an empty list.
If \(f\nmid N\), the solution set is empty.  Otherwise write
\[
 N=\prod_{\ell\mid N}\ell^{e_\ell},\qquad
 f=\prod_{\ell\mid N}\ell^{d_\ell},\qquad
 B_{f,N}=\prod_{\ell\mid N}\{d_\ell,\ldots,e_\ell\}.
 \tag{\(\ast\)}
\]
For each affine congruence let
\(g_i=\gcd(c_i,M_i)>0\).  If \(g_i\nmid b_i\), the solution
set is empty.  Otherwise put \(m_i=M_i/g_i\).
When \(m_i>1\), solve the single unit equation
\[
 r_i\equiv(c_i/g_i)^{-1}(b_i/g_i)\pmod{m_i}.
\]
When \(m_i=1\), put \(r_i=0\) in its one-element ring.
Successively combine the residues \((r_i\bmod m_i)\)
by Proposition~\ref{prop:crt}.  A failed compatibility test
gives an empty solution set.  If all tests pass, they give
one residue \(r\bmod L\), \(L=\lcm_i m_i\); for an empty
list take \(L=1,r=0\).

For \(\ell\mid N\), perform the finite residue-monoid
calculation
\begin{equation}\label{eq:monoid-dp}
 F_0(X)=[1],\qquad
 F_k(X)=F_{k-1}(X)
        \left(\sum_{j=d_{\ell_k}}^{e_{\ell_k}}[\ell_k^j]\right)
        \quad\text{in }\N[\Z/L\Z]_{\rm mult}.
\end{equation}
Here \([x][y]=[xy\bmod L]\), and the subscript specifies
the multiplicative monoid algebra of the full residue ring,
including its zero and nonunits.  Then the coefficient of
\([r]\) in the final product is exactly the number of
solutions.  A contribution to that coefficient with chosen
exponents \((j_\ell)\) gives the unique lift
\[
 u=\prod_{\ell\mid N}\ell^{j_\ell}.
\]
Retaining the contributing exponent choices gives every
solution once and gives their complete inverse coordinates
\(j_\ell=v_\ell(u)\).  These are precisely the positive
bounded lifts in the residue fibre
\begin{equation}\label{eq:numeric-fibre}
 r+L\Z,\qquad
 \left\lceil\frac{1-r}{L}\right\rceil
 \le\frac{u-r}{L}\le
 \left\lfloor\frac{N-r}{L}\right\rfloor,
\end{equation}
intersected with the box \((\ast)\) by its exact evaluation
map.  The reduction homomorphism \(\Z\to\Z/L\Z\) has kernel
\(L\Z\) and the displayed cosets as fibres.
\end{theorem}
\begin{proof}
The conditions \(u\mid N_j\) for every \(j\) are equivalent,
prime by prime, to \(u\mid N\); the conditions \(f_j\mid u\)
are equivalent to \(f\mid u\).  Hence \(f\nmid N\) is
impossible, and otherwise unique factorization gives exactly
the independent inequalities in \((\ast)\).  The empty
prime list \(N=f=1\) has one empty vector evaluating to one.
This includes every positive divisor and excludes zero.

A congruence \(c_i u-b_i=M_i z_i\) implies \(g_i\mid b_i\).
Under that divisibility it is equivalent to
\[
 (c_i/g_i)u\equiv b_i/g_i\pmod{m_i}.
 \tag{\(\dagger\)}
\]
When \(m_i>1\), the coefficient is prime to \(m_i\):
a common prime divisor after division would contradict the
definition of their full gcd \(g_i\).  Bezout therefore
gives the displayed inverse, proving equivalence to
\(u\equiv r_i\bmod m_i\) in both directions.
When \(c_i=0\), one has \(g_i=M_i,m_i=1\);
the initial test is exactly \(M_i\mid b_i\), after which
the congruence is vacuous.  Thus zero and nonunit
coefficients have also been treated exactly.
Proposition~\ref{prop:crt}, applied inductively to the
entire current residue and the next one, produces precisely
all simultaneous congruences, or proves that none exist.
It retains the gcd compatibility at every stage; the final
kernel is \(L\Z\) by divisibility by all the \(m_i\).

It remains to solve the finite budget in that fibre.
By definition of multiplication in the monoid algebra,
expanding \(F_k\) chooses one allowed exponent for each of
the first \(k\) labelled primes, multiplies their exact
residues, and adds one to its coefficient.  Induction on
\(k\), starting from the one empty choice, proves this
statement including multiplicities.  No cancellation or
inverse is taken in this expansion.  Thus the final
coefficient at \(r\) counts exactly those and only those
vectors in \((\ast)\) whose integer evaluation is \(r\bmod L\).
Distinct vectors give distinct positive integers by unique
factorization.  Their product evaluates to a divisor of
\(N\), divisible by \(f\), and satisfies all the reduced
congruences.  Reversing \((\dagger)\) proves all the
original affine congruences.  Conversely every solution
has one of those exact exponent vectors and necessarily
contributes to that coefficient.  This proves the
complete positive lifting assertion rather than assuming
it.  Finally \(1\le u\le N\) gives exactly the bounds
in \eqref{eq:numeric-fibre}; they do not replace the
prime-exponent inequalities, which remain in the
intersection already computed by the product.
\end{proof}

For the actual shell \(a\), the exterior substitution in
this theorem is
\[
 N_1=a^2,\quad f=1,\quad(c_1,b_1,M_1)=(4,-1,R_a),
\]
and the middle substitution is
\[
 N_1=a^2,\quad f=1,\quad(c_1,b_1,M_1)=(1,-a,R_a).
\]
One may replace the single modulus by any list of factors
whose least common multiple is \(R_a\), including
overlapping noncoprime factors; the simultaneous conditions
are then exactly the original divisibility.  The coefficients
of \eqref{eq:monoid-dp} compute the full \(E_a,M_a\).
Every recorded exponent choice reconstructs the original
positive denominators by Section~1.  A zero coefficient
proves emptiness of that fixed finite channel.  It does
not prove that all channels at all admissible shells
vanish or that at least one is positive for every prime.

For primitive pairs the equally exact two-coordinate
monoid calculation retains the exclusive use of each
prime budget.  If \(S=\prod\ell^{b_\ell}\), replace one
factor in \eqref{eq:monoid-dp} by
\[
 [(1,1)]+\sum_{j=1}^{b_\ell}
       \left([(\ell^j,1)]+[(1,\ell^j)]\right)
       \quad\text{in }\N[(\Z/L\Z)^2]_{\rm mult}.
\]
Multiplication is coordinatewise.  Its terms choose
exactly one of: no factor, a positive exponent on \(m\),
or a positive exponent on \(n\).  Unique factorization
therefore identifies the choices with the original
coprime pairs \(m,n\mid S\), bijectively.  Summing the
coefficients over \(r+s\equiv0\bmod R_a\), with
\(L=R_a\), counts exactly \(\mathcal W_a\); storing choices
returns each actual pair and then
\eqref{eq:reconstruct}.  If another divisor or congruence
mark is required, its exact test is imposed on the same
stored pair, or a further coordinate is kept in the
finite product.  Independent choices of representatives
are never substituted for the same integer coordinates.


\paragraph{Arithmetic lineage.}
The original finite-code and simultaneous-constraint input is the colleague's
continuation \cite{UmbrellaContribution}, especially its equations (71)--(100).
The intermediate source is \cite{PrimewiseCoordinates}.
The following construction proves the complete bounded code and witness
correspondences used here; it does not attribute this later proof to that input.
\section{Composition with the exact unary Boolean indicator}
\label{sec:unary-transport}

The parent text \emph{The complete primewise coordinate system},
in its section \emph{Exact unary counts and simultaneous logical
constraints}, defines a rectangular unary code, the indicator
\(\chi_p\), and truth polynomials for Boolean operations.
We use those exact coordinates here.  The additional result is the
complete bijection to the bounded divisor boxes and the computation
of every simultaneous truth fibre on that one retained domain.
The classical torsor identities proved in that text are not needed
again for this composition.

Keep the full range, with these names distinguished from
the shear matrix \(J\) and from a general divisor bound:
\[
 J_{\rm code}=6h,\quad M_{\rm code}=9h,\quad
 \Lambda=2J_{\rm code}M_{\rm code},\quad
 N_{\rm code}=2J_{\rm code}M_{\rm code}^{\,2}.
\]
For this section only abbreviate \(J_{\rm code},M_{\rm code}\)
by \(J_0,M_0\).  The full rectangle consists of
\[
 0\le j<J_0,\quad 1\le A,B\le M_0,\quad
 \varepsilon\in\{0,1\},
\]
with shell \(a=3h+j+1,\ R=4j+3,\ S=pa\).
Its code and decoder are
\begin{align}
 c&=(A-1)\Lambda+2M_0j+2(B-1)+\varepsilon,
       \label{eq:unary-code}\\
 A&=1+\lfloor c/\Lambda\rfloor,\quad
 v=c-\Lambda\lfloor c/\Lambda\rfloor,\quad
 j=\lfloor v/(2M_0)\rfloor,\notag\\
 B&=1+\lfloor v/2\rfloor-M_0j,\qquad
 \varepsilon=v-2\lfloor v/2\rfloor. \notag
\end{align}
Successive division by \(\Lambda,2M_0,2\) proves these are
inverse for \(0\le c<N_{\rm code}\).

Write \([P]\) for the integer \(0\) or \(1\) according as the
specified predicate \(P\) is false or true.  Define
\[
 \kappa_p(c)=[AB\mid a]\,[\gcd(A,B)=1],\qquad
 \chi_p(c)=\kappa_p(c)\,[R\mid A+p^{1-\varepsilon}B].
\]
This is precisely the parent's indicator: \(4AB\mid p+4j+3\)
is equivalent to \(AB\mid a\), while its gcd-floor bits
equal these divisibility and coprimality bits.
Set \(\mathcal C_p^{\rm bud}=\{c:\kappa_p(c)=1\}\) and
\(\mathcal C_p^+=\{c:\chi_p(c)=1\}\).

For completeness the parent's integer truth-polynomial convention is
\[
 \neg b=1-b,\quad b\land d=bd,\quad b\lor d=b+d-bd,
\]
\[
 b\Rightarrow d=1-b+bd,\qquad
 b\Longleftrightarrow d=1-b-d+2bd.
\]
For \(b=0\), the last two expressions are \(1\) and
\(1-d\); for \(b=1\), they are both \(d\).
The first three follow by the same two substitutions.
Induction on the expression therefore gives its exact
zero-or-one truth value.  We use that convention for \(F\) below.

\begin{theorem}[Unary codes and one retained tagged valuation box]
\label{thm:unary-box}
There is an explicit bijection
\[
 \Theta:\mathcal B_p:=
 \coprod_{\substack{0\le j<6h\\\varepsilon\in\{0,1\}}}
 \{(j,\varepsilon)\}\times
 \prod_{\ell\mid a_j}\{0,\ldots,2v_\ell(a_j)\}
 \ \longrightarrow\ \mathcal C_p^{\rm bud}.
\]
For \((j,\varepsilon,\boldsymbol e)\), its forward map is
\[
 u=\prod_{\ell\mid a_j}\ell^{e_\ell},\quad
 g=\gcd(a_j,u),\quad A=a_j/g,\quad B=u/g,\quad D=g^2/u,
\]
followed by \eqref{eq:unary-code}.
Its inverse decodes \(j,A,B,\varepsilon\), sets
\[
 D=a_j/(AB),\qquad u=B^2D,\qquad
 \boldsymbol e=(v_\ell(u))_{\ell\mid a_j}.
\]
All fibres of these inverse maps are singletons.
The preimage \(\mathcal B_p^+=\Theta^{-1}(\mathcal C_p^+)\)
is exactly
\[
 R_j\mid 4u+p^\varepsilon,
\]
on the same box.  This is the exterior condition for \(\varepsilon=0\)
and the middle condition for \(\varepsilon=1\).
Every point in this preimage reconstructs the positive integers
\begin{equation}\label{eq:unary-integer-lift}
 C=\frac{A+p^{1-\varepsilon}B}{R_j},\qquad
 (a_j,y,z)=(ABD,p^\varepsilon ACD,pBCD).
\end{equation}
On its image, the inverse from the ordered triple retains its first
denominator, obtains the reduced pair
\[
 (m,n)=\operatorname{red}\big((R_jy-S_j)/S_j\big),
\]
and then recovers \(A=m,\ B=n/p^{1-\varepsilon}\), the tag
\(\varepsilon=0\) when \(p\mid n\) and \(\varepsilon=1\) otherwise,
and the displayed inverse code and exponents.
\end{theorem}
\begin{proof}
Unique prime factorization makes the exponent box bijective to
\(\Div(a_j^2)\), with the displayed valuation inverse.
Write \(a_j=gA,u=gB\) with \(\gcd(A,B)=1\).
Since \(u\mid a_j^2\), one has \(B\mid gA^2\), hence
\(B\mid g\).  Therefore \(D=g/B=g^2/u\) is a positive
integer and \(a_j=ABD,u=B^2D\).
In particular \(AB\mid a_j\) and \(A,B\le a_j\le9h\);
so the forward code belongs to the specified rectangle
and has \(\kappa_p=1\).
Conversely \(D=a_j/(AB)>0\) gives
\(u=B^2D\mid a_j^2\), with quotient \(A^2D\);
coprimality gives \(\gcd(a_j,u)=BD\).
It follows that the two maps recover \(A,B,D,u\) exactly.
Unique prime factorization and the code decoder recover every
remaining label, proving both inverse compositions.

The integers \(A,p,R_j\) satisfy
\(\gcd(A,R_j)=\gcd(p,R_j)=1\).
Multiplying \(4u+p^\varepsilon\) by \(A\) modulo \(R_j\),
and using \(4ABD=p+R_j\), gives
\[
 A(4B^2D+p^\varepsilon)
       \equiv pB+p^\varepsilon A
       =p^\varepsilon(A+p^{1-\varepsilon}B)\pmod{R_j}.
\]
Both multipliers are units, so this is an equivalence.
For \(\varepsilon=1\), \(4u+p\equiv4(u+a_j)\pmod{R_j}\)
and \(4\) is a unit, giving exactly the middle gate.

On the passing locus \(C,D>0\) are integers and
\[
 4ABCD=pC+A+p^{1-\varepsilon}B.
\]
With common denominator \(pABCD\), the three reciprocals
in \eqref{eq:unary-integer-lift} have these three
numerators.  Their sum is therefore \(4/p\).
Moreover direct substitution yields
\[
 R_jy-S_j=p^\varepsilon A^2D,\qquad
 \frac{R_jy-S_j}{S_j}=\frac{A}{p^{1-\varepsilon}B}.
\]
This fraction is reduced because \(\gcd(A,B)=1\) and
\(A\le a_j<p\).  Also \(B<p\), so the divisibility
of its denominator by \(p\) recovers the tag.
This proves the asserted inverse on the exact image,
including the retained ordering.  The exterior chart chooses
the orientation \((A,pB)\); its transpose is recovered by
exchanging \(y,z\), as in Section~1.
In the middle channel the source chart has pair \((A,B)\),
whereas the map \(\mu_a(u)\) below has pair \((B,A)\);
their exact correspondence is the coordinate exchange \(P\),
with the same exchange of the final two denominators.
No orientation is silently identified.
\end{proof}

\begin{theorem}[Simultaneous truth fibres, unions, and intersections]
\label{thm:unary-truth-fibres}
Let \(P_1,\ldots,P_t\) be any finite list of specified arithmetic
predicates decidable by exact integer operations on
\(\mathcal C_p^+\), including divisibility, congruences,
equalities and inequalities on the reconstructed coordinates.
Define the evaluation map on that finite domain by
\[
 \rho(c)=([P_1(c)],\ldots,[P_t(c)])\in\{0,1\}^t.
\]
For \(\sigma\in\{0,1\}^t\), the full fibre is
\[
 H_\sigma=\{c\in\mathcal C_p^+:\rho(c)=\sigma\},\qquad
 \widehat H_\sigma=
 \{\xi\in\mathcal B_p^+:\rho(\Theta(\xi))=\sigma\}.
\]
The restriction of \(\Theta\) is a bijection
\(\widehat H_\sigma\to H_\sigma\), with exactly the
integer lift and inverse of Theorem~\ref{thm:unary-box}.

For any finite Boolean expression \(\Phi\) in these predicates,
including any simultaneous list of biconditionals, put
\(T_\Phi=\{\sigma:\Phi(\sigma)=1\}\).
Its full solution set is the disjoint union
\[
 H_\Phi=\coprod_{\sigma\in T_\Phi}H_\sigma,\qquad
 \Theta^{-1}(H_\Phi)=
           \coprod_{\sigma\in T_\Phi}\widehat H_\sigma .
\]
In particular, for any \(r\ge1\) solution subsets \(X_1,\ldots,X_r\)
of this same coded domain,
\[
 \Theta^{-1}\Big(\bigcap_iX_i\Big)=\bigcap_i\Theta^{-1}(X_i),
 \qquad
 \Theta^{-1}\Big(\bigcup_iX_i\Big)=\bigcup_i\Theta^{-1}(X_i).
\]
The iterated fibre product
\(X_1\times_{\mathcal C_p^+}\cdots
\times_{\mathcal C_p^+}X_r\) over the inclusions is exactly
\(\{(c,\ldots,c):c\in\bigcap_iX_i\}\).
Its inverse diagonal projection has singleton fibres.
The union projection
\(\coprod_iX_i\to\bigcup_iX_i\), \((i,c)\mapsto c\),
has exact fibre \(\{(i,c):c\in X_i\}\).
A disjoint inverse representation of the union is obtained
by choosing the least such \(i\).

The computable exact cardinalities are
\begin{align}
 |H_\sigma|&=\sum_{c\in\mathcal C_p^+}
     \prod_{\sigma_i=1}[P_i(c)]
     \prod_{\sigma_i=0}(1-[P_i(c)]),\label{eq:truth-cardinality}\\
 |H_\Phi|&=\sum_{\sigma\in T_\Phi}|H_\sigma|
       =\sum_{\xi\in\mathcal B_p^+}
                       [\Phi(\rho(\Theta(\xi)))],\notag\\
 \left|\bigcup_{i=1}^rX_i\right|
     &=\sum_{\varnothing\ne I\subseteq\{1,\ldots,r\}}
           (-1)^{|I|+1}\left|\bigcap_{i\in I}X_i\right|.\notag
\end{align}
Equivalently, writing \(F(c)\) for the parent's truth
polynomial, the exact selected tagged count and weighted count are
\[
 \sum_{c=0}^{N_{\rm code}-1}\chi_p(c)F(c),\qquad
 W_F=\sum_{c=0}^{N_{\rm code}-1}
                  (2-\varepsilon(c))\chi_p(c)F(c).
\]
Predicates involving \(C,D,y,z\) are evaluated on the passing
domain where these coordinates exist.  They can be assigned any
fixed bits outside that domain, since the factor \(\chi_p=0\)
removes those codes.  No fractional candidate is treated as an
integer lift.  For \(F=1\), \(W_F=2Q_p\).
For a general selection invariant under middle complementation
\(u\mapsto a_j^2/u\), \(W_F/2\) counts the selected unordered
pairs; without that invariance it is only the displayed
half-weighted tagged count and can be nonintegral.
\end{theorem}
\begin{proof}
Each code has exactly one truth vector, so the fibres \(H_\sigma\)
are pairwise disjoint and partition \(\mathcal C_p^+\).
The inverse bijection \(\Theta^{-1}\) preserves membership
in every predicate evaluated at that same code; this proves
the fibre statement and both set identities.
A tuple in the fibre product has equal images under the
inclusions, which means all of its entries are literally
the same code.  Conversely any member of the intersection
gives that diagonal tuple.  The union fibre follows from
the definition of the coproduct, and the least-index rule
is defined precisely on the union.

The product in \eqref{eq:truth-cardinality} is one exactly
when every bit agrees with \(\sigma\); otherwise one of
its factors is zero.  Summing proves its count.  Disjointness
and the bijection give the other two expressions for
\(|H_\Phi|\).  At a point belonging to \(k>0\) of the
\(X_i\), the inclusion--exclusion sum contributes
\(\sum_{s=1}^k(-1)^{s+1}\binom{k}{s}=1\);
at a point in none it contributes zero.  The equality follows
by expanding \((1-1)^k\) and then summing over the finite domain.
The parent's Boolean truth polynomials evaluate to the same bit
\([\Phi]\) at each code, so their displayed counts agree with
these fibres.

Finally the exterior chart represents both ordered orientations
of one unordered pair and the middle charts represent the two
distinct complementary orientations.  The preceding proofs and
the fixed-point-free middle complementation of Section~1 give
\(W_1=2Q_p\).  An invariant middle selection retains both or
neither of those orientations; division by two then gives the
unordered count.  If one orientation alone is retained its
weight is one, proving both the limitation and the possible
nonintegrality asserted above.
\end{proof}

\begin{proposition}[Constructive CRT evaluation inside every common truth fibre]
\label{prop:unary-crt-fibre}
Every fibre of Theorem~\ref{thm:unary-truth-fibres} can be
computed with its complete integer lifts, using the bounded
divisibility CRT of Theorem~\ref{thm:divisor-crt} for its
positive affine/divisibility atoms and exact finite tests for
all its other atoms.  Noncoprime moduli, nonunit or zero
affine coefficients, negations, and biconditionals are retained.

More explicitly, in one fixed component \((j,\varepsilon)\),
write \(u=\prod_{\ell\mid a_j}\ell^{e_\ell}\).
Include the original bound \(u\mid a_j^2\) and the original
gate \(4u\equiv-p^\varepsilon\pmod{R_j}\).
For a truth vector \(\sigma\), collect every atom required true
of the forms
\[
 u\mid N_i,\qquad f_i\mid u,\qquad
 \alpha_i u\equiv\beta_i\pmod{L_i},
 \quad N_i,f_i,L_i>0,\quad\alpha_i,\beta_i\in\Z,
\]
where the displayed parameters are fixed in that component.
Let \(\mathcal L_\sigma\) be the complete finite list returned
by Theorem~\ref{thm:divisor-crt} on those conditions.
For each \(u\) on this list reconstruct its unique code by
Theorem~\ref{thm:unary-box} and retain it exactly when
\[
 [P_i(c)]=\sigma_i\quad(1\le i\le t).
\]
This gives precisely \(H_\sigma\) in that component, without
multiplicity.  Summing the retained list lengths over
\(j,\varepsilon,\sigma\in T_\Phi\) gives \(|H_\Phi|\).
Storing the codes, exponents, and \eqref{eq:unary-integer-lift}
gives all positive integer lifts and their singleton inverse
fibres, not only that cardinality.
\end{proposition}
\begin{proof}
Each actual member of \(H_\sigma\) satisfies its positive
collected atoms and the original bound and gate.  The full
biconditional in Theorem~\ref{thm:divisor-crt} therefore places
its unique \(u\) in \(\mathcal L_\sigma\).
It passes the final truth-vector test, proving that none is lost.
Conversely each retained list element satisfies all original
arithmetic conditions and every required true or false atom,
so its code is in exactly the stated fibre.
Unique factorization and Theorem~\ref{thm:unary-box} preclude
duplicates in a fixed component; the retained \(j,\varepsilon\)
labels and disjoint truth vectors preclude duplicates between
components or vectors.

For completeness this is a finite algorithm for arbitrary
specified decidable atoms, even when none has the collected
affine form: in that case the list is the bounded divisor
list satisfying the original gate, and every additional atom
is tested by its specified exact rule.  If an atom's coefficients
depend on \(u\) or on another decoded coordinate, that atom is
evaluated in this last step, rather than incorrectly treating
a variable coefficient as constant in CRT.
A false divisibility or congruence atom is likewise tested
on each same integer \(u\); it is not replaced by an independent
choice of a failing representative.  The preliminary gcd
compatibilities, full residue fibre \(r+L\Z\), monoid
coefficients, and retained valuation intervals are exactly
those proved in Theorem~\ref{thm:divisor-crt}.
These calculations produce the entire list or prove it empty;
no existence assumption is added.  This proves termination,
both inclusions, and the reconstruction and counting statements.
\end{proof}

The previously proved actual-shell obstruction \(p=37,a=18\)
makes the common-domain requirement concrete.  On
\(\Div(324)\), the two predicates \(5\mid4u+1\) and
\(7\mid4u+1\) have respectively the nonempty truth sets
\(\{1,6,36,81\}\) and \(\{12,54\}\).
Their intersection on that same divisor domain is empty.
The code bijection transports precisely these two finite
subsets, so it also transports the empty intersection.
The unrestricted integer CRT lift \(u=96\) remains outside
the valuation budget.  Neither taking a Cartesian product
of these local truth sets nor introducing separate
witnesses changes this computed fibre.

\paragraph{Contribution provenance.} The simultaneous constraints continue \texttt{u/UmbrellaCorp\_HR}'s equations (78)--(80), (88)--(90), (96) and (98) \cite{UmbrellaContribution}. This credit identifies the original arithmetic input; the full shared-variable construction below supplies its complete proof and computed lift fibres.

\section{The colleague's shared-variable CRT construction, equations (78)--(98)}
\label{sec:shared-variable-crt}

The additional colleague attachment relayed by the user continues
equations (71)--(100).  It was read in full.  We retain its original
equation numbers in the attributions below.  The results of this
section prove the complete input and output CRT fibres in
(78)--(80) and (96), and construct all the bounded auxiliary
coordinates needed to make (98) an exact theorem with computable
positive integer returns.  No assertion in the attachment is
accepted merely because it has a theorem-status label.

\subsection{Affine input and prescribed output classes}

\begin{proposition}[Complete fibres for source equations (78)--(80)]
\label{prop:input-crt}
Fix a positive integer \(F\) and finitely many
\(\lambda_i,\beta_i\in\Z,\ m_i\in\Z_{>0}\).
The source choice \(F=p!\prod_\nu|E_\nu(p)|\), with
each \(E_\nu(p)\ne0\), is retained as a particular input.
The full input system is
\[
 X\equiv0\pmod{6F},\qquad
 \lambda_iX\equiv\beta_i\pmod{m_i}.
\]
Put \(g_i=\gcd(\lambda_i,m_i)\).  If some \(g_i\nmid\beta_i\),
the system is empty.  Otherwise let
\[
 n_i=m_i/g_i,\qquad
 a_i\equiv(\lambda_i/g_i)^{-1}(\beta_i/g_i)\pmod{n_i},
 \quad 0\le a_i<n_i,
\]
where \(a_i=0\) when \(n_i=1\).
It is solvable if and only if
\[
 \gcd(n_i,n_j)\mid a_i-a_j\quad\text{for every }i,j,
 \qquad \gcd(6F,n_i)\mid a_i\quad\text{for every }i.
\]
When these tests pass, the merge (80), with its stipulated
zero step for reduced modulus one, constructs a unique
residue \(a\in\{0,\ldots,M-1\}\), where
\(M=\lcm(6F,n_1,\ldots,n_s)\).  All integer solutions are
\(a+M\Z\); the associated homogeneous kernel is \(M\Z\).
Its positive solutions and any bounded intersection are exactly
\[
 \{a+Mt:t\in\Z,\ t\ge\lfloor-a/M\rfloor+1\},
\]
\[
 \{a+Mt:\lceil(U-a)/M\rceil\le t\le
                     \lfloor(V-a)/M\rfloor\}\quad(U\le V).
\]
The inverse lift parameter is \(t=(X-a)/M\).
In particular \(X=a+M>0\) satisfies every original condition.
\end{proposition}
\begin{proof}
The coefficient reduction, including \(\lambda_i=0\) and
\(n_i=1\), was proved in Theorem~\ref{thm:divisor-crt}.
The two-residue merge and its full inverse were proved in
Proposition~\ref{prop:crt}.  Necessity of the pairwise tests
follows by subtracting any two required congruences.

Here is a proof that those pairwise tests suffice for the
successive merges, even when the moduli overlap.
Suppose the previously merged modulus is
\(M_0=\lcm(d_0,\ldots,d_k)\), with merged residue \(a_0\);
include \(d_0=6F\) and its zero residue in this list.
For a new residue \(b\pmod n\), consider any prime \(\ell\).
Choose a previous \(d_i\) whose \(\ell\)-valuation is the
largest of the previous valuations.  The common exponent in
\(\gcd(M_0,n)\) is
\(\min(v_\ell(d_i),v_\ell(n))\).
The pairwise condition says that this power divides the
difference between \(b\) and the residue at \(d_i\).
Since \(a_0\) agrees with that residue modulo \(d_i\),
the same power divides \(b-a_0\).
Doing this at every prime proves
\(\gcd(M_0,n)\mid b-a_0\), so the next merge is allowed.
Induction proves all merges.

For the merge, put \(g=\gcd(M_0,n)\) and
\[
 h=\operatorname{rem}_{n/g}
       \left((M_0/g)^{-1}(b-a_0)/g\right).
\]
Then \(a_1=a_0+M_0h\) has both required residues and
the new modulus is \(M_1=M_0n/g\).
If \(0\le a_0<M_0\), the range \(0\le h<n/g\)
gives \(0\le a_1<M_1\).  Thus the claimed representative
range is retained, including \(h=0\) when \(n/g=1\).
The full kernel and fibres follow from the two-residue
proof and induction.  Finally dividing the inequalities
\(0<a+Mt\) or \(U\le a+Mt\le V\) by positive \(M\)
gives precisely the stated integer parameter bounds.
\end{proof}

\begin{proposition}[Prescribed output compatibility and its exact obstruction]
\label{prop:output-crt}
Retain a positive integer seed from the colleague's (88):
\[
 p=12h+1,\quad \varepsilon\in\{0,1\},\quad
 A,B,C,D_p>0,\quad
 4ABCD_p=A+p^{1-\varepsilon}B+pC.
\]
Put
\[
 K_0=4ABC,\quad
 \Theta=C+(1-\varepsilon)B,\quad
 \Omega=A+\varepsilon B,\quad
 g=\gcd(K_0,\Theta),\quad m=K_0/g.
\]
Here \(\Theta\) is the scalar from source equation (88);
it is distinct from the unary bijection of
Theorem~\ref{thm:unary-box}.
For every integer \(q\),
\[
 D_q=(\Omega+q\Theta)/K_0\in\Z
             \quad\Longleftrightarrow\quad q\equiv p\pmod m .
\]
For any \(\mathcal D>0\), the simultaneous output
\(q\equiv1\pmod{\mathcal D}\), \(D_q\in\Z\)
is possible exactly when
\[
 d:=\gcd(\mathcal D,m)\mid p-1.
\]
When this passes, one complete representative and modulus are
\[
 q_0=1+\mathcal D
  \operatorname{rem}_{m/d}
    \left((\mathcal D/d)^{-1}(p-1)/d\right),\qquad
 L_{\rm out}=\lcm(\mathcal D,m).
\]
The remainder is zero when \(m/d=1\).
Every integer output is \(q_0+L_{\rm out}t\), \(t\in\Z\);
its inverse parameter is \((q-q_0)/L_{\rm out}\).
All positive or bounded outputs are obtained by imposing
the corresponding exact linear bounds on \(t\), as above.

The obstruction has the additional exact form
\[
 \operatorname{den}\left(\frac{A+B+C}{4ABC}\right)
       =\frac{m}{\gcd(m,p-1)}>1,
\]
where \(\operatorname{den}\) means the positive denominator
of the reduced rational number.  Thus \(m\nmid p-1\);
\(\mathcal D=12\) is compatible, and any
\(\mathcal D\) divisible by \(m\) is incompatible.
Equivalently, source equation (96) is exactly
\[
 \min(v_\ell(\mathcal D),v_\ell(m))
                      \le v_\ell(p-1)\quad\text{for all primes }\ell.
\]
\end{proposition}
\begin{proof}
The seed identity is \(K_0D_p=\Omega+p\Theta\), so
\(D_q=D_p+(q-p)\Theta/K_0\).
Write \(\Theta=gw\); then \(\gcd(w,m)=1\).
The extra summand is integral exactly when \(m\mid q-p\).
Applying Proposition~\ref{prop:crt} to the two output
classes proves the compatibility test, displayed lift,
homogeneous kernel \(L_{\rm out}\Z\), and its full fibres.
Unique prime factorization proves the valuation form.

At \(q=1\), the same retained identity gives
\[
 \frac{A+B+C}{4ABC}
        =D_p-\frac{(p-1)w}{m}.
\]
Since \(w,m\) are coprime, its reduced denominator is
\(m/\gcd(m,p-1)\).
Each of \(A,B,C\) is at most \(ABC\), and all are positive;
hence the left side lies in \((0,3/4]\).
It is not an integer, which proves the strict inequality
for the denominator and all stated incompatibilities.
For \(\mathcal D=12\), its gcd with \(m\) divides \(12\),
which divides \(p-1\), proving that compatibility.

These calculations prescribe integer output classes.
They neither replace the condition that an output be prime
nor show that a prime divisor of one chosen cyclotomic
value belongs to the prescribed class. The exact incidence
between these two objects, with all input fibres and original
integer witnesses, is computed in the next theorem.
\end{proof}
\input{publication_input_output_incidence.tex}

\subsection{Canonical bounded auxiliary coordinates}

Fix \(p\) throughout the following construction; polynomial
coefficients may be specified integers depending on that fixed \(p\).
Let the original common variables be
\(x=(x_1,\ldots,x_n)\in\Z^n\), with
\[
 l_v\le x_v\le u_v,\qquad
 B_v=\max(|l_v|,|u_v|).
\]
Retain any specified decidable admissibility condition
\(\mathcal A(x)\) on this box, such as the prime label,
the exact shell range, the code decoder, or the valuation budgets.
Its full truth is always tested on one common integer tuple.
Consider a finite Boolean expression \(\Phi\) built from
polynomial equalities, inequalities, and divisibilities
\(d(x)\mid f(x)\), with \(d(x)>0\) on the retained domain.
Comparisons of two polynomials are expressed by their signed
difference, and a congruence is the corresponding divisibility
of that difference.  Implications and biconditionals are
included as Boolean connectives, rather than imposing
either side independently.

For each polynomial \(f=\sum_\alpha c_\alpha x^\alpha\),
define the explicit nonnegative integer bound
\[
 F_f=\sum_\alpha|c_\alpha|\prod_v B_v^{\alpha_v}.
\]
Then \(|f(x)|\le F_f\) throughout the original box.
A factor \(B_v^0\) is one, including when \(B_v=0\).

\begin{lemma}[Unique bounded sign and quotient encoding]
\label{lem:canonical-auxiliary}
The preceding finite truth system admits an explicit polynomial
residual system on a finite augmented box with these properties:
each admissible original tuple has exactly one auxiliary
extension satisfying the residual equations; that extension
records its exact truth bits; and its coordinates have the
following proved bounds.

For a sign test on \(f\), introduce
\(s_-,s_0,s_+\in\{0,1\}\) and
\(0\le t_-,t_+\le\max(F_f-1,0)\).
Use the residual equations
\begin{align}
 s_-+s_0+s_+-1&=0,\notag\\
 f-s_+(t_++1)+s_-(t_-+1)&=0,\label{eq:sign-encoding}\\
 (1-s_+)t_+&=0,\qquad(1-s_-)t_-=0.\notag
\end{align}
The bits for \(f<0,f=0,f>0\) are respectively
\(s_-,s_0,s_+\).

For \(d(x)\mid f(x)\), let \(G=\max(1,F_d)\ge d(x)\)
be the positive integer bound obtained from the absolute
coefficient sum for \(d\) on the original box.  Introduce
\[
 -F_f\le q\le F_f,\qquad 0\le r,s\le G-1,
\]
and residual equations
\begin{equation}\label{eq:quotient-encoding}
 f-dq-r=0,\qquad d-1-r-s=0.
\end{equation}
Then \(q=\lfloor f/d\rfloor\), \(r=f-dq\) is the
canonical remainder and \(s=d-1-r\).  Apply
\eqref{eq:sign-encoding} to \(r\), with bound \(G-1\);
its zero bit is exactly the divisibility bit.
All original Boolean operations are then evaluated once on
these common bits.  Alternatively, introducing one bit per
Boolean subexpression and equating it to its integer truth
polynomial gives a unique bounded polynomial extension as well.
\end{lemma}
\begin{proof}
The first sign equation and the bit domains make exactly one
of \(s_-,s_0,s_+\) equal one.  If it is \(s_+\), the
other equations give \(f=t_++1>0\), \(t_-=0\), and
\(t_+=f-1\).  If it is \(s_-\), they give
\(f=-(t_-+1)<0\), \(t_+=0\), and \(t_-=-f-1\).
If it is \(s_0\), both slacks are zero and \(f=0\).
Conversely these unique choices satisfy all equations and
bounds for each of the three signs; the cases \(F_f=0,1\)
are included by the displayed maximum.  This proves both
existence and uniqueness without an unspecified sign variable.

For division, Euclidean division by \(d>0\) gives unique
integers \(q,r\) with \(f=dq+r,\ 0\le r<d\).
Because \(|f|\le F_f\) and \(d\ge1\), one has
\(-F_f\le\lfloor f/d\rfloor\le F_f\).
Also \(r,s=d-1-r\) lie in \([0,G-1]\).
Conversely the two residual equations and nonnegative
\(r,s\) imply \(0\le r\le d-1\), hence exactly that
canonical Euclidean division, with unique \(s\).
The sign proof applied to \(r\) proves the stated bit.
The finitely many independent auxiliary blocks can therefore
be attached to the same original tuple, each uniquely.
The truth polynomials proved in the parent text take values
zero or one on input bits; evaluating the expression from
its leaves to its root gives a unique bit at every internal
node.  This proves the optional polynomial extension and
its bit bounds by induction on the expression.
\end{proof}

\subsection{The full bounded shared-variable theorem behind equation (98)}

Let \(z=(x,\text{all auxiliaries})\) denote this complete
augmented tuple, whose proved coordinate intervals are
\(\ell_v\le z_v\le u_v\), \(1\le v\le k\).
Keep the original admissibility condition on \(x\) and
all indicated bit domains.  Put
\(B'_v=\max(|\ell_v|,|u_v|)\).
List all residual polynomials from
Lemma~\ref{lem:canonical-auxiliary} and any required
polynomial equalities as
\[
 f_a(z)=\sum_\alpha c_{a,\alpha}(p)z^\alpha.
\]
Required equalities have their own residuals; an equality
appearing inside a negation or biconditional instead has
its computed sign bit and is not incorrectly required true.
Set
\begin{equation}\label{eq:shared-residual-bound}
 H=\max\left(\{0\}\cup
   \left\{\sum_\alpha|c_{a,\alpha}(p)|
                  \prod_v(B'_v)^{\alpha_v}:a\right\}\right).
\end{equation}
Let \(D_{\rm fixed}\ge1\) be any specified common multiple of
all fixed congruence moduli, as in the colleague's equation (98).
Their least common multiple is an explicit possible choice,
with value one for an empty list.  Retain this chosen value and set
\[
 L=D_{\rm fixed}(H+1)>H.
\]
Variable divisor values are handled by
\eqref{eq:quotient-encoding}; they are not treated as
fixed factors of \(D_{\rm fixed}\).

\begin{theorem}[Constructive bounded common lifts and global truth aggregation]
\label{thm:shared-variable-crt}
Choose any nonempty finite list of positive moduli
\(d_1,\ldots,d_s\) whose least common multiple is \(L\).
Overlapping noncoprime moduli are allowed.
Define the vector residue map
\[
 \pi:\Z^k\longrightarrow\prod_{i=1}^s(\Z/d_i\Z)^k.
\]
Its image is exactly the tuples
\(\boldsymbol r=(r_1,\ldots,r_s)\) whose corresponding
coordinates agree modulo every \(\gcd(d_i,d_j)\).
Its additive kernel is \(L\Z^k\).
For each compatible tuple, scalar CRT in each coordinate
constructs its representative
\(r\in\{0,\ldots,L-1\}^k\), and its full integer fibre is
\[
 \pi^{-1}(\boldsymbol r)=r+L\Z^k.
\]
The complete bounded fibre is explicitly
\begin{equation}\label{eq:shared-lift-fibre}
 \left\{r+Lt:
 \left\lceil\frac{\ell_v-r_v}{L}\right\rceil
 \le t_v\le
 \left\lfloor\frac{u_v-r_v}{L}\right\rfloor
 \ \text{for every }v\right\}.
\end{equation}
Its inverse parameter is \(t_v=(z_v-r_v)/L\).
Admissibility and bit-domain restrictions are imposed
on these actual common lifts, retaining every one that passes.

Enumerate exactly those compatible residue tuples for which
\[
 f_a(r_i)\equiv0\pmod{d_i}
       \quad\text{for all residuals }a\text{ and all }i.
\]
For each such tuple, enumerate
\eqref{eq:shared-lift-fibre}, retain the original admissibility
and bit domains, and evaluate the complete Boolean expression
\(\Phi\) on that common lift's canonical bits.
Projection to \(x\) is a bijection from the retained tuples
with \(\Phi=1\) to the original bounded truth-system
solutions.  Its inverse is the unique auxiliary construction
of Lemma~\ref{lem:canonical-auxiliary}.
Thus this algorithm determines existence, all solutions and
their exact integer lifts, including positive ES denominators
when composed with Theorem~\ref{thm:unary-box}.

In particular define
\[
 n(\boldsymbol r)=
 \sum_{\substack{t\text{ in }\eqref{eq:shared-lift-fibre}}}
 [\mathcal A(x(r+Lt))]
 [\text{all specified bit domains at }r+Lt]\,
 [\Phi(\text{recorded bits at }r+Lt)].
\]
For residue tuples that pass the local residual tests above,
the exact solution count is \(\sum_{\boldsymbol r}n(\boldsymbol r)\).
There is no duplicate projection in this count.
The truth fibres, intersections, unions, multiplicities of union
projections and inclusion--exclusion counts of
Theorem~\ref{thm:unary-truth-fibres} hold here on this same
solution domain, with the same common lifted tuple.
\end{theorem}
\begin{proof}
The pairwise sufficiency proof in
Proposition~\ref{prop:input-crt}, applied in each coordinate,
proves the residue image statement.  A vector maps to zero
precisely when every coordinate is divisible by every \(d_i\),
equivalently by \(L\); this proves the kernel.
The coordinatewise CRT representative and kernel give all
integer fibres and the displayed inverse.  Dividing the
original interval inequalities by positive \(L\) proves
\eqref{eq:shared-lift-fibre}, including empty intervals.
The case \(L=1\) has the unique zero residue and
kernel \(\Z^k\), so this description also includes it.

Integer polynomials preserve congruences, since sums and
products preserve them.  Therefore any common lift of a
tuple passing all local residual tests satisfies
\(L\mid f_a(z)\) for every \(a\).
Throughout the proved augmented box,
\eqref{eq:shared-residual-bound} gives
\(|f_a(z)|\le H<L\).  The only multiple of positive
\(L\) in that interval is zero; hence every residual
vanishes as an integer.  Conversely vanishing integer
residuals pass every local residual test.

For an admissible retained lift, the canonical auxiliary
lemma now proves that its recorded bits are the exact
integer truth values of the original predicates, including
all negative remainders and signs.  Thus the single
evaluation \(\Phi=1\) is precisely the original Boolean
requirement.  Every original solution has its unique
auxiliary extension; reducing that extension modulo the
\(d_i\) places it in one of the enumerated compatible
tuples, and its unique lift parameter places it in
\eqref{eq:shared-lift-fibre}.  None is lost.
Conversely each retained lift projects to a solution.
Two retained lifts cannot project to the same \(x\),
because the auxiliary extension is unique, and one
augmented vector has exactly one residue tuple and one
parameter vector \(t\).  This proves both inverse
compositions, the absence of duplicate projections,
and the count.

All products and sums are finite: there are at most
\(L^k\) compatible residue tuples, and each coordinate
interval in the bounded fibre is finite.
Thus this is a terminating exact construction, even if
its enumeration is large.  Applying the previously proved
set-fibre identities on these retained tuples proves
the final union and intersection assertions.
\end{proof}

Two distinctions in the literal content of (98) are essential
and are now calculated, rather than left as lifting assumptions.
First, \(L>H\) makes residual zero-detection exact but need not
make coordinate reduction injective.  For a zero residual,
\(D_{\rm fixed}=2,\ H=0,\ L=2\) and \(-2\le x\le2\),
the zero residue has precisely the three lifts \(-2,0,2\).
Formula \eqref{eq:shared-lift-fibre} returns all three.
Second, a negated composite divisibility is tested after
global aggregation.  In the original shell
\(p=13,a=7,R=15,u=1\), one has \(4u+1=5\):
the congruence passes modulo five and fails modulo three.
Thus \(\neg(15\mid4u+1)\) is true, but
\(\neg(5\mid4u+1)\land\neg(3\mid4u+1)\) is false.
A canonical remainder bit, or the exact global truth
test on the same lift, preserves this difference.

To apply the theorem directly to the complete ES system,
take as original finite domain the unary codes
\(0\le c<N_{\rm code}\) with their exact decoded budget
conditions, or the disjoint tagged valuation boxes of
Theorem~\ref{thm:unary-box}.  Retain every component label,
code, exponent and ordering.  Theorem~\ref{thm:unary-truth-fibres}
evaluates arbitrary additional specified predicates there.
For polynomial residual evaluation, finite decoded
coordinates may instead be included as original bounded
variables with the linear code identity
\eqref{eq:unary-code}, their rectangle bounds, and the
proved integer admissibility tests.
The division and sign construction above supplies every
auxiliary bound required in (98).
The unique code and divisor inverse then gives the original
positive integer denominators, with the exact orientation
in \eqref{eq:unary-integer-lift}.
This proves the requested composition while retaining the
full \(Q_p\) range; \(A\le2\) remains a separate optional
predicate on this same domain.
\begin{corollary}[The actual full primewise system as one polynomial truth system]
\label{cor:actual-prime-truth-system}
For the original prime \(p=12h+1\), use the single rectangle
\[
 0\le j<6h,\quad 1\le A,B\le M_0=9h,\quad
 \varepsilon\in\{0,1\}.
\]
Set, without changing any original shell variable,
\[
 a=3h+j+1,\quad R=4j+3,\quad S=pa,\quad
 T=A+\bigl(p+(1-p)\varepsilon\bigr)B.
\]
The following finite Boolean expression consists solely of
divisibility atoms on that same tuple:
\begin{equation}\label{eq:actual-prime-boolean}
 \Phi_{\rm ES}=
 [AB\mid a]\ \land\
 \bigwedge_{d=2}^{M_0}
       \neg\bigl([d\mid A]\land[d\mid B]\bigr)
       \ \land\ [R\mid T].
\end{equation}
Its successful tuples are exactly the full unary hit set
\(\mathcal C_p^+\) under \eqref{eq:unary-code}.
All divisors in these tests are positive throughout the rectangle.
The canonical quotient construction in
Lemma~\ref{lem:canonical-auxiliary}, applied to \(a/(AB)\)
and \(T/R\), gives the exact variables \(D,C\).
On the truth set they satisfy the explicit bounds
\[
 1\le D\le M_0,\qquad
 1\le C\le C_{\max}:=\left\lfloor\frac{(p+1)M_0}{3}\right\rfloor .
\]
Consequently the original positive integer lifts satisfy
\[
 a\le M_0,\quad u=B^2D\le a^2\le M_0^2,\qquad
 y,z\le pM_0^2C_{\max}.
\]
Appending any specified finite collection of equalities,
positive-divisor congruences, inequalities, implications or
biconditionals on these reconstructed coordinates therefore
gives a fully bounded shared-variable system to which
Theorem~\ref{thm:shared-variable-crt} applies with all
auxiliary bounds constructed above.
The full tagged count and \(2Q_p\) are the respective
unweighted and weighted sums already proved in
Theorem~\ref{thm:unary-truth-fibres}.
\end{corollary}
\begin{proof}
For \(\varepsilon=0,1\), the polynomial
\(p+(1-p)\varepsilon\) equals \(p^{1-\varepsilon}\)
exactly.  Thus \(T\) is the original positive gate numerator.
Also \(a>0,AB>0,R\ge3\), and each fixed \(d\ge2\) is
positive.  If \(\gcd(A,B)>1\), that gcd is an integer in
\([2,M_0]\) dividing both coordinates, so the corresponding
conjunct in \eqref{eq:actual-prime-boolean} fails.
Conversely a failing conjunct supplies a common divisor
greater than one.  The middle conjunction is therefore
exactly \(\gcd(A,B)=1\), proving that the whole expression
is precisely \(\chi_p=1\).

At a hit the two canonical remainders are zero, so their
canonical quotients are \(D=a/(AB)\) and \(C=T/R\).
They are positive because their numerators are positive.
The rectangle gives \(D\le a\le M_0\) and
\(T\le A+pB\le(p+1)M_0\); division by \(R\ge3\)
proves the \(C\) bound.  Since \(a=ABD\),
\[
 u=B^2D\mid A^2B^2D^2=a^2
\]
with exact positive quotient \(A^2D\), proving the
retained divisor budget and its numerical bound.
The reconstruction \eqref{eq:unary-integer-lift},
together with \(p^\varepsilon\le p\) and the established
bounds on \(A,B,C,D\), proves the bounds for both
\(y=p^\varepsilon ACD\) and \(z=pBCD\).
These formulas have already been proved to give the
original ES identity with its exact ordered inverse.

For an appended predicate use these same canonical \(C,D\)
and reconstructed integer coordinates.  Every polynomial
in them has a computable coefficient bound from the
displayed finite intervals; the sign and division lemma
then supplies bounded unique auxiliaries for that predicate.
Equivalently, substitute the coordinate polynomials
\(y=(1+(p-1)\varepsilon)ACD\) and \(z=pBCD\),
retaining their original values and their proved bounds.
The code identity gives a unique code and inverse, so no
additional tuple is introduced by using either presentation.
Applying the shared-variable theorem proves all lifts and
truth fibres; the previously proved multiplicities give
the final count statements.
\end{proof}

\input{publication_raw_scale_hit_incidence.tex}

\paragraph{Boundary-source lineage.}
The same-tag exchange antecedent is \cite{SignedMarginAntecedent}, within
the arithmetic programme originating in \cite{UmbrellaContribution}.
The four-candidate construction below writes out all maps and proofs needed
for its expanded conclusions.
% Wrapper-free proof; uses only amsmath, amssymb and theorem environments.
\section{Four retained coordinate candidates and the exact residual-three first hit}
\label{sec:boundary-four-candidates}

The antecedent \emph{Signed margins, exact witness swaps, and the
prime-output sieve}, Theorem \texttt{thm:margin-swap-bands}, proves
the same-tag exchange, its shell criterion, its oriented divisor map,
and the initial forbidden bands. The source file is
\texttt{signed\_margin\_sieve.tex}, SHA256
\texttt{d51347ad95517c59d4922b526bc55ac1e7a9487f18503f415943b9870ddda6ef}.
We compute here the two additional tag candidates, every failed
candidate's exact rational obstruction, the complete integral orbit
fibres, and the actual least code on every occupied residual-three
shell. All raw coordinates and their source marks remain present.

Fix a prime \(p=12h+1\), \(h\ge1\), and a full original shell
\[
 0\le j<6h,\qquad a=3h+j+1,\qquad R=4j+3=4a-p,
 \qquad S=pa.
\]
Let \(\mathcal R_{p,j}\) consist of the complete records
\[
\begin{gathered}
 X=(p,j,a,R,S,A,B,C,D,\varepsilon),\\
 A,B,D\in\mathbb Z_{>0},\quad ABD=a,\\
 \varepsilon\in\{0,1\},\quad
 C=\frac{A+p^{1-\varepsilon}B}{R}\in\mathbb Q_{>0}.
\end{gathered}
\]
Let \(\mathcal I_{p,j}\) be its exact subset \(C\in\mathbb Z\).
No coprimality restriction is imposed on \(A,B\) in these raw
domains. Since \(a<p\), \(\gcd(a,R)=\gcd(a,p)=1\);
therefore every one of \(A,B,D,p\) is a unit modulo \(R\).
For the first equality, a prime dividing \(a,R\) would divide
\(p=4a-R\), contradicting \(a<p\). For the second use primality
and \(0<a<p\). Finally \(R\) is odd and
\(\gcd(p,R)=1\), since \(p\mid R\) would imply \(p\mid a\).
These unit statements will justify each cancellation explicitly.

\begin{theorem}[The four rational candidates with their complete inverse]
\label{thm:boundary-four-rational}
For \(s,t\in\{0,1\}\), set
\[
 (\alpha,\beta)=
 \begin{cases}(A,B)&s=0,\\(B,A)&s=1,\end{cases}
 \qquad \eta=\varepsilon+t-2\varepsilon t,
 \qquad C_{s,t}=\frac{\alpha+p^{1-\eta}\beta}{R}.
\]
The map \(T_{s,t}:\mathcal R_{p,j}\to\mathcal R_{p,j}\)
retains \(p,j,a,R,S,D\) and replaces
\((A,B,C,\varepsilon)\) by
\((\alpha,\beta,C_{s,t},\eta)\). Its inverse is precisely
\(T_{s,t}\). The four maps form an action of
\((\mathbb Z/2\mathbb Z)^2\): composition adds both displayed
bits modulo two. Every individual map has singleton fibres.
The full rational orbit of \(X\) has four distinct records if
\(A\ne B\), and two if \(A=B\); its complete stabilizer is
\(\{(0,0)\}\) in the former case and
\(\{(0,0),(1,0)\}\) in the latter case.

For every candidate retain its original ordered rational witness and
both oriented residuals:
\begin{align*}
 (x_{s,t},y_{s,t},z_{s,t})
   &=(a,p^\eta\alpha C_{s,t}D,p\beta C_{s,t}D),\\
 \frac1{x_{s,t}}+\frac1{y_{s,t}}+\frac1{z_{s,t}}&=\frac4p,\\
 d_{y,s,t}=Ry_{s,t}-S&=p^\eta\alpha^2D,\\
 d_{z,s,t}=Rz_{s,t}-S&=p^{2-\eta}\beta^2D,
 \qquad d_{y,s,t}d_{z,s,t}=S^2.
\end{align*}
Thus these residuals are positive integers even when one or both
witness denominators are nonintegral. Writing
\(u=B^2D\), the retained target coordinate is
\(u_{s,t}=u\) if \(s=0\), and
\(u_{s,t}=a^2/u=A^2D\) if \(s=1\), with the tag
\(\eta\) still attached. The map to the graph containing the
source \(X\), both bits, and this full target has inverse projection
to \((X,s,t)\), hence singleton fibres; no witness ordering or
choice among equal rational candidates is forgotten.

Set \(k=\gcd(A,B)\). The raw-coordinate correspondence is
\[
 (A_0,B_0,C_0,D_0,k)
   =(A/k,B/k,C/k,k^2D,k),\qquad
 (A,B,C,D)=(kA_0,kB_0,kC_0,D_0/k^2).
\]
Here \(\gcd(A_0,B_0)=1\), \(k^2\mid D_0\), and
\(C_0=(A_0+p^{1-\varepsilon}B_0)/R\).
The maps \(T_{s,t}\) commute with this correspondence, retaining
the very same \(k\) and \(D_0\). Above a fixed coprime record,
the full raw-scale fibre is exactly
\(\{k>0:k^2\mid D_0\}\), including when \(C_0\) is rational.
The integral locus is equivalent to \(C_0\in\mathbb Z\), and
on that locus all displayed witnesses agree under this raw-scale
correspondence coordinate by coordinate.
\end{theorem}
\begin{proof}
The product \(\alpha\beta D=ABD=a\) and the target coefficient
equation put each target in the specified rational domain. Applying
the same bit twice restores \(A,B,\varepsilon\). The coefficient
is the unique rational number satisfying its defining equation, so
it also returns to the original \(C\). The swap and tag operations
commute on these three discrete coordinates, which proves the group
law and all individual inverses. A fixed point must have \(t=0\),
because changing the tag changes that retained coordinate. For
\(s=1,t=0\) it must, and conversely does, have \(A=B\).
These are all four possible stabilizer elements, giving the orbit
counts directly.

The target equation reads
\[
 4\alpha\beta C_{s,t}D
     =\alpha+p^{1-\eta}\beta+pC_{s,t}.
\]
Over the positive rational common denominator
\(p\alpha\beta C_{s,t}D\), the three reciprocals have
numerators \(pC_{s,t},p^{1-\eta}\beta,\alpha\), respectively.
Their sum is consequently \(4/p\). Multiplying
\(RC_{s,t}=\alpha+p^{1-\eta}\beta\) by
\(p^\eta\alpha D\), and then by \(p\beta D\), gives
the two residual identities after subtracting
\(S=p\alpha\beta D\). Their product is
\(p^2\alpha^2\beta^2D^2=S^2\). The formulas for
\(u_{s,t}\) follow from \(a^2=A^2B^2D^2\); applying the
same complement twice restores \(u\), still with its appropriate
tag. Projection from the fully specified graph restores exactly
the source and the two chosen bits, including at fixed points.

The raw-scale formulas are inverse because
\(\gcd(kA_0,kB_0)=k\). Substitution gives
\(A_0B_0D_0=a\) and the stated rational coefficient. The
square-divisor condition is exactly what makes \(D_0/k^2\)
an integer. Thus every scale in the displayed fibre occurs, and
no other scale does. On the original integral locus, reducing
\(RC=A+p^{1-\varepsilon}B\) modulo \(k\) gives
\(k\mid RC\); since \(k\mid a\), \(\gcd(k,R)=1\), so
\(k\mid C\). This proves \(C_0\) integral. Conversely,
\(C=kC_0\) is integral when \(C_0\) is. The same argument
works for every target because its \(k\) is unchanged.
Finally direct substitution gives
\(ABD=A_0B_0D_0\),
\(p^\varepsilon ACD=p^\varepsilon A_0C_0D_0\), and
\(pBCD=pB_0C_0D_0\). The target calculation has the same
identities with \(\alpha/k,\beta/k,C_{s,t}/k\). This proves
the asserted commutation and every retained witness equality.
\end{proof}

\begin{theorem}[Exact obstruction classes, denominators, and four integral fibres]
\label{thm:boundary-obstruction-fibres}
For an integral source \(X\in\mathcal I_{p,j}\), the following
table gives the target numerator modulo \(R\), the common exact
positive reduced denominator \(\delta_{s,t}\) of
\(C_{s,t},y_{s,t},z_{s,t}\), and the complete integrality test.
The original positive rational numbers themselves remain the
fractions defined in the preceding theorem.
\[
\begin{array}{c|c|c|c}
 (s,t)\text{ and source tag}&
 \alpha+p^{1-\eta}\beta\pmod R&\delta_{s,t}&
 T_{s,t}X\in\mathcal I_{p,j}\\\hline
 (0,0),\ \varepsilon=0\text{ or }1&0&1&\text{always}\\
 (1,0),\ \varepsilon=1&0&1&\text{always}\\
 (1,0),\ \varepsilon=0&(1-p^2)B&
 R/\gcd(R,p^2-1)&R\mid p^2-1\\
 (0,1)\text{ or }(1,1),\ \varepsilon=0&(1-p)B&
 R/\gcd(R,p-1)&R\mid p-1\\
 (0,1),\ \varepsilon=1&(p-1)B&
 R/\gcd(R,p-1)&R\mid p-1\\
 (1,1),\ \varepsilon=1&-(p-1)B&
 R/\gcd(R,p-1)&R\mid p-1
\end{array}
\]
Each integral restriction
\[
 T_{s,t}:\{X\in\mathcal I_{p,j}:T_{s,t}X\in\mathcal I_{p,j}\}
   \longrightarrow
 \{X\in\mathcal I_{p,j}:T_{s,t}X\in\mathcal I_{p,j}\}
\]
is a bijection with the same inverse \(T_{s,t}\) and singleton
fibres. If the test fails, the precise rational target and its
inverse in \(\mathcal R_{p,j}\) still apply.

More explicitly, the additive morphism
\[
 \psi_R:\mathbb Z/R\mathbb Z\longrightarrow
      (\mathbb Q/\mathbb Z)[R],\qquad
 \overline n\longmapsto n/R+\mathbb Z
\]
is an isomorphism. The target obstruction is exactly
\(\omega_{s,t}=\psi_R(\overline{\alpha+p^{1-\eta}\beta})
 =C_{s,t}+\mathbb Z\), and has additive order
\(\delta_{s,t}\). For any integer multiplier \(m\) occurring
in the second column, put \(g_m=\gcd(R,m)\) and
\(\delta_m=R/g_m\), with \(\gcd(R,0)=R\).
The map \(\psi_R\circ[m]\) has image exactly
\((\mathbb Q/\mathbb Z)[\delta_m]\), kernel exactly
\(\delta_m\mathbb Z/R\mathbb Z\), and every nonempty fibre
is exactly \(\overline{b_0}+\delta_m\mathbb Z/R\mathbb Z\).
In the table the retained input \(B\) is a unit modulo \(R\),
so its output has the full order \(\delta_m\).
\end{theorem}
\begin{proof}
For \(\varepsilon=0\), the integral source gives
\(A\equiv-pB\pmod R\). Its same-tag swap numerator is
\(B+pA\equiv(1-p^2)B\). Both tag-changing numerators are
\(A+B\equiv(1-p)B\). For \(\varepsilon=1\), the source
gives \(A\equiv-B\pmod R\); its same-tag swap has numerator
\(A+B\equiv0\), while its tag changes have numerators
\(A+pB\equiv(p-1)B\) and
\(B+pA\equiv-(p-1)B\). The identity target has its original
integral numerator. This proves every residue entry independently
of any representative choice.

For an integer \(N>0\), the exact denominator of \(N/R\)
is \(R/\gcd(R,N)\). To prove this statement, write
\(N=gN_1,R=gR_1\), where \(g=\gcd(R,N)\) and
\(\gcd(N_1,R_1)=1\); then \(dN/R\) is an integer exactly
when \(R_1\mid d\), by Bezout. Its least positive such \(d\)
is \(R_1\). Every factor \(B,p,\alpha,\beta,D\) is a
unit modulo \(R\), so
\[
 \gcd(R,mB)=\gcd(R,m),\quad
 \gcd(R,p^\eta\alpha DN)=\gcd(R,N),\quad
 \gcd(R,p\beta DN)=\gcd(R,N).
\]
Apply the denominator calculation to each displayed numerator and
its two indicated multiples. This proves the same exact denominator
for the coefficient and both ordered witness denominators, all table
values, and both directions of every integrality assertion.
The restricted map is its own inverse because the rational map is
its own inverse; if its source and target are integral, the same
condition holds with their roles exchanged.

Changing \(n\) by a multiple of \(R\) changes \(n/R\) by
an integer, so \(\psi_R\) is well-defined and additive. Its
kernel is zero. Every class annihilated by \(R\) is represented
by a rational \(x\) with \(Rx=n\in\mathbb Z\), and is
therefore \(\psi_R(\overline n)\). This proves its stated
codomain and both parts of the isomorphism. The denominator
calculation already proves the order of each \(N/R+\mathbb Z\).
For \(m\ne0\), write \(R=g_m\delta_m\),
\(m=g_mm_1\), \(\gcd(\delta_m,m_1)=1\).
Then \(R\mid mb\) is equivalent to \(\delta_m\mid b\),
giving the complete kernel. Bezout gives an integer inverse of
\(m_1\) modulo \(\delta_m\), so the fractions
\(mb/R=m_1b/\delta_m\) attain every class
\(n/\delta_m+\mathbb Z\); these are precisely the stated
image. Two inputs have the same output exactly when their
difference is in the proved kernel, giving every fibre. For
\(m=0\), the map is zero, \(\delta_m=1\), and these same
descriptions respectively give the zero image, whole domain, and
whole nonempty fibre. Finally \(\gcd(B,R)=1\) gives
\(\gcd(B,\delta_m)=1\), so multiplication by the retained
\(B\) leaves the exact order \(\delta_m\) unchanged.
\end{proof}

\begin{theorem}[Complete integral orbits and the original code-lowering map]
\label{thm:boundary-orbit-code}
Retain either original radix
\[
 (J_H,M_H)=(3h,6h),\qquad (J_F,M_F)=(6h,9h),\qquad
 \Lambda_\rho=2J_\rho M_\rho.
\]
For \(\rho=H\), impose its actual shell domain \(j<3h\);
for \(\rho=F\), keep the whole domain above. Codes use the
retained coprime coordinates of the raw record:
\[
 c_\rho(X)=(A_0-1)\Lambda_\rho+2M_\rho j
                    +2(B_0-1)+\varepsilon.
\]
The rectangle \(1\le A_0,B_0\le M_\rho\),
\(0\le j<J_\rho\), \(\varepsilon\in\{0,1\}\) has
the following full code inverse on
\(0\le c<M_\rho\Lambda_\rho\):
\[
\begin{gathered}
 A_0=1+\lfloor c/\Lambda_\rho\rfloor,\quad
 r=c-(A_0-1)\Lambda_\rho,\quad
 j=\lfloor r/(2M_\rho)\rfloor,\\
 w=r-2M_\rho j,\quad B_0=1+\lfloor w/2\rfloor,\quad
 \varepsilon=w-2\lfloor w/2\rfloor.
\end{gathered}
\]
On a hit this recovers \(a=3h+j+1\), \(R=4j+3\),
\(S=pa\), \(D_0=a/(A_0B_0)\), and
\(C_0=(A_0+p^{1-\varepsilon}B_0)/R\).
Keeping the scale \(k\) then recovers the entire raw record.
Every integral target in the same rational orbit has code
\begin{equation}\label{eq:boundary-four-code-difference}
 c_\rho(T_{s,t}X)-c_\rho(X)
  =s(B_0-A_0)(\Lambda_\rho-2)+(\eta-\varepsilon).
\end{equation}
Let \(\mathcal O_I(X)\) be the rational orbit intersected with
the integral domain. Its complete list and its smallest-code member
are as follows; the coefficient at every listed member is always
its forced coefficient and \(D,k,p,j,a,R,S\) are retained.
\[
\begin{array}{c|c|c}
\text{shell and tag}&\text{members, as }(A,B,\varepsilon)&
                 \text{smallest-code coordinates}\\\hline
 R\mid p-1& (A,B,0),(B,A,0),(A,B,1),(B,A,1)&
               (\min(A,B),\max(A,B),0)\\
 R\nmid p-1,\ R\mid p^2-1&
              (A,B,\varepsilon),(B,A,\varepsilon)&
               (\min(A,B),\max(A,B),\varepsilon)\\
 R\nmid p^2-1,\ \varepsilon=1&(A,B,1),(B,A,1)&
               (\min(A,B),\max(A,B),1)\\
 R\nmid p^2-1,\ \varepsilon=0&(A,B,0)&(A,B,0)
\end{array}
\]
The first row has four distinct members. The second row has two
except when \(A=B\), when it has one. The third row has two
distinct members. The fourth row has one. The map selecting the
listed smallest member, while retaining \(k\), has full fibre
exactly that displayed integral orbit. If \(k\) is also forgotten,
the full fibre above its coprime smallest member is the disjoint
union of these orbits over every \(k>0\) with \(k^2\mid D_0\).
Retaining the source and chosen bits instead gives the singleton
graph inverse of Theorem~\ref{thm:boundary-four-rational}.

Consequently any globally least hit in either original radix lies
among the listed orbit minima. In particular on a shell
\(R\mid p-1\) its tag is zero and \(A_0<B_0\).
On every middle hit the swap lowers the original code exactly when
\(A_0>B_0\). On an exterior hit with \(A_0>B_0\), the
same-tag swap gives a lower integral hit exactly on
\(R\mid p^2-1\); elsewhere
the rational companion and its nontrivial obstruction remain the
exact relation. These statements specify both the retained positive
integer domains and the failed integral fibres.
\end{theorem}
\begin{proof}
The product \(A_0B_0D_0=a\) gives
\(1\le A_0,B_0\le a\le M_\rho\) on the specified shell
range. Swapping retains these same bounds. Thus all candidates
remain in the same code rectangle. The inner digit
\(2(B_0-1)+\varepsilon\) runs bijectively through
\(0,\ldots,2M_\rho-1\). Adding \(2M_\rho j\) runs
bijectively through \(0,\ldots,\Lambda_\rho-1\).
Adding \((A_0-1)\Lambda_\rho\) therefore runs bijectively
through the full displayed interval. Successive Euclidean division
gives exactly the three remainders \(r,w,\varepsilon\) above,
so both encode--decode composites are identities on every original
digit. The hit's defining budget and gate make the displayed
\(D_0,C_0\) positive integers; their formulas and the raw-scale
inverse restore every other recorded coordinate.
The tag increment is precisely
\(\eta-\varepsilon\); the difference of the outer and inner
digits is \((B_0-A_0)\Lambda_\rho+2(A_0-B_0)\) if
\(s=1\), and zero otherwise. This proves the full difference
formula without changing the radix or any shell digit.

The four-candidate table of the preceding theorem gives exactly the
four rows of integral members. If \(R\mid p-1\), equality
\(A=B\) would imply, from either source tag, \(R\mid2A\).
But \(A\) and two are units modulo the odd integer \(R\ge3\),
a contradiction. A middle hit with \(A=B\) gives that same
contradiction even without \(R\mid p-1\). In the second row
an exterior hit with \(A=B\) is permitted, and the swap then
fixes the entire record. There are no other coincidences, by the
stabilizer calculation above.

One has \(\Lambda_\rho\ge36\), so
\(\Lambda_\rho-2>1\). Thus any nonzero integer decrease of
\(A_0\) under a swap decreases the code by more than either
possible tag increment; at a fixed ordered pair tag zero precedes
tag one by exactly one. This proves each claimed smallest member
and its uniqueness as a record. Distinct integral orbit minima
cannot share a fibre: if two records have the same minimum, they
belong to the rational orbit of that minimum. Conversely every
integral member of that orbit has precisely the same computed
minimum. This proves the complete fibre. The raw-scale inverse
already established gives exactly the extra disjoint choices when
\(k\) is forgotten, and proves that no other raw preimages occur.

If a globally least hit were not minimal within its own integral
orbit, the computed smaller member would be another original hit
with a strictly smaller original code, a contradiction. All further
statements now follow from the displayed minima, the exclusion of
\(A=B\) in the first row, and the exact difference formula.
\end{proof}

\begin{theorem}[The exact global first code on the occupied residual-three domain]
\label{thm:boundary-r3-global-minimum}
Put \(a_3=3h+1=(p+3)/4\), and retain its actual prime factors.
Define the explicit prime set
\[
 \mathcal P_2(a_3)=\{\ell:\ell\text{ prime},\ \ell\mid a_3,
                                  \ \ell\equiv2\pmod3\}.
\]
The residual-three shell has an integral hit exactly on the domain
\(\mathcal P_2(a_3)\ne\varnothing\). On this precisely
specified domain let \(\ell_*=\min\mathcal P_2(a_3)\).
The global least hit in each of the two original radices is
\[
 (A_0,j,B_0,\varepsilon)=(1,0,\ell_*,0),\qquad
 c_{*,H}=c_{*,F}=2(\ell_*-1),
\]
with every original coordinate given by
\[
 D_0=a_3/\ell_*,\quad C_0=(1+p\ell_*)/3,\quad
 u=a_3\ell_*,\quad R=3,\quad S=pa_3,
\]
\[
 (x,y,z)=(a_3,C_0D_0,p\ell_*C_0D_0),\qquad
 d_y=D_0,\quad d_z=p^2\ell_*^2D_0.
\]
Its complete raw-scale fibre is
\[
 (A,B,C,D,\varepsilon)
    =(k,k\ell_*,kC_0,D_0/k^2,0),\qquad k>0,\quad k^2\mid D_0.
\]
If \(\mathcal P_2(a_3)=\varnothing\), the entire original
\(j=0\) hit fibre is empty in both radices; the theorem makes
no assertion that the later shells are empty.

For a passing boundary source with coprime digits
\((A_0,j,B_0,\varepsilon)=(m+1,0,1,0)\), its same-tag
swap has code \(2m\), retaining
\[
 D_0=a_3/(m+1),\quad
 C_0=(m+1+p)/3,\quad C_0^{\leftrightarrow}=(1+p(m+1))/3.
\]
The global formula strengthens this to
\(c_{*,\rho}=2(\ell_*-1)\le2m\), with equality exactly
when \(m+1=\ell_*\). The swapped ordered witness is
\((a_3,C_0^{\leftrightarrow}D_0,
 p(m+1)C_0^{\leftrightarrow}D_0)\), and its source is
recovered by the proved involution with the same raw scale.
\end{theorem}
\begin{proof}
Here \(a_3\equiv1\pmod3\), so none of its prime factors is
three. For any raw hit, its coprime coordinates satisfy
\(A_0B_0\mid a_3\) and, because \(p\equiv1\pmod3\),
the exact gate \(A_0+B_0\equiv0\pmod3\), independently
of the retained tag. Both coordinates are units modulo three;
therefore at least one of \(A_0,B_0\) is two modulo three.
Its prime factorization contains a prime two modulo three: if
every factor were one, its product would be one. That prime divides
\(a_3\), proving necessity and proving that no omitted raw scale
can create a hit on the failed domain.

Conversely any \(\ell\in\mathcal P_2(a_3)\) gives
\(A_0=1,B_0=\ell,D_0=a_3/\ell\) and
\(C_0=(1+p\ell)/3\), all positive integers. Their equation
\(3C_0=1+p\ell\), together with \(4a_3=p+3\), is exactly
the original central equation. The preceding witness calculation
therefore proves it is a hit with the displayed ordering, residuals,
and full scale fibre. This establishes the occupied-domain
equivalence before making any least-code assertion.

On this domain select \(\ell_*\). Every hit with \(A_0\ge2\)
has code at least \(\Lambda_\rho\). A hit with \(A_0=1\)
and \(j\ge1\) has code at least \(2M_\rho\). The selected
hit has code
\(2(\ell_*-1)\le2(a_3-1)<2M_\rho\le\Lambda_\rho\),
where \(a_3\le M_\rho\) and \(J_\rho\ge1\).
It therefore precedes all these other digits. The remaining hits
have \(A_0=1,j=0\), hence \(B_0\mid a_3\) and
\(B_0\equiv2\pmod3\). Such a positive integer has at least
one prime factor in \(\mathcal P_2(a_3)\), so
\(B_0\ge\ell_*\). Equality at the smallest code is attained
with tag zero, and tag one is one integer later. Thus no code
precedes \(2(\ell_*-1)\), and the decoder makes its digits
unique. This proves the actual global minimum in both rectangles,
including comparison with every later shell.

For the boundary source, its budget says \(m+1\mid a_3\)
and its gate says \(m+1\equiv2\pmod3\). In particular
\(m+1\ge2\), and it has a prime factor at least \(\ell_*\)
and at most \(m+1\). The shell has \(3\mid p-1\), so its
integral swap and both coefficient formulas follow from the
four-candidate computation. Substitution into the unchanged radix
gives the exact swapped code \(2m\). The global formula then
gives the inequality, and equality of the two displayed integer
codes is equivalent to \(\ell_*=m+1\). The witness and inverse
are the specified candidate formulas and preserve every raw scale.
\end{proof}

\paragraph{Two exact residual-seven companions.}
For \(R=7\), the exterior same-tag integral restriction is
precisely \(p\equiv1\) or \(6\pmod7\): the prime seven
divides \((p-1)(p+1)\) exactly when it divides one of the two
factors. The tag-changing restriction is precisely
\(p\equiv1\pmod7\). The maps and fibres in the preceding
theorems apply on each of these exact residue domains.
At \(p=373\), \(a=95\), \(R=7\), the two original exterior
records \((A,B,C,D,\varepsilon)=(19,1,56,5,0)\) and
\((5,1,54,19,0)\) have \(p\equiv2\pmod7\).
Their same-tag swaps have respectively
\[
 C^{\leftrightarrow}=7088/7,\quad
 (x',y',z')=(95,35440/7,251163280/7),
\]
and
\[
 C^{\leftrightarrow}=1866/7,\quad
 (x',y',z')=(95,35454/7,66121710/7).
\]
These follow by substituting the two retained source records in
\(C^{\leftrightarrow}=(B+pA)/R\),
\(y'=BC^{\leftrightarrow}D\), and
\(z'=pAC^{\leftrightarrow}D\). Since
\(\gcd(7,p^2-1)=\gcd(7,3)=1\), the theorem proves that
every displayed coefficient and denominator has exact denominator
seven and that both rational witnesses have reciprocal sum \(4/373\).
The tag-changing candidates also have denominator seven because
\(\gcd(7,p-1)=1\). Thus the complete integral orbit of each
of these exterior records is its singleton, while all three other
labelled candidates and their inverse maps remain fully specified.
For completeness, trial division proves the asserted primality of
373: the remainders at \(2,3,5,7,11,13,17,19\) are respectively
\(1,1,3,2,10,9,16,12\), all nonzero; these are all primes below
\(\sqrt{373}<20\).


\section{Exact cross-shell transport at a fixed rational ratio}

For each \(a\), the complete positive rational shell is parametrized
by \(\tau\in\Q_{>0}\):
\[
 y_a(\tau)=\frac{pa(1+\tau)}{R_a},\qquad
 z_a(\tau)=\frac{pa(1+\tau)}{R_a\tau}.
\]
Substitution proves the original reciprocal identity, and
\(\tau=y_a/z_a\) is the inverse coordinate.  For \(a,b\in A_p\)
this gives the bijection of rational shells
\begin{equation}\label{eq:rational-transport}
 T_{a,b}(a,y,z)=
 \left(b,\frac{bR_a}{aR_b}y,\frac{bR_a}{aR_b}z\right).
\end{equation}
Its inverse is \(T_{b,a}\), and \(T_{b,c}T_{a,b}=T_{a,c}\).
All constants remain in this expression.

\begin{theorem}[Complete integral domain of ratio-preserving transport]\label{thm:ratio}
Let \(a,b\in A_p\), \(g=\gcd(a,b)\), and \(L=\lcm(R_a,R_b)\).
The exact domain on which \eqref{eq:rational-transport} carries a
positive integer witness to a positive integer witness consists of
\begin{equation}\label{eq:common-pairs}
 \{(m,n):m,n\mid pg,\ \gcd(m,n)=1,\ L\mid m+n\},
\end{equation}
with the source reconstructed by \eqref{eq:reconstruct}.
Its image is given by the same reduced pair at \(b\).
The restricted map is a bijection between these source and target
subsets, with the same inverse and one-element fibres.

If \(b=a+d\) for \(d>0\), the exact retained budgets and modulus
compatibility are
\[
 g=\gcd(a,d),\qquad
 v_\ell(pg)=\min(v_\ell(pa),v_\ell(pb)),\qquad
 \gcd(R_a,R_b)=\gcd(R_a,d).
\]
\end{theorem}
\begin{proof}
The ratio \(\tau=y/z\) is unchanged by the displayed nonzero
rational factor, so it has the same unique reduced positive pair
at both ends.  Lemma~\ref{lem:units} makes simultaneous integrality
equivalent to \(m,n\) dividing both \(pa,pb\) and both moduli
dividing \(m+n\).  Since \(a,b<p\) and \(p\) is prime,
\(\gcd(pa,pb)=p\gcd(a,b)=pg\).  This gives exactly
\eqref{eq:common-pairs} and proves both directions.  For such
a pair put \(k_a=(m+n)/R_a\), \(k_b=(m+n)/R_b\).
Their reconstructed denominators satisfy
\[
 \frac{(pb/n)k_b}{(pa/n)k_a}
  =\frac{bR_a}{aR_b}
  =\frac{(pb/m)k_b}{(pa/m)k_a},
\]
which proves the claimed image without assuming a rational
factor is an integer.  The inverse and fibres follow by
recovering the unchanged ratio.  Finally
\(\gcd(a,a+d)=\gcd(a,d)\), while
\[
 \gcd(R_a,R_b)=\gcd(R_a,4d)=\gcd(R_a,d)
\]
because \(R_a\) is odd.  Unique factorization gives the
coordinatewise minimum of every retained exponent budget.
\end{proof}

\begin{corollary}[All adjacent persistent ratios]\label{cor:adjacent}
Suppose \(a,a+1\in A_p\) and put \(R=R_a\).  The common integer
witness domain of \(T_{a,a+1}\) is empty unless
\begin{equation}\label{eq:persistent}
 R(R+4)\mid p+1.
\end{equation}
When \eqref{eq:persistent} holds, it consists of exactly the two
ordered pairs \((p,1),(1,p)\).  Thus there is exactly one
unordered common ratio pair.  For \((p,1)\), the two triples are
\[
 (a,pa(p+1)/R,a(p+1)/R),\qquad
 (a+1,p(a+1)(p+1)/(R+4),(a+1)(p+1)/(R+4)).
\]
Each quotient here is an integer by \eqref{eq:persistent}.
\end{corollary}
\begin{proof}
The preceding theorem gives \(g=1\) and
\(\gcd(R,R+4)=1\), so the common budget is \(\Div(p)\).
The full list of coprime ordered pairs is
\((1,1),(p,1),(1,p)\).  The first has sum two and cannot pass
either modulus, which is at least three.  Each of the other
two has sum \(p+1\) and passes both precisely under
\eqref{eq:persistent}.  Reconstruction gives the displayed
triples and their complements.
\end{proof}

In particular there is a sharp height restriction:
an adjacent persistent ratio can occur only if
\[
 R(R+4)\le p+1.
\]
This is a restriction on this exact transport, not on the
existence of other witnesses in either shell.  At \(p=769\),
\(a=194,R=7\), the next shell is \(195,11\), and
\(p+1=770=10\cdot7\cdot11\); hence both ratios do persist.
Primality of 769 is verified by trial division by
\(2,3,5,7,11,13,17,19,23\), all primes at most its square root.

\section{A primitive Euclidean successor with its full arithmetic domain}

Use the two original positive unimodular shears
\[
 L=\begin{pmatrix}1&1\\0&1\end{pmatrix},\qquad
 J=\begin{pmatrix}1&0\\1&1\end{pmatrix}.
\]
They act on the column \((m,n)^{\mathsf T}\) as
\(L(m,n)=(m+n,n)\) and \(J(m,n)=(m,m+n)\).  Their integer
inverses are respectively
\(\left(\begin{smallmatrix}1&-1\\0&1\end{smallmatrix}\right)\)
and \(\left(\begin{smallmatrix}1&0\\-1&1\end{smallmatrix}\right)\).
Both preserve primitivity, since an integer divides
\(m+n,n\) exactly when it divides \(m,n\), and likewise for \(J\).

\begin{proposition}[No same-shell arithmetic shear]\label{prop:no-shear}
Neither \(L\) nor \(J\) carries a pair in \(\mathcal W_a\)
into the primitive divisor supply \(\mathcal P(S_a)\) of the
same shell.
\end{proposition}
\begin{proof}
The new coordinate is \(m+n>0\), which is divisible by \(R_a\).
If it divided \(S_a\), then \(R_a\) would divide \(S_a\),
contradicting \(\gcd(R_a,S_a)=1\) and \(R_a\ge3\).
There is no further gcd cancellation, because the new pair
is primitive.  This proves the assertion for both maps.
\end{proof}

\begin{theorem}[Complete adjacent arithmetic shear classification]\label{thm:shear}
Let \(a,a+1\in A_p\), \(R=R_a\), and \((m,n)\in\mathcal W_a\).
Then \(L(m,n)\in\mathcal W_{a+1}\) exactly for the following data:
\begin{equation}\label{eq:shear-domain}
 n=1,\qquad m\mid a,\qquad m+1\mid a+1,\qquad
 R\mid m+1,\qquad R+4\mid m+2.
\end{equation}
All such data have the unique parametrization
\begin{align}
 m&=\frac{R(R+5)}4-1+R(R+4)t,\qquad t\in\Z_{\ge0},
                    \label{eq:shear-m}\\
 a&=m+c\,m(m+1),\qquad c\in\Z_{\ge0},\qquad
 p=4a-R,              \label{eq:shear-ap}
\end{align}
retaining \(R\ge3,R\equiv3\pmod4\), the original prime condition
\(p=12h+1\), and \(a,a+1\in A_p\).
Conversely every datum with these stated arithmetic properties
gives the adjacent shear and actual positive integer triples by
\eqref{eq:reconstruct}.  The classification for \(J\) is obtained
by interchanging \(m,n\); its unchanged coordinate must be one.
\end{theorem}
\begin{proof}
Suppose first that both pairs pass.  The unchanged integer \(n\)
divides both \(pa,p(a+1)\), whose greatest common divisor is \(p\).
Thus \(n=1\) or \(p\).  If \(n=p\), primitivity gives \(p\nmid m\);
then \(m+p\), also prime to \(p\), would divide \(a+1\).
But \(m+p>p>a+1\), impossible.  Hence \(n=1\).
If \(p\mid m\), the new coordinate \(m+1\) is prime to \(p\)
and greater than \(p\); once again its divisibility in \(p(a+1)\)
would force it to divide the smaller \(a+1\).  Thus \(p\nmid m\),
and \(m\mid pa\) implies \(m\mid a\).  Now
\(m+1\le a+1<p\), so \(p\nmid m+1\), and its target supply
condition gives \(m+1\mid a+1\).  The two residual conditions
are exactly the last two in \eqref{eq:shear-domain}.
Conversely \eqref{eq:shear-domain} supplies both divisor
conditions, primitivity, positivity, and the two target
congruences.  This proves the equivalence.

For the complete parameter calculation, put \(m+1=Rk\).
The next congruence is \(Rk+1\equiv0\pmod{R+4}\), or
\(4k\equiv1\pmod{R+4}\).  Because \(R+5\) is divisible by
four, the least positive inverse of four modulo \(R+4\) is
\((R+5)/4\), giving
\(k=(R+5)/4+(R+4)t\).  Since \(m>0\), no negative \(t\) is
possible: decreasing the least positive value by \(R+4\)
makes \(k<0\).  This is \eqref{eq:shear-m}, and the numerator
there was already proved divisible by four.
Next write \(a=mr\).  The condition \(m+1\mid a+1\) says
\(-r+1\equiv0\pmod{m+1}\), so \(r=1+c(m+1)\).
Positivity of \(r\) makes \(c\ge0\); this is
\eqref{eq:shear-ap}.  Every division and reverse substitution
has now been checked.  Finally the coordinate exchange
\(P(m,n)=(n,m)\) satisfies \(J=PLP\) and preserves the source
and target conditions, proving the last assertion with the
denominators exchanged as in Lemma~\ref{lem:units}.
\end{proof}

\begin{corollary}[An explicit arithmetic successor family]\label{cor:family}
For every \(c\in\Z_{\ge0}\) for which
\(p=73+1680c\) is prime, put \(a=20+420c\).
There is a passing pair \((20,1)\) at \(a,R_a=7\),
and \(L(20,1)=(21,1)\) is a passing pair at \(a+1,R_{a+1}=11\).
In particular the shears do sometimes give actual adjacent
integer transport even though Proposition~\ref{prop:no-shear}
excludes same-shell transport.
\end{corollary}
\begin{proof}
Here \(p=12(6+140c)+1\), so
\[
 3h+1=19+420c,\qquad 9h=54+1260c.
\]
Both \(a=20+420c\) and \(a+1\) belong to this interval.
One has \(20\mid a=20(1+21c)\) and
\(21\mid a+1=21(1+20c)\), while the two coordinate sums are
\(21=3\cdot7\) and \(22=2\cdot11\).
This proves all the conditions in \eqref{eq:shear-domain},
with the stated prime restriction retained.  At \(c=0\),
trial division by \(2,3,5,7\) proves that \(73\) is prime, and
the actual triples are
\[
 (20,4380,219),\qquad(21,3066,146).
\]
Their reciprocal identities follow directly from
Lemma~\ref{lem:units}, or by substituting \(S_a=1460,k_a=3\)
and \(S_{a+1}=1533,k_{a+1}=2\) in \eqref{eq:reconstruct}.
\end{proof}

\begin{theorem}[No two consecutive positive Euclidean shell steps]
\label{thm:no-two-shears}
For the fixed prime \(p=12h+1\), form the directed graph whose
vertices are the fully labelled pairs \((a,m,n)\) with
\(a\in A_p\), \((m,n)\in\mathcal W_a\), and whose edges
are exactly the passing transitions
\[
 (a,m,n)\longmapsto(a+1,m+n,n)
 \quad\hbox{or}\quad
 (a,m,n)\longmapsto(a+1,m,m+n).
\]
There is no directed path of length two in this graph.
More precisely, each nontrivial connected component of the
underlying undirected graph consists of exactly one directed edge.
There do exist edges, for example the one at \(p=73,a=20\)
in Corollary~\ref{cor:family}.
\end{theorem}
\begin{proof}
Theorem~\ref{thm:shear} shows that an \(L\)-edge must have
source \((m,1)\) and target \((m+1,1)\).  A following
\(J\)-edge would require the unchanged first coordinate
of that target to equal one; this is impossible because
\(m+1>1\).  Interchanging coordinates excludes a \(J\)
edge followed by an \(L\)-edge.

Suppose instead that there were two \(L\)-edges.  Their
three passing pairs would be \((m+i,1)\) at shells
\(a+i\), for \(i=0,1,2\), with all three shells in
the original allowed interval.  Put \(K=a-m\).
The full shear classification gives
\[
 m+i\mid a+i,\qquad R_a+4i\mid m+i+1
 \quad(i=0,1,2).
\]
Subtracting \(m+i\) in the first condition proves
\(m+i\mid K\).  One of the three consecutive positive
integers \(m,m+1,m+2\) is divisible by three, so
\(3\mid K\).  The second condition implies
\[
 R_a+4i\mid 4(m+i+1)-(R_a+4i)
                 =4(m+1)-R_a=:C .
\]
One of the three residues \(R_a,R_a+4,R_a+8\) is
divisible by three, since their increment is one
modulo three.  Hence \(3\mid C\).  But the unchanged
prime and shell identity give
\[
 C=4m+4-(4a-p)=p+4-4K\equiv2\pmod3,
\]
using \(p\equiv1\pmod3\) and \(3\mid K\).
This contradiction excludes two \(L\)-edges.
The coordinate swap preserves every divisor condition,
sum congruence, and shell label, so it excludes two
\(J\)-edges by the same argument.  All four possible
two-edge types have now been excluded.  The verified
\(p=73\) edge proves the stated sharpness.
For the stronger component assertion, an outgoing \(L\)-edge
requires \(n=1\), and an outgoing \(J\)-edge requires \(m=1\).
Both could occur only at \((1,1)\), which cannot pass the
gate \(R_a\mid2\), since \(R_a\ge3\).  Each individual
map has a unique image, so every outdegree is at most one.
An incoming \(L\)-edge has target first coordinate strictly
greater than its second, whereas an incoming \(J\)-edge has
target second coordinate strictly greater than its first.
Their inverse subtractions are unique.  Thus every
indegree is at most one.  The absence of a length-two
path also excludes a vertex having both a predecessor
and a successor.  Every nonisolated vertex therefore
belongs to precisely one edge, proving the assertion.
\end{proof}

This graph retains a passing arithmetic witness at every
intermediate step.  The theorem makes no assertion about
matrix words whose intermediate pairs leave that domain,
or about the different fixed-denominator successor below.

The family is stated for every prime specialization; this proof
does not assert that a progression supplies primes by an elementary
local lifting argument.  The successor classification is additional
cross-shell information: it names every permitted input, every
output, and its inverse on its image.  It supplies no new witness
from an empty input shell.  In particular a height based only on
the determinant of an orchard subdivision cannot infer that an
empty leaf has a permitted arithmetic successor.

\section{Arithmetic localization along the original norm subdivision}

We now retain the colleague's original norm endpoints and its corrected
subdivision together with actual marked return matrices.  This gives a
stronger inheritance statement than the mere inclusions of child
return sets in the parent.

Let \(F(r,s)=r^2+2s^2\), and retain positive primitive
\[
 v_0=(u,v),\quad w_0=(x,y),\quad
 F(v_0)=p,\quad F(w_0)=p+2,\quad
 \Delta=uy-vx,\quad T=ux+2vy .
\]
The assignment of the two norms and the sign of \(\Delta\) are never
changed.  Expansion gives
\[
 T^2+2\Delta^2
 =u^2x^2+4uvxy+4v^2y^2
   +2u^2y^2-4uvxy+2v^2x^2
 =(u^2+2v^2)(x^2+2y^2)=p(p+2).
\]
If \(\Delta=0\), the two vectors are positive rational multiples;
integrality and the gcd-one hypotheses force this multiple to be
one by reducing its numerator and denominator.  Their norms would
then agree, impossible.  Thus \(|\Delta|\ge1\) and
\(T^2\le(p+1)^2-3\).  The integers \(p+1-T,p+1+T\)
are positive and their product is \(1+2\Delta^2\).
If \(|\Delta|=1\), the product is three; its two positive factors
have sum four, which would give \(p=1\).  Hence
\(|\Delta|\ge2\).  These facts concern the original norm
endpoints, not new vectors inserted into their cone.

For the oriented cone generators use \(V_0=v_0,W_0=w_0\)
when \(\Delta>0\) and \(V_0=w_0,W_0=v_0\) when \(\Delta<0\).
This records an order change, while retaining \(v_0,w_0\) and
their respective norm labels.  Write
\[
 C(V,W)=\{rV+sW:r,s\in\mathbb R_{\ge0}\}.
\]
All generators below are primitive vectors with strictly positive
integer coordinates.  They satisfy
\[
 \det(V,W)>0 .
\]
Primitivity is an explicit hypothesis of the following subdivision lemma.

\begin{lemma}[Canonical subdivision with coordinate bounds]\label{lem:split}
For \(D=\det(V,W)>1\), there are unique integers \(r\) and
integer vector \(E\) satisfying
\[
 1\le r<D,\qquad W=rV+DE,\qquad\det(V,E)=1.
\]
The vector
\[
 Z=V+E=\frac{(D-r)V+W}{D}
\]
is positive, primitive, and interior to \(C(V,W)\), with
\(\det(V,Z)=1\), \(\det(Z,W)=D-r\).
Recursion on the right child terminates in at most \(D\)
unimodular leaves.  In every coordinate all inserted vectors
are at most the maximum of that coordinate on the initial
two generators.
\end{lemma}
\begin{proof}
Bezout gives \(E_0\in\Z^2\) with \(\det(V,E_0)=1\).
The integral basis \(V,E_0\) writes \(W=r_0 V+D E_0\).
For a unique \(k\), the replacement \(E=E_0+kV\),
\(r=r_0-Dk\) gives \(0\le r<D\).
Primitivity of \(W\) in this integral basis implies
\(\gcd(r,D)=1\); as \(D>1\), \(r\ne0\).
The formula for \(Z\) is a positive combination with both
coefficients nonzero.  It proves positivity and interiority.
Its determinant with \(V\) is one, proving primitivity.
The other determinant is
\(\det(V+E,rV+DE)=D-r\), by bilinearity and
\(\det(E,V)=-1\).
Any other \(E\) with determinant one differs from \(E_0\)
by an integer multiple of \(V\), so the specified remainder
also proves uniqueness.  Both child determinants are
smaller positive integers, hence the process terminates.
Inductively the number of leaves is at most
\(1+(D-r)\le D\), with one leaf at \(D=1\).
Finally the two coefficients of \(V,W\) in \(Z\) sum to
\((D-r+1)/D\le1\).  Each coordinate of \(Z\) is therefore
at most the corresponding maximum of \(V,W\).  Induction
proves the assertion on every later child.
\end{proof}

For a tag \(e\in\{1,p\}\), a tagged marked return in \(C\)
is the full tuple \((q,k,t,b,e)\) of positive integers
with its retained tag such that
\begin{equation}\label{eq:marked-return}
 M=\begin{pmatrix}q&p\\k&t\end{pmatrix},\qquad
 qt-pk=1,\qquad q+e=4bk,\qquad (q,k),(p,t)\in C .
\end{equation}
Here \(q\) is a positive integer, with no primality
assumption.  The tag is retained even when both marks
hold for the same matrix.

\begin{lemma}[Exact finite test and return to original denominators]
\label{lem:marked-finite}
Suppose \(C\) has slope interval
\(0<\beta\le r/s\le\alpha<\infty\).
All tuples \eqref{eq:marked-return} are obtained exactly
once by choosing
\[
 \lceil p/\alpha\rceil\le t\le\lfloor p/\beta\rfloor,\qquad
 k\mid et+1,\quad k>0,\quad \delta=(et+1)/k,
\]
testing
\[
 \delta\equiv-p\pmod{4t},\qquad
 \beta\le (pk+1)/(tk)\le\alpha,
\]
and setting \(b=(p+\delta)/(4t)\),
\(q=4bk-e=(pk+1)/t\).
Every quotient in this reconstruction is integral.
For \(e=p\) the original integer witness is
\((bk,bt,pbkt)\); for \(e=1\) it is
\((bt,pbk,pbkt)\).
\end{lemma}
\begin{proof}
The second column gives the stated finite interval for \(t\).
Substitution of \(q=4bk-e\) in the determinant equation gives
\((4bt-p)k=et+1>0\), so \(\delta=4bt-p>0\) and the
divisor and congruence conditions hold.  The first column
gives the other slope test.
Conversely the congruence makes \(b\) a positive integer,
and
\[
 t(4bk-e)=(p+\delta)k-et=pk+1
\]
proves the positive integrality of \(q\), the determinant,
and the mark.  Both slope tests place the columns in \(C\).
The constructions recover \(e,t,k,b,q\) exactly and are
inverse.  For the asserted triples, their common positive
denominator \(pbkt\) makes the two identities respectively
\[
 4bkt=pt+pk+1\quad(e=p),\qquad
 4bkt=pk+t+1\quad(e=1),
\]
which are precisely the determinant equations after the
mark is substituted.  Thus they are actual original
unit-fraction witnesses.
\end{proof}

\begin{theorem}[Disjoint arithmetic localization in the norm subdivision]
\label{thm:orchard-partition}
At every split of Lemma~\ref{lem:split} starting from
the original norm cone \(C(V_0,W_0)\), the set of tagged
returns is the disjoint union of the two child return sets.
Consequently it is exactly the disjoint union over all
terminal leaves of this canonical subdivision.
The maps keep \(q,p,k,t,b,e\) unchanged, so the integer
witness reconstructed in Lemma~\ref{lem:marked-finite}
and every divisor-budget mark attached to that witness
are unchanged.  The inverse leaf assignment is unique.
\end{theorem}
\begin{proof}
First consider any split ray \(Z=(c,d)\) whose first
coordinate satisfies \(c\le p\).  Since
\(qt-pk=1\), the first return column has strictly larger
slope than the second:
\[
 \frac qk-\frac pt=\frac1{kt}>0.
\]
If the columns were in opposite strict interiors of
the two child cones, they would therefore straddle \(Z\)
in this order.  Because their matrix has determinant
one, the exact coordinates of \(Z\) in that integral
basis are
\[
 Z=A(q,k)+B(p,t),\qquad
 A=ct-dp,\qquad B=qd-kc .
\]
Strict straddling says \(A>0,B>0\).  Both are integers,
so \(A,B\ge1\), and the first coordinate gives
\(c=Aq+Bp\ge q+p>p\), a contradiction.
Thus both columns lie in one closed child cone, including
all cases where one lies on the splitter.  A return in
both closed child cones would have both columns on their
intersection ray, contradicting determinant one.
This proves disjoint union for such a split, not just
the inclusions from either child.

It remains to prove this numerical hypothesis for the
entire original construction.  The norm equations and
strictly positive second coordinates give
\[
 u^2<p,\qquad x^2<p+2 .
\]
For \(p>2\), \(\sqrt{p+2}<p\), since
\(p^2-p-2=(p-2)(p+1)>0\).  Both initial first coordinates
are therefore less than \(p\).  The coordinate bound of
Lemma~\ref{lem:split} makes every later \(c\) less than
\(p\) too.  Repeated disjoint union proves the theorem,
as the subdivision terminates.  Each inclusion is the
identity on the full return tuple; its reconstruction
therefore retains precisely the same integer denominators.
Any added mark that is a predicate on that retained
tuple restricts the disjoint union to a disjoint union.
Given a parent return, at each split the preceding
argument identifies its unique child, giving the
unique inverse assignment all the way to a leaf.
\end{proof}

This theorem composes the original norm inequalities,
the corrected coordinate-preserving subdivision, and
the integral determinant-one return arithmetic.
It also explains why geometric inheritance alone
does not force all leaves occupied: disjoint
localization preserves the parent return count, rather
than creating a return in a prescribed empty leaf.

For the supplied \(p=1201\) example retain
\[
 v_0=(7,24),\quad w_0=(31,11),\quad
 \Delta=-667,\quad T=745,\quad
 p+1-T=457,\quad p+1+T=1947.
\]
Their norms and all displayed integers are verified by
squaring and multiplication.  The oriented canonical
chain from \((31,11)\) to \((7,24)\), obtained by the
explicit algorithm, is
\begin{align*}
 &(31,11),(14,5),(11,4),(8,3),(5,2),(2,1),\\
 &(1,1),(1,2),(1,3),(3,10),(5,17),(7,24).
\end{align*}
Each consecutive determinant is one; its exact split
remainders are recorded by the accompanying checker.
The supplied tagged return
\[
 M=\begin{pmatrix}23&1201\\9&470\end{pmatrix},\qquad
 e=p,\quad b=34,
\]
satisfies \(23\cdot470-1201\cdot9=1\) and
\(23+1201=4\cdot34\cdot9\).
Its unique leaf is \(C((8,3),(5,2))\), as shown by
the full integer coordinates
\[
 (23,9)=(8,3)+3(5,2),\qquad
 (1201,470)=52(8,3)+157(5,2),\qquad
 1\cdot157-3\cdot52=1.
\]
The reconstructed original triple is
\[
 (306,15980,172727820).
\]
The contained empty cone in the supplied correction,
generated by \((24,35),(13,19)\), is a different
unimodular subcone; it is not a leaf of this canonical
chain.  In particular its second coordinate \(35\)
exceeds the canonical invariant bound \(24\).
There is no incompatibility between these facts.

For completeness its full empty-return proof is short.
Its slope interval is
\(\beta=13/19,\alpha=24/35\), and its determinant is one.
Its containment in the original norm cone follows from
\[
 (24,35)=\frac{821v_0+331w_0}{667},\qquad
 (13,19)=\frac{446v_0+179w_0}{667},
\]
with positive coefficients.  From \(q+e=4bk\),
\(b\ge1\), and \(q\le\alpha k\) one obtains
\((4-\alpha)k\le e\).
The \(e=1\) bound is less than one, so that tag is
impossible.  For \(e=p\), the full bounds are
\[
 1752\le t\le1755,\qquad
 k\le\left\lfloor42035/116\right\rfloor=362.
\]
The determinant congruence has the following least
positive solutions:
\[
 \begin{array}{c|r|r}
 t&k_0&(1201k_0+1)/t\\\hline
 1752&407&279\\
 1753&759&520\\
 1754&1237&847\\
 1755&434&297
 \end{array}.
\]
Each equality proves that \(1201\) is a unit modulo
\(t\), and hence that every solution is
\(k=k_0+tz\), \(z\in\Z\).
The indicated \(k_0\) is between one and \(t-1\) and
exceeds 362 in every row.  No positive solution
satisfies the full bound.  Thus the entire local
return set is empty, while the distinct canonical
leaf above retains the given witness.


\paragraph{Divisor-category lineage.}
The precise source for the next construction is
Connes--Consani \cite[Section 7.4, Lemma 7.7 and Proposition 7.8]{CCPericyclic}.
The finite shell specialization, its primitive-pair return and its corrections
are the calculations proved below.
% Wrapper-free contribution. Requires amsmath, amssymb, amsthm and theorem environments.
% All mathematical assertions below are proved in this file.
\section{Embedded arithmetic-site points and common primitive shell pairs}
\label{sec:connes-primitive}

\paragraph{Primary source and precise specialization.}
Alain Connes and Caterina Consani, \emph{Cyclic theory and the pericyclic
category}, in \emph{Cyclic Cohomology at 40: Achievements and Future Prospects},
Proceedings of Symposia in Pure Mathematics 105 (2023), pp.~103--122,
introduce the divisor category $\mathcal R$ in \S7.4.
The chapter DOI is \texttt{10.1090/pspum/105/01897}.
The category has one morphism $n\to m$ when $m\mid n$.
Their Lemma~7.7, printed p.~120, identifies its topos points with embedded
groups $\mathbb Z\subset H\subset\mathbb Q$; Proposition~7.8, printed
p.~121, calculates the further passage to an abstract ordered group.
The supplied volume's corresponding PDF pages are 131 and 132.
Both complete Connes--Consani chapters, printed pp.~83--102 and 103--122,
were read before making this selection.  The following is a finite
specialization with an explicit primitive-pair return, proved independently
here.  No claim that the source already contains the shell theorem is made.

For a positive integer $N=\prod_\ell\ell^{e_\ell}$ start with the actual
finite box
\[
 \mathcal B_N=\prod_{\ell\mid N}\{0,\ldots,e_\ell\},\qquad
 \mathcal D_N=\left\{\prod_{\ell\mid N}\ell^{j_\ell}:
             (j_\ell)\in\mathcal B_N\right\}.
\]
Thus $\mathcal D_N$ is precisely the positive divisors of $N$.
The functor $F_N:\mathcal R\to\mathrm{Sets}$ sends $n$ to one point
when $n\mid N$ and to the empty set otherwise.  It is well defined by
transitivity of divisibility.  Its category of elements is nonempty,
and the least common multiple of two divisors of $N$ is a divisor of $N$
mapping to both; there are no distinct parallel arrows.  These are the
filtering conditions used in the source's Lemma~7.7.  Its associated
embedded group is
\[
 H_N=\bigcup_{n\in\mathcal D_N}\frac1n\mathbb Z\subset\mathbb Q.
\]
This is a specified construction from the box, with an exact return map;
the box is not discarded in favour of an unmarked group.

\begin{proposition}[Finite embedded points and exact divisor return]
\label{prop:connes-finite-return}
One has
\[
 H_N=\frac1N\mathbb Z,\qquad
 \mathcal D_N=\{d\in\mathbb Z_{>0}:1/d\in H_N\}.
\]
For positive $N,M$, if $g=\gcd(N,M)$ and $L=\operatorname{lcm}(N,M)$,
then
\[
 H_N\cap H_M=H_g,\qquad H_N+H_M=H_L.
\]
In particular there is an exact sequence of abelian groups
\[
 0\longrightarrow H_g\xrightarrow{t\mapsto(t,-t)}H_N\oplus H_M
 \xrightarrow{(x,y)\mapsto x+y}H_L\longrightarrow0.
\]
Every fibre of the last arrow is a translate of the displayed image of
$H_g$.  Also $H_N\subset H_M$ holds exactly when $N\mid M$.
\end{proposition}
\begin{proof}
For $n\mid N$, $1/n=(N/n)/N$, and $n=N$ occurs in the union.
Thus $H_N=(1/N)\mathbb Z$.  The equality $1/d=k/N$ with $k\in\mathbb Z$
is equivalent to $dk=N$, and for $d>0$ it is equivalent to $d\mid N$.
This proves the exact return, including all prime valuations and their bounds.
An $x\in H_N\cap H_M$ has $Nx,Mx\in\mathbb Z$.
Choose integers $r,s$ with $rN+sM=g$; then $gx\in\mathbb Z$, proving
$x\in H_g$.  Conversely $g\mid N,M$ implies $H_g\subset H_N\cap H_M$.
The sum is generated by $1/N=(M/g)/L$ and $1/M=(N/g)/L$.
Since $\gcd(M/g,N/g)=1$, a Bezout combination gives $1/L$, proving the
sum formula.  The kernel and its fibres follow by subtracting two lifts
of the same sum: a difference in the kernel is $(t,-t)$ with
$t\in H_N\cap H_M$.  Finally $H_N\subset H_M$ implies $1/N\in H_M$,
so the exact divisor return gives $N\mid M$; the converse follows from
$1/N=(M/N)/M$.
\end{proof}


For such an embedded group define its primitive reciprocal pairs by
\[
 \mathcal P(H)=\{(1/m,1/n)\in H^2:
                  m,n\in\mathbb Z_{>0},\ \gcd(m,n)=1\}.
\]
In the original finite box this means that for each prime the two
valuations lie between zero and its exponent in $N$, and at least one
valuation is zero.  Thus the primitivity decoration retains an explicit
exclusive-use condition on each finite exponent budget.

Fix the same prime $p=12h+1$ throughout, and let
$a,b\in\{3h+1,\ldots,9h\}$.  Keep
\[
 S_a=pa,\quad R_a=4a-p,\qquad S_b=pb,\quad R_b=4b-p.
\]
For each shell $t\in\{a,b\}$ define the ordered marked pair set
\[
 \mathcal Q_t=\{(m,n)\in\mathbb Z_{>0}^2:
 m\mid S_t,\ n\mid S_t,\ \gcd(m,n)=1,\ R_t\mid m+n\}.
\]
Both $R_t$ are positive odd integers at least $3$.
Also $\gcd(R_t,S_t)=1$: a common prime divisor with $t$ would divide
$p$ but $0<t<p$; a common divisor with $p$ would divide $4t$, whereas
$p$ is an odd prime not dividing $t$.

\begin{theorem}[Embedded intersection with full shell marks]
\label{thm:connes-common-pairs}
Set
\[
 d=\gcd(a,b),\qquad L_R=\operatorname{lcm}(R_a,R_b).
\]
Then
\[
 H_{S_a}\cap H_{S_b}=H_{pd},
\]
and reciprocal embedding gives the exact bijection
\[
 \mathcal Q_a\cap\mathcal Q_b
 \longleftrightarrow
 \{(1/m,1/n)\in\mathcal P(H_{pd}):L_R\mid m+n\}.
\]
Its map is $(m,n)\mapsto(1/m,1/n)$ and its inverse is coordinatewise
reciprocal.  In integer coordinates the common pair set is precisely
\[
 \{(m,n):m\mid p\gcd(a,b),\ n\mid p\gcd(a,b),\
          \gcd(m,n)=1,\ L_R\mid m+n\}.
\]
The fibres of this correspondence are singletons.  For noncoprime shell
moduli the retained sum modulus is exactly
$L_R=R_aR_b/\gcd(R_a,R_b)$.
\end{theorem}
\begin{proof}
Since $p$ is the same factor in both $S_a,S_b$,
$\gcd(S_a,S_b)=p\gcd(a,b)=pd$.  Apply
Proposition~\ref{prop:connes-finite-return} to get the intersection.
An integer $m$ divides both $S_a,S_b$ exactly when $1/m$ lies in both
embedded groups, equivalently in $H_{pd}$; the same holds for $n$.
The primitivity condition is unchanged under this comparison because
the integer labels $m,n$ are recovered exactly by reciprocals.
For a positive integer $z$, divisibility by both $R_a$ and $R_b$ is
equivalent to divisibility by their least common multiple: for each prime
its valuation must be at least the maximum of the two exponents.
Applying this to $z=m+n$ proves the mark and the displayed formula for
$L_R$.  Both reciprocal maps undo each other, proving bijectivity and
the asserted fibres.
\end{proof}

\begin{proposition}[Explicit return of every common pair to both witnesses]
\label{prop:connes-pair-witness}
For $(m,n)$ in the common pair set and $t=a,b$, define
\[
 k_t=\frac{S_t}{mn}\frac{m+n}{R_t},\qquad
 B_t=k_tm,\quad C_t=k_tn.
\]
These are positive integers and
\[
 \frac4p=\frac1t+\frac1{B_t}+\frac1{C_t},\qquad
 \frac{B_t}{C_t}=\frac mn.
\]
Conversely every ordered positive integral denominator pair $(B_t,C_t)$
at shell $t$ determines exactly one pair in $\mathcal Q_t$ by
\[
 k_t=\gcd(B_t,C_t),\qquad m=B_t/k_t,\quad n=C_t/k_t.
\]
Thus an ordered rational ratio occurs at both shells exactly when its
unique primitive pair lies in the common set of
Theorem~\ref{thm:connes-common-pairs}.  The two shell witness pairs for
such a ratio are then exactly the displayed formulas.
\end{proposition}
\begin{proof}
The coprimality of $m,n$ and their separate divisibility into $S_t$ give
$mn\mid S_t$: at each prime at most one of the two positive valuations
occurs, and that valuation is at most the exponent in $S_t$.
Also $R_t\mid m+n$.  Hence $k_t,B_t,C_t$ are positive integers, and
\[
 \frac1{B_t}+\frac1{C_t}
 =\frac{m+n}{k_tmn}=\frac{R_t}{S_t}=\frac4p-\frac1t.
\]
The ratio equality is immediate from the positive factor $k_t$.
Conversely the reciprocal identity gives
$R_t B_t C_t=S_t(B_t+C_t)$.  Set $k=\gcd(B_t,C_t)$,
$m=B_t/k$, and $n=C_t/k$; the definitions give $\gcd(m,n)=1$ and
\[
 R_t k mn=S_t(m+n).
\]
No prime dividing $mn$ divides $m+n$, so $mn\mid S_t$.  Dividing by
$mn$ gives $R_tk=(S_t/(mn))(m+n)$; as
$\gcd(R_t,S_t/(mn))=1$, it follows that $R_t\mid m+n$.
Thus $(m,n)\in\mathcal Q_t$, and the same equation forces exactly
$k=(S_t/(mn))((m+n)/R_t)$.
For completeness, uniqueness of the primitive pair for a positive ratio
follows as well: if $m/n=m'/n'$ with both pairs coprime, then
$mn'=m'n$ implies $m\mid m'$ and $m'\mid m$, so $m=m'$ and $n=n'$.
This proves the final assertion without identifying differently decorated
or nonprimitive presentations.
\end{proof}

\begin{corollary}[Consecutive-shell intersection]
\label{cor:connes-consecutive}
If $b=a+1$ and both shells are in the allowed interval, write $R=R_a$.
Then the two shells have a common ordered primitive pair exactly when
\[
 R(R+4)\mid p+1.
\]
If this divisibility holds their common pairs are exactly $(p,1)$ and
$(1,p)$; otherwise the common set is empty.
For $(p,1)$ the exact denominator pair at shell $t=a,a+1$ is
\[
 \left(\frac{pt(p+1)}{4t-p},\frac{t(p+1)}{4t-p}\right),
\]
and the other primitive pair interchanges its entries.
\end{corollary}
\begin{proof}
Here $\gcd(a,a+1)=1$, so the embedded intersection is $H_p$ and the
retained integer divisors are $1,p$.  Its coprime pairs are
$(1,1),(1,p),(p,1)$.  The first has sum $2$, which is not divisible by
$R\ge3$.  The other two have sum $p+1$.
Also $\gcd(R,R+4)=\gcd(R,4)=1$ since $R$ is odd.  The preceding theorem
therefore imposes exactly $R(R+4)\mid p+1$.
For $(m,n)=(p,1)$, the formula for $k_t$ becomes
$t(p+1)/(4t-p)$; multiplying by $p$ or $1$ gives the displayed pair.
\end{proof}

\begin{proposition}[What the abstract-point isomorphism retains]
\label{prop:connes-abstract-loss}
For positive $N,M$, the unique positive ordered-group isomorphism
$H_N\to H_M$ is
\[
 \phi_{N,M}(x)=\frac NM x.
\]
On reciprocal divisors its exact integral-return subset and map are
\[
 u=N/e\longmapsto v=M/e
 \qquad(e\mid\gcd(N,M)).
\]
The inverse sends $M/e$ back to $N/e$; these fibres are singletons.
All finite embedded groups $H_N$ have the same abstract ordered-group
isomorphism class, so the fibre of that passage consists of every finite
positive budget $N$.
For the shells above, applying $\phi_{pa,pb}$ coordinatewise to the
reciprocal of a primitive pair yields the reciprocal integer labels
\[
 (m',n')=\left(\frac ba m,\frac ba n\right).
\]
This is an integral primitive pair only when $a=b$.  In particular, for
two distinct shells this isomorphism cannot itself give a successor on
the actual primitive pairs, although their embedded intersection has the
complete correspondence and positive witness return just proved.
\end{proposition}
\begin{proof}
The least positive element of $H_N$ is $1/N$, and an ordered-group
isomorphism must send it to the least positive element $1/M$ of $H_M$.
Additivity forces the given formula, which is indeed an isomorphism.
For $u\mid N$ write $e=N/u$.  Then
$\phi_{N,M}(1/u)=e/M$ is a reciprocal positive integer $1/v$ exactly
when $e\mid M$, in which case $v=M/e\mid M$.
Together with $e\mid N$ this gives precisely the stated retained subset
and its inverse.  Existence of the isomorphism for every $N,M$ proves
the assertion about the isomorphism-class fibre; as this passage is a
map on embedded objects, its kernel relation is all pairs of finite
positive budgets, not the kernel of an additive homomorphism.
For shell pairs $\phi_{pa,pb}(1/m)=a/(bm)$, which gives the displayed
integer label $m'=bm/a$ if integral, and likewise for $n$.
Write $a=dk$, $b=d\ell$, where $d=\gcd(a,b)$ and $\gcd(k,\ell)=1$.
Integrality of both labels forces $k\mid m,n$.  Their original primitivity
therefore forces $k=1$.  Now $(m',n')=(\ell m,\ell n)$ has greatest
common divisor $\ell$, so its primitivity forces $\ell=1$.
Consequently $a=b=d$.  Conversely if $a=b$ the map is the identity
and retains every primitive pair.  This proves the exact obstruction for
this specific isomorphism, alongside the surviving embedded correspondence.
\end{proof}

\subsection{Composing the finite point with the middle-channel shear}
\label{subsec:connes-middle-shear}
For an allowed shell $a$ retain the full divisor set
$\mathcal D_{a^2}$ and the two marks
\[
 E_a=\{u\mid a^2:R_a\mid4u+1\},\qquad
 M_a=\{u\mid a^2:R_a\mid u+a\}.
\]
The following maps have different domains and will all be kept in the
composition: the ordered-group isomorphism on $H_{a^2}$, its reciprocal
divisor return, and a divisor-to-primitive-pair map depending on the shell.

\begin{proposition}[Full middle-channel primitive return]
\label{prop:connes-middle-return}
There is a bijection
\[
 \mu_a:\mathcal D_{a^2}\longrightarrow
 \{(m,n):m\mid a,\ n\mid a,\ \gcd(m,n)=1\},
\]
\[
 \mu_a(u)=\left(\frac{u}{\gcd(u,a)},\frac{a}{\gcd(u,a)}\right).
\]
Its inverse is $(m,n)\mapsto am/n$ and its fibres are singletons.
It preserves the exact mark by
\[
 R_a\mid u+a\quad\Longleftrightarrow\quad R_a\mid m+n.
\]
For a marked divisor its denominator return is
\[
 \left(\frac{p(a+u)}{R_a},\frac{p(a+a^2/u)}{R_a}\right),
\]
and the ordered ratio of these entries is exactly $u/a=m/n$.
\end{proposition}
\begin{proof}
Put $g=\gcd(u,a)$.  Then $m=u/g$, $n=a/g$ are coprime.
For a prime $\ell$ let $e=v_\ell(a)$ and $j=v_\ell(u)$.
Since $u\mid a^2$, $0\le j\le2e$.  The valuations of $m,n$ are,
respectively, $j-\min(j,e)$ and $e-\min(j,e)$, both in $[0,e]$;
at least one is zero.  Thus $m,n\mid a$ and the entire exponent data
are retained by the formulas.  Conversely, if $m,n\mid a$ are coprime,
$u=am/n$ is integral and $a^2/u=an/m$ is integral.  Writing $a=gn$
gives $u=gm$ and $\gcd(u,a)=g\gcd(m,n)=g$, proving both inverse identities.
Also $u+a=g(m+n)$ and $g\mid a$, so $\gcd(g,R_a)=1$; this proves
the mark equivalence.  Proposition~\ref{prop:connes-pair-witness}
applies to this primitive pair.  Its $k$ is
$pa(m+n)/(mnR_a)$, and substitution of $u=am/n$ shows that
$km=p(a+u)/R_a$ and $kn=p(a+a^2/u)/R_a$.
Thus both are positive integers giving the required shell witness.
Finally dividing the displayed entries, or using $(km)/(kn)$, gives
$m/n=u/a$ exactly.
\end{proof}

\begin{theorem}[Exact marked domain of the arithmetic-site successor]
\label{thm:connes-marked-successor}
Let $b=a+1$ be another allowed shell and put $R=4a-p$.
The ordered-group isomorphism $H_{a^2}\to H_{(a+1)^2}$ has reciprocal
divisor return only at
\[
 u=a^2\longmapsto v=(a+1)^2.
\]
Neither endpoint is $E$-marked.  Its full domain of successful $E/M$
marked returns consists exactly of the middle-channel endpoints for
which
\[
 R(R+4)\mid p+4.
\]
For every such endpoint the maps compose as
\[
 \begin{array}{ccc}
 a^2&\longmapsto&(a+1)^2\\
 \mu_a\downarrow&&\downarrow\mu_{a+1}\\
 (a,1)&\xrightarrow{\ L\ }&(a+1,1),
 \end{array}
 \qquad
 L=\begin{pmatrix}1&1\\0&1\end{pmatrix},
\]
where the lower arrow acts on column vectors.
The exact positive integer denominator map is
\[
 \left(a,\frac{pa(a+1)}R,\frac{p(a+1)}R\right)
 \longmapsto
 \left(a+1,\frac{p(a+1)(a+2)}{R+4},\frac{p(a+2)}{R+4}\right).
\]
Both triples satisfy the Erd\H{o}s--Straus equation for the unchanged
prime $p$.  The inverse on this marked image retains $p$ and sends
$(a+1,R+4,(a+1)^2)$ back to $(a,R,a^2)$; its fibres are singletons.

Equivalently, all successful parameter triples $(p,R,a)$ have the form
\[
 \begin{split}
 &R\ge3,\quad R\equiv3\pmod4,\quad t\in\mathbb Z_{\ge0},\\
 &p=R(R+4)(4t+1)-4,\qquad p\text{ prime},\qquad p\equiv1\pmod{12},\\
 &a=\frac{R(R+5)}4-1+R(R+4)t.
 \end{split}
\]
For these parameters both $a$ and $a+1$ are automatically in the
prescribed shell interval.
\end{theorem}
\begin{proof}
The reciprocal divisor return calculation in
Proposition~\ref{prop:connes-abstract-loss}, applied to
$N=a^2$, $M=(a+1)^2$ with greatest common divisor $1$, gives exactly
the top-divisor map.
We prove the $E$ assertion without dropping the sign or the factor $4$.
Any positive odd integer $R\equiv3\pmod4$ has a prime divisor
$q\equiv3\pmod4$: otherwise every prime factor would be $1$ modulo
$4$, or every occurrence of a $3$-modulo-$4$ prime would have even
exponent, and the product would be $1$ modulo $4$.
If $R\mid4a^2+1$, then modulo this $q$ the nonzero residue $x=2a$
satisfies $x^2=-1$.  The residues $1,x,-1,-x$ are distinct and form a
subgroup of order $4$ in $(\mathbb Z/q\mathbb Z)^\times$.
Its cosets partition the $q-1$ nonzero residues into sets of four,
forcing $4\mid q-1$, a contradiction.  Thus $R\nmid4a^2+1$.
The same proof with $a+1,R+4$ excludes the next $E$ mark.
The middle mark at the first endpoint is
\[
 R\mid a^2+a\quad\Longleftrightarrow\quad R\mid a+1
 \quad\Longleftrightarrow\quad R\mid p+4.
\]
The first equivalence uses $\gcd(a,R)=1$; the second uses
$4(a+1)=p+R+4$ and $\gcd(4,R)=1$.
At the next endpoint the same identities give
$R+4\mid(a+1)^2+(a+1)$ exactly when $R+4\mid a+2$,
equivalently $R+4\mid p+4$.
Since $\gcd(R,R+4)=1$, both marks hold exactly when
$R(R+4)\mid p+4$.  This calculates the successful domain; no lifting
condition is left unevaluated.
Proposition~\ref{prop:connes-middle-return} gives
$\mu_a(a^2)=(a,1)$ and $\mu_{a+1}((a+1)^2)=(a+1,1)$.
Multiplication by $L$ sends $(a,1)^{\mathsf T}$ to
$(a+1,1)^{\mathsf T}$, so the displayed composition commutes.
The denominator formulas are the middle-channel return at these two
divisors, and that proposition proves their integrality, positivity,
and reciprocal identity.  Retaining $p,a,R$ gives the stated inverse.
This composition is not the coordinatewise map $\phi_{pa,pb}$ on
primitive reciprocal pairs: it changes the ratio from $a$ to $a+1$
through the two different shell-dependent maps $\mu_a,\mu_{a+1}$.
Consequently it does not contradict the no-return assertion for that
different map in Proposition~\ref{prop:connes-abstract-loss}.

For the parameter description put $k=(p+4)/(R(R+4))$.
It is a positive integer; both $p+4$ and $R(R+4)$ are $1$ modulo
$4$, so $k\equiv1\pmod4$ and uniquely $k=4t+1$ with $t\ge0$.
Then $p=R(R+4)(4t+1)-4$ and $a=(p+R)/4$ give the displayed formula.
Conversely those formulas imply $R=4a-p$ and the successful divisibility.
Set $h=(p-1)/12$.  Since $R\ge3$ and $R\equiv3\pmod4$,
$a=(p+R)/4\ge(p+3)/4=3h+1$ and is integral.
Also
\[
 2p-(R+7)=2R(R+4)(4t+1)-R-15
 \ge2R^2+7R-15>0\qquad(R\ge3).
\]
Thus $p+R+4\le3p-3$, which gives
$a+1\le(3p-3)/4=9h$.  This proves the required interval assertions.
\end{proof}

\paragraph{Exact successful example.}
For $R=7$ and $t=0$ the parameter formula gives $p=73$, $h=6$,
$a=20$, and $a+1=21$, both in $\{19,\ldots,54\}$.
The divisor map is $400\mapsto441$, the primitive map is
$(20,1)\mapsto(21,1)$, and the two triples are
\[
 (20,4380,219)\longmapsto(21,3066,146).
\]
Indeed $p+4=77=7\cdot11$,
$400+20=420=7\cdot60$, and $441+21=462=11\cdot42$.
Their exact reciprocal identities are
\[
 \frac1{20}+\frac1{4380}+\frac1{219}=\frac4{73},\qquad
 \frac1{21}+\frac1{3066}+\frac1{146}=\frac4{73},
\]
as follows by putting the first sum over denominator $4380$,
giving $(219+1+20)/4380=240/4380=4/73$, and the second over
$3066$, giving $(146+1+21)/3066=168/3066=4/73$.


\subsection{The full shear as an affine correction of the finite-point map}

The preceding middle-endpoint map is the zero-correction part
of a fully explicit correspondence for every shear in
Theorem~\ref{thm:shear}.  We compute that correction and its
composition, keeping the finite divisor bounds.

\begin{proposition}[Affine embedded-group map and its exact divisor fibre]
\label{prop:affine-point}
For positive integers \(N,M\) and any \(c\in\Z\), define
\[
 \Psi_{N,M,c}:H_N\longrightarrow H_M,\qquad
 x\longmapsto\frac NM x-\frac cM .
\]
It is an affine bijection of the embedded sets, with inverse
\[
 y\longmapsto\frac MN y+\frac cN.
\]
Its fibres are singletons.  It is an additive homomorphism
exactly when \(c=0\); otherwise no homomorphism kernel
is being assigned to it.  Its zero fibre is \(\{c/N\}\),
and its positive-image domain is
\(\{x\in H_N:Nx>c\}\).  Its exact return on positive
reciprocal divisors is
\begin{equation}\label{eq:affine-divisor-return}
 u=N/e\longmapsto v=M/(e-c),\qquad
 e\mid N,\quad e>c,\quad e-c\mid M.
\end{equation}
Thus the retained index fibre is the finite intersection
\(\Div(N)\cap(c+\Div(M))\), with no unbounded
completion substituted for either side.  The complete
composition rule is
\[
 \Psi_{M,L,d}\circ\Psi_{N,M,c}=\Psi_{N,L,c+d}
 \qquad(N,M,L>0,\ c,d\in\Z).
\]
For compositions restricted to reciprocal divisors,
the intermediate condition \(e-c\mid M,\ e-c>0\)
remains part of the domain in addition to the two
endpoint conditions.
\end{proposition}
\begin{proof}
Write \(x=e/N\), \(e\in\Z\).  Its image is
\((e-c)/M\), which is in \(H_M\).  Conversely any
\(f/M\) is the image of \((f+c)/N\).  These formulas
prove both inverse compositions and singleton fibres.
The image of zero is \(-c/M\), so additivity
requires \(c=0\), and for \(c=0\) the formula is
an additive homomorphism.  The formula \((e-c)/M\)
also proves the zero-fibre and positive-image assertions.
For \(u\mid N\) put \(e=N/u\); it is a positive
divisor of \(N\).  The image of \(1/u\) is
\((e-c)/M\).  It is \(1/v\) with \(v\) a positive
integer precisely when \(e-c>0\) and \(e-c\mid M\).
Then \(v=M/(e-c)\mid M\), proving
\eqref{eq:affine-divisor-return} and its converse.
Finally substitution gives
\[
 \frac ML\left(\frac NM x-\frac cM\right)-\frac dL
       =\frac NLx-\frac{c+d}{L}.
\]
For a two-step reciprocal return the intermediate
point has index \(e-c\); its positive divisor
condition is exactly the one stated, independent
of whether the endpoint has index \(e-c-d\).
\end{proof}

\begin{theorem}[Exact affine-point realization of every positive shear]
\label{thm:full-connes-shear}
For any \(L\)-edge of Theorem~\ref{thm:shear}, use
its unique parameters \(m,c\):
\[
 a=m+c\,m(m+1),\qquad b=a+1,\qquad c\ge0.
\]
The source and target middle divisors are
\[
 u=am,\qquad v=b(m+1).
\]
They divide \(a^2,b^2\), respectively.  The following
maps compose exactly:
\[
 \frac1u
 \xmapsto{\ \Psi_{a^2,b^2,c}\ }\frac1v,\qquad
 \mu_a(u)=(m,1)\xmapsto{\ L\ }(m+1,1)=\mu_b(v).
\]
Their retained integral indices satisfy
\[
 e=\frac{a^2}{u}=\frac am=1+c(m+1),\qquad
 e'=\frac{b^2}{v}=\frac b{m+1}=1+cm,\qquad e-e'=c.
\]
Both middle marks, all exponent bounds, and the
positive integer denominator reconstruction are
exactly the original shear marks and reconstruction.
The uncorrected ordered-group isomorphism gives this
same reciprocal-divisor correspondence precisely
when \(c=0\).  If \(c>0\), the uncorrected map has
no positive reciprocal integer return at this source.
\end{theorem}
\begin{proof}
The source divisibility \(m\mid a\) and target
divisibility \(m+1\mid b\) imply \(u\mid a^2\)
and \(v\mid b^2\) by their displayed definitions.
Moreover
\(\gcd(am,a)=a\) and \(\gcd(b(m+1),b)=b\),
so the exact maps \(\mu_a,\mu_b\) of
Proposition~\ref{prop:connes-middle-return}
give the stated primitive pairs.
Their middle marks are
\[
 R_a\mid u+a=a(m+1)
     \ \Longleftrightarrow\ R_a\mid m+1,
\]
and
\[
 R_b\mid v+b=b(m+2)
     \ \Longleftrightarrow\ R_b\mid m+2 ,
\]
by the respective unit assertions.  These are
precisely the two gates of the full shear
classification, not additional assumptions.
The parameter formula gives both index equalities;
the equality \(e-e'=c\) then shows
\[
 \Psi_{a^2,b^2,c}(1/u)=\frac{e-c}{b^2}
                         =\frac{e'}{b^2}=\frac1v .
\]
It also explicitly proves the positive divisor
return conditions of Proposition~\ref{prop:affine-point}.
The maps \(\mu_a,\mu_b\) have already been proved
to retain their whole divisor valuations and give
the exact integer witnesses, so this composition
retains those data as well.
Without the correction, equality to \(1/v\)
would instead require \(e=e'\), equivalently \(c=0\).
This proves both the full affine statement and the
precise special case of the ordered-group
isomorphism.  When \(c>0\), the uncorrected image is
\(e/b^2\) with \(e=1+c(m+1)>1\) and \(e\mid a\).
Since \(\gcd(a,b)=1\), \(\gcd(e,b^2)=1\), so
\(e\nmid b^2\).  Its reciprocal cannot be an integer.
\end{proof}

Here is the full conjugate correspondence for \(J\),
including the different reciprocal-coordinate formula.
On the positive reciprocal-divisor set of \(N\), define
\[
 \iota_N(x)=\frac1{Nx}.
\]
For \(x=1/u\), \(u\mid N\), its value is
\(u/N=1/(N/u)\), so it implements the involutive divisor
complement \(u\mapsto N/u\).  Direct substitution gives
\begin{equation}\label{eq:conjugate-affine}
 \iota_M\circ\Psi_{N,M,c}\circ\iota_N(x)
       =\frac{x}{1-cx}.
\end{equation}
Its exact positive reciprocal-divisor domain and label map are
\[
 u\mid N,\quad u>c,\quad u-c\mid M,
 \qquad u\longmapsto v=u-c.
\]
The inverse is \(v\mapsto v+c\), restricted to
\(v\mid M\), \(v+c>0\), \(v+c\mid N\);
all fibres are singletons.  These assertions follow
also by inserting \(x=1/u\) in \eqref{eq:conjugate-affine},
which gives \(1/(u-c)\).  The condition \(u>c\)
both excludes its pole and proves positivity.

For the \(J\)-edge \((1,m)\mapsto(1,m+1)\), use the same
parameters \(a=m+c\,m(m+1)\), \(b=a+1\).
The middle divisors are now
\[
 u=\frac am=1+c(m+1),\qquad
 v=\frac b{m+1}=1+cm=u-c.
\]
They are positive divisors of \(a^2,b^2\), respectively,
so the exact domain above is satisfied with \(N=a^2,M=b^2\).
Divisor complementation satisfies
\(\mu_a(a^2/w)=P\mu_a(w)\), where
\(P(m,n)=(n,m)\): indeed the reduced ratio
\((a^2/w)/a=a/w\) is the inverse of \(w/a\).
Thus the proved conjugation is exactly \(J=PLP\).
It retains both middle congruences, their complete
valuation bounds, and exchanges the last two original
denominators under \eqref{eq:reconstruct}.

For an explicit nonzero correction retain the prime
\(p=1753\), \(a=440\), \(b=441\), \(R_a=7\),
\(R_b=11\), \(m=20\), \(c=1\).  Primality follows
by testing the primes at most \(\sqrt{1753}<42\),
namely \(2,3,5,7,11,13,17,19,23,29,31,37,41\);
none divides \(1753\).  Here the \(L\) map is
\[
 \frac1{8800}\xmapsto{\Psi_{193600,194481,1}}
 \frac1{9261},
\]
with indices \(22\mapsto21\); its original witnesses are
\[
 (440,2313960,115698)\longmapsto
 (441,1546146,73626).
\]
The conjugate \(J\) map sends \(1/22\) to \(1/21\)
and exchanges the last two entries in these triples.
All these equalities follow by substitution in the
proved formulas, which retain the original prime and shells.

The integer difference \(c\) has the additive
composition law computed above.  It is not a
global Lyapunov function for ES failure: it may
be zero on an actual edge, and
Theorem~\ref{thm:no-two-shears} proves that within
the exact positive shear domains at a fixed
prime, no two such edges can even be composed.
The full affine maps do compose on their
embedded groups; their restrictions retain
the additional intermediate divisor and
residual tests, exactly as stated.

\section{A strictly decreasing denominator with exact arithmetic return}
\subsection{An arithmetic successor preserving one original denominator}
\label{sec:exact-audit-fixed-y}

Let \(p=12h+1\) be prime, let
\(A_p=\{3h+1,\ldots,9h\}\), and choose \(a,a+1\in A_p\).
Retain the original residuals and supplies
\[
 b=a+1,\qquad R=4a-p,\qquad R_b=R+4,\qquad
 S_a=pa,\qquad S_b=p(a+1),\qquad A=a(a+1).
\]
For \(c\in\{a,b\}\), let
\[
 \mathcal I_c=\left\{(y,z)\in\mathbb Z_{>0}^2:
       \frac4p=\frac1c+\frac1y+\frac1z\right\}.
\]
Every retained coordinate, including the ordered unmarked denominators,
is part of these objects.

\begin{theorem}[Exact integral fixed-denominator successor]
\label{thm:exact-audit-fixed-y}
Define the two finite domains
\begin{align*}
 \mathcal D_a&=\{(y,z)\in\mathcal I_a:A+z\mid A^2\},\\
 \mathcal D_b^-&=\{(y,z')\in\mathcal I_b:
                       0<z'<A,\ A-z'\mid A^2\}.
\end{align*}
Then
\begin{equation}\label{eq:exact-audit-fixed-y-map}
 T_a:\mathcal D_a\longrightarrow\mathcal D_b^-,\qquad
 (y,z)\longmapsto\left(y,\frac{Az}{A+z}\right)
\end{equation}
is a bijection, with inverse
\begin{equation}\label{eq:exact-audit-fixed-y-inverse}
 (y,z')\longmapsto\left(y,\frac{Az'}{A-z'}\right).
\end{equation}
Every forward value satisfies \(0<z'<\min(z,A)\). Thus this
successor strictly decreases the original denominator \(z\) on its
complete integral-return domain; every fibre is a singleton.

For the source reduced positive primitive ratio \(\tau=m/n=y/z\),
the target ratio is
\begin{equation}\label{eq:exact-audit-fixed-y-ratio}
 \tau'=\frac{a(R+4)\tau+p}{(a+1)R}.
\end{equation}
Its raw integer pair, reduction factor, and primitive return are exactly
\begin{equation}\label{eq:exact-audit-fixed-y-gcd}
 U=a(R+4)m+pn,\quad V=(a+1)Rn,\quad
 g=\gcd(U,V),\quad (m',n')=(U/g,V/g).
\end{equation}
For every source witness, this primitive target obeys
\[
 m',n'\mid p(a+1),\qquad R+4\mid m'+n'
\]
if and only if the proved divisibility \(A+z\mid A^2\) holds.
In particular the reduction factor and all target exponent bounds are
retained, not replaced by projective-ratio existence.
\end{theorem}
\begin{proof}
Keeping the same original \(y\), subtraction of the equations for
shells \(a\) and \(b=a+1\) forces
\[
 \frac1{z'}=\frac1z+\frac1a-\frac1{a+1}
           =\frac1z+\frac1A.
\]
Its unique positive rational solution is \(z'=Az/(A+z)\).
Conversely this identity proves the original equation at \(b\),
because
\[
 \frac1b+\frac1y+\frac1{z'}
 =\frac1b+\frac1y+\frac1z+\frac1a-\frac1b=\frac4p.
\]
Since \(A,z>0\), both \(z'/z=A/(A+z)<1\) and
\(z'/A=z/(A+z)<1\), proving strict positivity and both bounds.
The exact identity
\begin{equation}\label{eq:exact-audit-fixed-y-factor}
 z'=A-\frac{A^2}{A+z},\qquad (A+z)(A-z')=A^2
\end{equation}
shows, for integer source \(z\), that \(z'\) is an integer exactly
when \(A+z\mid A^2\). In that case \(A-z'>0\) is a positive
integer divisor of \(A^2\); thus the forward point lies in
\(\mathcal D_b^-\).

Conversely, for a point in \(\mathcal D_b^-\), put
\[
 z=\frac{Az'}{A-z'}=\frac{A^2}{A-z'}-A.
\]
It is a positive integer by \(0<z'<A\) and the displayed divisor
condition. Its reciprocal satisfies \(1/z=1/z'-1/A\), so
\[
 \frac1a+\frac1y+\frac1z
 =\frac1a+\frac1y+\frac1{z'}-\frac1a+\frac1b=\frac4p.
\]
Also \(A+z=A^2/(A-z')\) divides \(A^2\). Thus it lies in
\(\mathcal D_a\). Equation \eqref{eq:exact-audit-fixed-y-factor}
proves both inverse compositions. The domains are finite: on the
source, \(A+z\) is a positive divisor of the fixed positive integer
\(A^2\), and each \(z\) fixes \(y\) uniquely through the original
equation; the target has \(1\le z'<A\), again with at most one
\(y\) for each \(z'\).

The source shell identity gives
\(y=pa(1+\tau)/R\), since \(y/z=\tau\) and
\(R/pa=1/y+1/z=(1+\tau)/y\). Consequently
\begin{align*}
 \tau'=\frac y{z'}
 &=\frac yz+\frac yA
 =\tau+\frac{p(1+\tau)}{(a+1)R}\\
 &=\frac{((a+1)R+p)\tau+p}{(a+1)R}
 =\frac{a(R+4)\tau+p}{(a+1)R},
\end{align*}
where the last equality uses the original identity \(R+p=4a\).
This proves \eqref{eq:exact-audit-fixed-y-ratio}; substitution
\(\tau=m/n\) and division by the explicitly retained \(g\) prove
\eqref{eq:exact-audit-fixed-y-gcd}. All four integers \(U,V,m',n'\)
are positive, and \(\gcd(m',n')=1\).

To verify the target budget claim without assuming any lift, recall its
elementary exact reconstruction. For an integer target witness,
\[
 d_b=(R+4)y-pb>0,\qquad e_b=(R+4)z'-pb>0,\qquad
 d_b e_b=(pb)^2.
\]
Positivity follows from \((R+4)/pb=1/y+1/z'\); the product follows
by expansion of the same identity. Moreover
\(d_b/(pb)=y/z'=m'/n'\). If \(v_\ell(pb)=e_\ell\), then
\(0\le v_\ell(d_b)\le2e_\ell\). Reducing \(d_b/(pb)\) gives
\[
 v_\ell(m')=\max(v_\ell(d_b)-e_\ell,0),\quad
 v_\ell(n')=\max(e_\ell-v_\ell(d_b),0),
\]
so both are at most \(e_\ell\). Thus \(m',n'\mid pb\).
Since \(\gcd(pb,R+4)=1\), reducing
\(d_b\equiv-pb\pmod{R+4}\) gives \(R+4\mid m'+n'\).

In the reverse direction those bounded primitive conditions reconstruct
positive integers
\[
 Y=\frac{pb}{n'}\frac{m'+n'}{R+4},\qquad
 Z=\frac{pb}{m'}\frac{m'+n'}{R+4}.
\]
Their ratio is \(m'/n'\), and their reciprocal sum is
\((R+4)/(pb)\), exactly the target rational shell. A positive
ratio \(m'/n'\) and that reciprocal sum determine the ordered rational
pair uniquely by these same formulas; hence \((Y,Z)=(y,z')\).
Thus \(z'\) is integral, which by
\eqref{eq:exact-audit-fixed-y-factor} is precisely \(A+z\mid A^2\).
This proves both directions of the budget equivalence without an
unproved arithmetic-return premise.
\end{proof}

At \(p=13,a=4,R=3\), the actual source pair \((m,n)=(13,2)\)
gives
\[
 (a,y,z)=(4,130,20),\qquad A=20,\qquad
 (A+z,A-z')=(40,10),\qquad (b,y,z')=(5,130,10).
\]
The raw target pair is
\((U,V)=(4\cdot7\cdot13+13\cdot2,\,5\cdot3\cdot2)
=(390,30)\); its \(g=30\) gives \((m',n')=(13,1)\).
The original identities can also be checked over denominators
\(260\) and \(130\): their reciprocal numerators are respectively
\(65+2+13=80\) and \(26+1+13=40\).

At \(p=37,a=12,R=11\), the actual full source witness
\((m,n)=(3,74)\) gives \((a,y,z)=(12,42,1036)\).
Here
\[
 A=156,\quad A+z=1192,\quad A^2=24336
 =20\cdot1192+496.
\]
Thus the exact positive rational successor has
\[
 z'=\frac{20202}{149},\quad
 (U,V)=(3278,10582),\quad g=22,\quad(m',n')=(149,481).
\]
Its target supply is \(p(a+1)=481=37\cdot13\), and the retained
numerator \(149\) does not divide it. Indeed \(481=3\cdot149+34\).
Even though \(15\mid149+481=630\), the actual divisor budget fails.
This computes the failure of this specific successor on an occupied
source shell, while the successful \(p=13\) case proves that its
integer-return domain is nonempty.


\section{The Navier--Stokes auxiliary cover and exact arithmetic sampling}

The independently proved calculations below use the literal auxiliary
torus matrix in OpenAI's public manuscript
\href{https://cdn.openai.com/pdf/32d9f210-8b73-45e0-91bc-82a30aef8a9a/navier-stokes.pdf}
{\emph{Finite Time Blowup for Navier--Stokes}}, Section~6.
Two different files were served at that URL on 8 September 2026;
both versions are identified by their complete digests in
\texttt{ns\_operator\_bridge/SOURCES.json}. The first has 165 pages
and the later one 166. The matrix (6.2), its derivative coordinates
(6.3)--(6.7), and Lemma~6.2 are the specific source passages used
here. The proofs below establish the stated operator identities
directly; they do not import the manuscript's full blowup claim as a
proved result of this repository. The pinned formal source is
\href{https://github.com/openai/NavierStokesAndEuler/tree/8937a8f4cbc7abaab5e9e97d1cc7f5d2319d9538}
{revision 8937a8f4cbc7abaab5e9e97d1cc7f5d2319d9538}; no formal
verification process was run for this contribution.

The related work has distinct authorship and provenance.
\href{https://cims.nyu.edu/~tristanb/statement.pdf}{Tristan Buckmaster's statement}
credits joint forced-fluid work with Levent Alp\"oge and the
foundational forced-blowup programme of Diego C\'ordoba and
Luis Mart\'inez-Zoroa. The
\href{https://cims.nyu.edu/~tristanb/ipm.pdf}{IPM manuscript}
lists Alp\"oge, Buckmaster and Matei P. Coiculescu; the
\href{https://cims.nyu.edu/~tristanb/boussinesq.pdf}{Boussinesq manuscript}
lists Alp\"oge and Buckmaster and credits the earlier Boussinesq
work of C\'ordoba, Andr\'es La\'in-Sanclemente and Mart\'inez-Zoroa.
The later OpenAI PDF also discusses the hypodissipative work of
C\'ordoba, Mart\'inez-Zoroa and Fan Zheng. These are source
attributions, not additional PDE theorems asserted here. Exact
authorship evidence, passage locations, reference-number differences,
and the separately stated formal-review statuses are retained in
\texttt{ns\_operator\_bridge/ATTRIBUTION.md}.


\paragraph{Source operators and intellectual context.}
The specific auxiliary matrix and original directions used below are from
\cite[Section 6.1, equations (6.1)--(6.7), and Lemma 6.2]{OpenAINS};
the separately pinned source is \cite{OpenAISource}. Earlier forced-fluid
research is credited in \cite{BuckmasterStatement,CMZEuler,CMZIPM,CMZHypo}.
The separate smooth-forcing IPM and Boussinesq manuscripts are
\cite{ABCIPM,ABBoussinesq}; their authorship and theorem domains are not
interchanged. None of these complete fluid theorems is assumed in the
arithmetic operator proofs below.
\subsection{The auxiliary torus cover on exact finite arithmetic fibres}
\label{sec:ns-torus-finite}

The source of the operator in this subsection is the public manuscript
\emph{Finite Time Blowup for Navier--Stokes}, whose displayed author is
OpenAI, Section~6, equation~(6.2) on printed page~63 and Lemma~6.2,
equation~(6.19), on printed page~67.  The preserved source is the PDF
with SHA-256
\texttt{8c8a94ad9ac824c8b605b9827cadf7beaca}\allowbreak%
\texttt{48bd10b380de3cfc872a2c37afa81};
the associated source pin is
\texttt{8937a8f4cbc7abaab5e9e97d1cc7f5d2319d9538}.
The following proofs concern this explicitly specified integer operator
and its arithmetic counting consequences. They do not use a
Navier--Stokes existence or blowup theorem. Authorship of other results
is not inferred from this operator calculation.

Retain, in the coordinates and order of the source,
\[
 \begin{gathered}
 J=\begin{pmatrix}3&1\\1&5\end{pmatrix},\qquad \det J=14,\\
 v_r=(1,1-\sqrt2),\qquad v_t=(\sqrt2-1,1),\\
 Jv_r=(4-\sqrt2)v_r,\qquad Jv_t=(4+\sqrt2)v_t.
 \end{gathered}
\]
For example the two components of $Jv_r$ are $4-\sqrt2$ and
$6-5\sqrt2=(4-\sqrt2)(1-\sqrt2)$; those of $Jv_t$ are
$3\sqrt2-2=(4+\sqrt2)(\sqrt2-1)$ and $4+\sqrt2$.
No eigenvector is divided by its length. Put
$\mathbb T=\mathbb R/\mathbb Z$ and
$\Phi:\mathbb T^2\to\mathbb T^2$, $\Phi(x,y)=(3x+y,x+5y)$.

\begin{theorem}[Directional operators in the unchanged source coordinates]
\label{thm:ns-directional-intertwining}
Let $\mathcal P$ be the complex trigonometric polynomials on
$\mathbb T^2$, and let $\mathcal P_0$ consist of those with zero
constant Fourier coefficient. For $m\geq0$ define
$P_m:\mathcal P\to\mathcal P$ by $(P_m f)(Y)=f(J^mY)$.
Writing $f(Y)=\sum_{n\in\mathbb Z^2}\widehat f(n)
e^{2\pi i n\cdot Y}$ with finite support gives the exact formula
\[
 P_m f(Y)=\sum_n\widehat f(n)e^{2\pi i((J^m)^Tn)\cdot Y}.
                                                               \tag{T0a}
\]
Thus $P_m$ is injective; its image is precisely the polynomials
whose frequencies belong to $(J^m)^T\mathbb Z^2$, and its inverse
on that image has coefficient at $n$ equal to the coefficient at
$(J^m)^Tn$ of the image polynomial. It preserves the constant
coefficient and the probability Haar average.

For $(v,\lambda)=(v_r,4-\sqrt2)$ or $(v_t,4+\sqrt2)$, define
$D_v=v\cdot\partial_Y$. Then
\[
 D_vP_m=\lambda^mP_mD_v,\qquad
 D_v^{-1}P_m=\lambda^{-m}P_mD_v^{-1}\quad\text{on }\mathcal P_0,
                                                               \tag{T0b}
\]
where the exact inverse with its original Fourier constant is
\[
 D_v^{-1}f(Y)=\sum_{n\ne0}
       \frac{\widehat f(n)}{2\pi i(v\cdot n)}e^{2\pi i n\cdot Y}.
                                                               \tag{T0c}
\]
The kernel of $D_v:\mathcal P\to\mathcal P$ is exactly the
constant polynomials; its image is $\mathcal P_0$; and every
fibre at $f\in\mathcal P_0$ is $D_v^{-1}f+\mathbb C$.
\end{theorem}
\begin{proof}
The equality $n\cdot J^mY=((J^m)^Tn)\cdot Y$ proves (T0a).
The integer frequency map is injective because $\det J^m=14^m$
is nonzero. This proves the image, inverse and singleton fibres
of $P_m$, and the zero frequency is preserved in both directions.
Integrating each character over $[0,1]^2$ proves the averaging
claim directly. By the chain rule
$v\cdot\partial_Y f(J^mY)=(J^mv)\cdot(\partial f)(J^mY)$;
the displayed eigenvector computations give $J^mv=\lambda^mv$.
This proves the first identity (T0b) with the actual source vectors.
For nonzero $n=(n_1,n_2)\in\mathbb Z^2$, neither
$n_1+(1-\sqrt2)n_2$ nor $(\sqrt2-1)n_1+n_2$ is zero.
Otherwise, unless both integer coordinates were zero, solving
that equation would make $\sqrt2$ rational. To recall the exact
obstruction, a reduced fraction $a/b$ with square two satisfies
$a^2=2b^2$, forcing $a$ and then $b$ even, a contradiction.
Each nonconstant Fourier mode of $D_v$ therefore has nonzero
multiplier $2\pi i(v\cdot n)$. This proves (T0c), the kernel,
the image and every asserted fibre. On a pulled-back mode the
denominator is
$2\pi i(v\cdot(J^m)^Tn)=\lambda^m2\pi i(v\cdot n)$.
Its reciprocal proves the second identity (T0b).
\end{proof}

\begin{theorem}[The cover and its complete finite-torsion restriction]
\label{thm:ns-torus-fibres}
For a target $(a,b)\in\mathbb T^2$, choose real representatives
$\widetilde a,\widetilde b$. Its complete fibre under $\Phi$ is
\[
 \left\{\left(
 \frac{5\widetilde a-\widetilde b+k}{14},
 \widetilde a-3\frac{5\widetilde a-\widetilde b+k}{14}
 \right)\pmod{\mathbb Z^2}:k=0,\ldots,13\right\}.       \tag{T1}
\]
In particular $\Phi$ is surjective and
$\ker\Phi=\{(k/14,-3k/14):k=0,\ldots,13\}$.

For an integer $N\geq1$, identify $\mathbb T[N]$ with
$\mathbb Z/N\mathbb Z$ by $j\mapsto j/N+\mathbb Z$, with inverse
$z\mapsto N\widetilde z\pmod N$ for any real representative
$\widetilde z$ of $z$. Changing that representative by an integer
changes $N\widetilde z$ by a multiple of $N$.
The restriction $J_N$ to
$(\mathbb Z/N\mathbb Z)^2$ has image and kernel
\[
 \begin{split}
 I_N&=\{(a,b):g\mid5a-b\},\qquad g=\gcd(N,14),\\
 K_N&=\{(kN/g,-3kN/g)\pmod N:0\leq k<g\}.
 \end{split}                                                    \tag{T2}
\]
Divisibility in (T2) is independent of the chosen integer target
representatives. A target outside $I_N$ has empty fibre. For a target
inside $I_N$, put $n=N/g$. If $n=1$ put $x_0=0$; otherwise put
\[
 x_0\equiv (14/g)^{-1}(5a-b)/g\pmod n.
\]
Its entire fibre is
\[
 \{(x_0+kn,a-3x_0-3kn)\pmod N:0\leq k<g\}.           \tag{T3}
\]
Consequently $|K_N|=g$ and $|I_N|=N^2/g$.
\end{theorem}
\begin{proof}
The first equation $3x+y=a$ forces $y=a-3x$ in either group.
Substitution in the second equation gives $14x=5a-b$.
On $\mathbb T$ this equation has exactly the fourteen solutions in
(T1): after choosing real representatives the difference from
$(5\widetilde a-\widetilde b)/14$ belongs to
$\frac1{14}\mathbb Z/\mathbb Z$. Their first coordinates are distinct.
Changing representatives leaves the solution equations unchanged, so
leaves the displayed full fibre unchanged. On $\mathbb Z/N\mathbb Z$
the equation is solvable precisely when $g\mid5a-b$: necessity follows
because $g$ divides both $14$ and $N$; sufficiency follows by dividing
the integer congruence by $g$, since $\gcd(14/g,n)=1$.
The resulting single class modulo $n$ has exactly the $g$ lifts in
(T3). For $n=1$ that divided congruence has its single zero class.
At the zero target this gives (T2)'s kernel. Every occupied fibre has
$g$ elements and the domain has $N^2$ elements, proving the image
cardinality. The coordinate identification with torsion commutes with
$J$ because every matrix entry is an integer.
\end{proof}

\begin{theorem}[Exact averaging, aliasing, and the complete-preimage repair]
\label{thm:ns-torus-average}
For every function $f:(\mathbb Z/N\mathbb Z)^2\to\mathbb C$,
\[
 \frac1{N^2}\sum_{x,y\bmod N} f(3x+y,x+5y)
   =\frac g{N^2}\sum_{(a,b)\in I_N}f(a,b).             \tag{T4}
\]
This equals the uniform mean of $f$ on all targets for every $f$ if
and only if $g=1$. For a frequency $(u,v)\bmod N$, the character
$\exp(2\pi i(ua+vb)/N)$ pulls back to frequency
$(3u+v,u+5v)$. Its mean after pullback is one precisely at
\[
 (u,v)=(kN/g,-3kN/g)\pmod N,\qquad 0\leq k<g,          \tag{T5}
\]
and is zero at every other frequency.

For any finite subgroup $D\subset\mathbb T^2$, let
$H=\Phi^{-1}(D)$. Then $H$ is a finite subgroup,
$|H|=14|D|$, and $\Phi:H\to D$ has the complete fibres (T1).
Define typed linear maps
\[
 \begin{aligned}
 P:\mathbb C^D&\longrightarrow\mathbb C^H,& Pf(h)&=f(\Phi h),\\
 L:\mathbb C^H&\longrightarrow\mathbb C^D,& Lv(d)&=
       \frac1{14}\sum_{\Phi h=d}v(h).
 \end{aligned}                                                   \tag{T6}
\]
Then $LP$ is the identity, $PL$ is the projection onto functions
constant on every deck orbit $h+\ker\Phi$, and
\[
 \frac1{|H|}\sum_{h\in H}Pf(h)=\frac1{|D|}\sum_{d\in D}f(d),
 \qquad
 \langle Pf,v\rangle_H=\langle f,Lv\rangle_D.          \tag{T7}
\]
Here $\langle f,v\rangle_D=|D|^{-1}\sum f\overline v$,
and likewise on $H$. The full fibre of $L$ at $f$ is
$Pf+\ker L$, where
$\ker L=\{v:\sum_{\Phi h=d}v(h)=0\text{ for every }d\}$.
The map $P$ is injective, and its fibre at $v$ is the singleton
$\{Lv\}$ if $v$ is deck-invariant, and is empty otherwise.
\end{theorem}
\begin{proof}
Sum $f$ by the fibres (T3) to obtain (T4). If $g=1$ the image is the
whole group. If $g>1$, the function indicating the zero target has
pullback mean $g/N^2$ and full target mean $1/N^2$, proving failure
without any regularity assumption. For characters, the geometric
sum $\sum_{j=0}^{N-1}\exp(2\pi i kj/N)$ is $N$ if $N\mid k$
and zero otherwise: multiplication by its ratio leaves the sum
unchanged, so a ratio different from one forces it to vanish. Apply this to
both input coordinates and use (T2), since $J$ is symmetric.

The inverse image $H$ is a subgroup because $\Phi$ is a homomorphism.
It is the disjoint union of the fourteen-element fibres over $D$.
On every such fibre $Pf$ is the constant $f(d)$, which proves
$LP=1$ and the first identity (T7). Formula (T6) makes $PLv$ the
arithmetic mean of $v$ on its deck orbit. It therefore has image
exactly the deck-invariant functions and squares to itself.
Summing $f(d)\overline{v(h)}$ on these same fibres proves the
inner-product identity with the stated factors $|H|=14|D|$.
The kernel and all claimed fibres follow directly: $L(v-Pf)=0$
is equivalent to $Lv=f$, while $Pf=v$ implies and, for
deck-invariant $v$, is implied by $f=Lv$.
\end{proof}

For the specific subgroup $D=\mathbb T[N]^2$, (T1) gives points
of $H$ with denominators dividing $14N$, and $|H|=14N^2$.
Thus (T7) changes the sampling domain to this explicitly given
preimage; it does not assert that the original $N^2$ samples acquire
new values. The continuous identity in source equation~(6.19)
has the same Fourier calculation: on $\mathbb T^2$ a frequency
$n\in\mathbb Z^2$ has nonzero mean exactly when $n=0$, and
$J^Tn=0$ exactly when $n=0$ because $\det J=14$.
For completeness this identity holds for all continuous functions.
Set $K_M(t)=M^{-1}|\sum_{j=0}^{M-1}e^{2\pi ijt}|^2$.
Expansion shows that this is a nonnegative trigonometric polynomial;
integrating its expansion shows its mass is one, since exactly the
$M$ diagonal terms survive. Convolution with
$K_M(t_1)K_M(t_2)$ is a trigonometric polynomial and converges
uniformly to each continuous function. Indeed, outside any fixed
neighborhood of zero each one-dimensional kernel is bounded by
$1/(M\sin^2(\pi\delta))$; uniform continuity bounds the integral
inside that neighborhood and this bound controls the complement.
The finite Fourier identity therefore passes to the integral limit.
This argument uses the original torus and its probability Haar
measure, not a replacement of a finite exponent box by a subgroup.

\begin{proposition}[Powers and product covers, with all lift labels]
\label{prop:ns-torus-powers}
For $m\geq0$ the map $\Phi_m=J^m$ has torus fibres of size $14^m$.
For a target $d$ with real representative $\widetilde d$, the exact
parametrization is
\[
 \mathbb Z^2/J^m\mathbb Z^2\longrightarrow\Phi_m^{-1}(d),
 \qquad [k]\longmapsto J^{-m}(\widetilde d+k)\pmod{\mathbb Z^2}.
                                                               \tag{T8}
\]
Its inverse sends a real lift $\widetilde h$ to
$[J^m\widetilde h-\widetilde d]$. In particular a full preimage of
a finite subgroup $D$ has $14^m|D|$ elements, and (T6)--(T7)
hold with $\Phi_m$ and $14^m$. For a block diagonal product of $s$
copies of $J^m$, use $s$ independent copies of (T8); every fibre
has $14^{ms}$ elements and the averaging statements have that factor.

On fixed $N$-torsion, define $S_0(d)=\{d\}$ and
$S_{j+1}(d)=\coprod_{w\in S_j(d)}J_N^{-1}(w)$, with every
one-step fibre given by (T3). Then $S_j(d)=\{h:J_N^jh=d\}$
with no repeated elements. These exact recursion fibres must be
used in place of an assumed size $\gcd(N,14)^j$.
\end{proposition}
\begin{proof}
Changing $k$ by $J^m z$ changes (T8) by $z$ and hence changes no
torus point. Conversely equality of torus points implies this exact
change in $k$. Every lift in the fibre gives an integral difference
$J^m\widetilde h-\widetilde d$, proving surjectivity and the inverse.
One-step surjectivity and fourteen-element fibres imply by induction
that the fibre of $\Phi_m$ consists of $14^m$ points: the preimages
of distinct intermediate points are disjoint. This proves the
cardinality of the quotient as well. Fibre counting now proves all
averaging assertions exactly as in (T6)--(T7), and independent
coordinate choices prove the product statement. In the torsion
recursion each candidate $h$ has the unique intermediate point
$J_Nh$, so the displayed union is disjoint and the induction gives
precisely $J_N^{j+1}h=d$. For example, at $N=2$, $J_N$ is the
matrix whose four entries are one; its square is zero. Its kernel
has two elements at $j=1$ and four elements at every $j\geq2$,
which also verifies why the proposed exponential size is incorrect.
\end{proof}

\begin{lemma}[The actual finite dual and its coordinate embedding]
\label{lem:ns-finite-dual}
Let $G$ be a finite abelian multiplicative group and
$\widehat G=\operatorname{Hom}(G,\{z\in\mathbb C:|z|=1\})$.
Then $|\widehat G|=|G|$, its characters separate the elements of
$G$, and
\[
 \frac1{|G|}\sum_{\chi\in\widehat G}\chi(u)
      =\begin{cases}1&u=1,\\0&u\ne1.\end{cases}       \tag{T9}
\]
Choose and retain an ordered generating tuple $g_1,\ldots,g_s$ of
$G$ (the list of all elements is permitted). The map
\[
 E:\widehat G\to\mathbb T^s,\qquad
 E(\chi)_j=\theta_j\pmod{\mathbb Z},\quad
       e^{2\pi i\theta_j}=\chi(g_j)                  \tag{T10}
\]
is an injective homomorphism. Its inverse on the image sends
$\theta$ to the character
$g_1^{n_1}\cdots g_s^{n_s}\mapsto
\exp(2\pi i\sum n_j\theta_j)$. The image consists exactly of
the vectors for which $\sum n_j\theta_j=0$ in $\mathbb T$
for every relation $\prod g_j^{n_j}=1$.
\end{lemma}
\begin{proof}
We give the character construction to make the dual and its fibres
explicit. Suppose a character $\chi$ is already defined on a
subgroup $B$, and choose $g\notin B$. Let $r$ be the smallest
positive integer such that $g^r\in B$. Every element of
$\langle B,g\rangle$ has a unique presentation $bg^j$ with
$b\in B$ and $0\leq j<r$: equality of two such presentations
would place $g^{j-j'}$ in $B$. For each of the $r$ distinct complex
roots $z$ of $z^r=\chi(g^r)$, define
$\widetilde\chi(bg^j)=\chi(b)z^j$.
The relation at $j=r$ proves multiplicativity, and restriction and
the value at $g$ recover $\chi,z$ uniquely. Starting at the trivial
subgroup and adjoining elements until reaching $G$ constructs
exactly $\prod r=|G|$ characters, because the subgroup order also
multiplies by $r$ at each step. To separate $u\ne1$ of order $d$,
first define the character on $\langle u\rangle$ by
$u\mapsto e^{2\pi i/d}$, then extend by this construction.
Multiplication by a character nontrivial at $u$ permutes
$\widehat G$ and multiplies the sum in (T9) by a number other
than one; hence that sum is zero. At $u=1$ every summand is one.
Finally a homomorphism is determined by its values on the retained
generators. The relation equations in (T10) are exactly what makes
the inverse formula independent of presentation; multiplication
then proves it is a character. This proves every image, inverse,
and singleton-fibre assertion.
\end{proof}

On $\widehat G^2$ the source matrix is the exact endomorphism
\[
 J_G(\chi,\psi)=(\chi^3\psi,\chi\psi^5).             \tag{T11}
\]
Under the coordinate embedding
$E_2(\chi,\psi)=((E\chi)_1,(E\psi)_1,\ldots,
(E\chi)_s,(E\psi)_s)$ it is the restriction of the $s$-block
torus map $\Phi^{\oplus s}$. Thus $E_2J_G=\Phi^{\oplus s}E_2$
coordinate by coordinate, with an explicitly proved inverse for
$E_2$ on its image. This is the precise morphism between the
continuous cover and the finite character group.

\begin{theorem}[Exact arithmetic aliasing and the original packet fibres]
\label{thm:ns-es-packet}
For a function $f$ on $\widehat G^2$, define
\[
 \widehat f(u,v)=\frac1{|G|^2}
   \sum_{\alpha,\beta\in\widehat G}
       f(\alpha,\beta)\overline{\alpha(u)\beta(v)}.
\]
Then its mean after the finite restriction (T11) is exactly
\[
 \frac1{|G|^2}\sum_{\chi,\psi}f(J_G(\chi,\psi))
   =\sum_{\substack{u\in G\\u^{14}=1}}
           \widehat f(u,u^{-3}).                     \tag{T12}
\]
Each one-coordinate marginal of $J_G$ is uniform, although their
joint distribution need not be uniform. The complete kernel is
$\{(\chi,\chi^{-3}):\chi^{14}=1\}$; a target $(\alpha,\beta)$
is in the image precisely when $\alpha^5\beta^{-1}$ is a
fourteenth power in $\widehat G$. Its complete fibre consists of
$(\chi,\alpha\chi^{-3})$ for all fourteenth roots
$\chi^{14}=\alpha^5\beta^{-1}$.

Retain an original shell $p=12h+1$ prime, $h\geq1$,
$3h+1\leq a\leq9h$, and $R=4a-p$. Put $G=U(R)$ and retain
the entire prime exponent box
\[
 {\cal B}_a=\prod_{\ell\mid a}\{0,\ldots,2v_\ell(a)\},\quad
 U_a(e)=\prod_{\ell\mid a}\ell^{e_\ell},\quad
 c_a(g)=\#\{e\in{\cal B}_a:[U_a(e)]_R=g\}.
\]
Here $R\geq3$ is odd, $a<p$, and $\gcd(a,R)=1$, since any
common prime divisor divides $p=4a-R$ and cannot divide $a<p$.
In particular every listed prime factor is a unit. Define the
actual packet function at target $g\in G$ by
\[
 F_{a,g}(\chi)=\chi(g^{-1})
       \prod_{\ell\mid a}\sum_{e=0}^{2v_\ell(a)}\chi([\ell]_R)^e.
                                                               \tag{T13}
\]
The exact one-step count is
\[
 c_a(g)=\frac1{|G|}\sum_\chi F_{a,g}(\chi)
       =\frac1{|G|^2}\sum_{\chi,\psi}F_{a,g}(\chi^3\psi).
                                                               \tag{T14}
\]
Every original exponent is retained in (T13), and the counted
fibre maps bijectively to the original divisors of $a^2$ congruent
to $g$, with inverse $u\mapsto(v_\ell(u))_{\ell\mid a}$.

For two such finite packets on this same group, with coefficients
$c_a,c_b$ and targets $g,h$, where all factors of $b$ are also
units modulo this same $R$, the exact joint mean is
\[
 \frac1{|G|^2}\sum_{\chi,\psi}
    F_{a,g}(\chi^3\psi)F_{b,h}(\chi\psi^5)
   =\sum_{\substack{u\in G\\u^{14}=1}}
          c_a(gu)c_b(hu^{-3}).                       \tag{T15}
\]
All summands are nonnegative integers. Their complete pair fibre
consists of the original exponent pairs $(e,f)$ with
$[U_a(e)]_R=gu$, $[U_b(f)]_R=hu^{-3}$, indexed by $u^{14}=1$.
The index is uniquely recovered as $u=g^{-1}[U_a(e)]_R$.

More generally write the first row of the original integer power
$J^m$ as $(A_m,B_m)$ and set $d_m=\gcd(A_m,B_m)$, with
$J^0$ having row $(1,0)$ and $d_0=1$. Then
\[
 \frac1{|G|^2}\sum_{\chi,\psi}
    F_{a,g}(\chi^{A_m}\psi^{B_m})
       =\sum_{u^{d_m}=1}c_a(gu).                     \tag{T16}
\]
The powers and their integer row entries are retained, so the
additional torsion terms are visible at every level.
\end{theorem}
\begin{proof}
Formula (T9) makes the functions $(\alpha,\beta)\mapsto
\alpha(u)\beta(v)$ an orthonormal family of $|G|^2$ elements
in the $|G|^2$-dimensional space of all functions on
$\widehat G^2$. They are therefore a basis and have the displayed
Fourier coefficients. The pullback of this basis element is
\[
 (\chi^3\psi)(u)(\chi\psi^5)(v)
       =\chi(u^3v)\psi(uv^5).
\]
By (T9) its mean is one exactly when $u^3v=uv^5=1$.
The first equation forces $v=u^{-3}$; substituting in the second
gives $u^{-14}=1$. This proves (T12). For fixed $\chi$ the first
output $\chi^3\psi$ runs bijectively through $\widehat G$ as
$\psi$ varies, so its marginal is uniform. For fixed $\psi$ the
second output $\chi\psi^5$ likewise runs bijectively as $\chi$
varies. Solving the first target equation gives
$\psi=\alpha\chi^{-3}$; substituting in the second yields
$\beta=\alpha^5\chi^{-14}$, proving the image and full fibres.
The zero-target fibre yields the stated kernel.

Expanding the finite product (T13) gives precisely
$\sum_{e\in{\cal B}_a}\chi(g^{-1}[U_a(e)]_R)$.
Applying (T9) proves its mean is $c_a(g)$; applying the first
marginal calculation proves the second equality in (T14).
Unique prime factorization proves the divisor bijection and the
displayed inverse on every fibre. In the product of the two
expanded packets, (T9) keeps exactly the exponent pairs whose
quotients $u=g^{-1}[U_a(e)]_R$ and $v=h^{-1}[U_b(f)]_R$
satisfy $u^3v=uv^5=1$. The preceding calculation gives precisely
the indexed fibres in (T15). No summand or exponent allocation is
discarded. Finally the mean of
$\chi(u)^{A_m}\psi(u)^{B_m}$ is one precisely when
$u^{A_m}=u^{B_m}=1$. These equations are equivalent to
$u^{d_m}=1$: Bezout's integer relation yields one implication,
and divisibility of both row entries by $d_m$ yields the other.
Expansion of the single packet now proves (T16).
\end{proof}

\begin{corollary}[Complete-preimage coefficient extraction and an original-shell discrepancy]
\label{cor:ns-es-cover-repair}
Let $D=E_2(\widehat G^2)\subset\mathbb T^{2s}$ and let
$H_m=\{z\in\mathbb T^{2s}:(J^m)^{\oplus s}z\in D\}$ be
the full inverse image under the block map $(J^m)^{\oplus s}$. For any function $f$ on
$\widehat G^2$, the exact identity is
\[
 \frac1{|H_m|}\sum_{z\in H_m}
   f\!\left(E_2^{-1}((J^m)^{\oplus s}z)\right)
   =\frac1{|G|^2}\sum_{\chi,\psi}f(\chi,\psi),
 \qquad |H_m|=14^{ms}|G|^2.                           \tag{T17}
\]
All lifts are given by the product of (T8), and $E_2^{-1}$ is
the coordinate formula (T10); thus every map in (T17) is typed
on its full domain. For $f(\chi,\psi)=F_{a,g}(\chi)$,
(T17) equals $c_a(g)$; for
$f(\chi,\psi)=F_{a,g}(\chi)F_{b,h}(\psi)$ it equals
$c_a(g)c_b(h)$, with no extra terms.

In the original shell $p=13$, $a=4$, $R=3$, the full box is
$e_2=0,1,2,3,4$. The residues of its five divisors
$1,2,4,8,16$ are $1,2,1,2,1$. Hence $c_a(1)=3$ and
$c_a(2)=2$, and both original targets $-4^{-1}$ and $-a$ are
$2$ in $U(3)$. Formula (T15) with both packets and both targets
equal to these values is $2\cdot2+3\cdot3=13$, whereas
the independent target count is $2\cdot2=4$.
Moreover
$J^2=\left(\begin{smallmatrix}10&8\\8&26\end{smallmatrix}\right)$,
so (T16) at $m=2$ is $c_a(2)+c_a(1)=5$, whereas the
original coefficient remains $2$. Formula (T17) returns $2$
and $4$, respectively, at every cover level.
\end{corollary}
\begin{proof}
Apply the product-cover assertion of Proposition~\ref{prop:ns-torus-powers}
to the subgroup $D$. Lemma~\ref{lem:ns-finite-dual} proves
$|D|=|G|^2$ and the precise inverse used in (T17). Every target
therefore occurs $14^{ms}$ times, which proves the identity by
fibre counting. The one-packet conclusion is (T14); the two-packet
conclusion follows by separating the independent finite sums.
The displayed residues directly give the two coefficients in
the example. Both elements of $U(3)$ satisfy $u^{14}=1$;
in (T15) their two contributions are respectively $4$ and $9$.
The first row of $J^2$ has gcd $2$, and both group elements
satisfy $u^2=1$, proving its stated mean. All original five
exponents and both nonzero coefficient fibres have been listed.
\end{proof}

The bridge therefore supplies an exact finite sampling rule for the
original shell packet and an explicit obstruction to transferring
joint averages or repeated covers on an unchanged finite torsion set.
It proves neither global shell occupancy nor a PDE assertion.

\subsection{Exact character correction for every original sieve fibre}
\label{sec:ns-character-sieve}

The torus operator in the preceding subsection acts on the arithmetic
characters of each original shell. We now calculate this action and all
of its fibres, including those needed when its degree shares factors with
the finite character group. No Navier--Stokes existence assertion is a
hypothesis of these calculations.

Fix the original prime $p=12h+1$, $h\geq1$, and an original shell
$3h+1\leq a\leq9h$. Retain $R=4a-p$, $S=pa$, the distinct increasing
prime factors $\ell_1,\ldots,\ell_s$ of $a$, and their original positive
valuations $v_1,\ldots,v_s$. Put
\[
 G=(\mathbb Z/R\mathbb Z)^\times,\qquad
 \widehat G=\operatorname{Hom}(G,\mathbb R/\mathbb Z),\qquad
 \langle\chi,g\rangle=\exp(2\pi i\chi(g)).
\]
We write $G$ multiplicatively and $\widehat G$ additively. Thus
$\chi(gg')=\chi(g)+\chi(g')$ in $\mathbb R/\mathbb Z$, with the
displayed $2\pi i$ retained in every complex character. Let $n=|G|$.
The exact change from the preceding subsection's multiplicative
dual is the isomorphism
$\chi\mapsto(g\mapsto\exp(2\pi i\chi(g)))$; its inverse takes the
unique class in $\mathbb R/\mathbb Z$ with that exponential at each
$g$. The exponential has kernel $\mathbb Z$ on $\mathbb R$ and
maps onto the complex unit circle. Consequently both inverse
compositions are identities and the character equations are preserved.
Thus the two dual conventions are explicitly related, with singleton
fibres; neither is silently substituted for the other.
Every element $g\in G$ has order dividing $n$: its cyclic subgroup
partitions $G$ into cosets, each with that order. Consequently
$\chi(g)\in n^{-1}\mathbb Z/\mathbb Z$. The complete set $\widehat G$
is therefore finite and can be obtained by enumerating all functions
$G\to n^{-1}\mathbb Z/\mathbb Z$ and retaining exactly the identity
and product equations. This is a finite construction on the actual
unit multiplication table; it imposes no cyclicity hypothesis on $G$.

For completeness, these characters separate $G$. If $g$ has order
$o>1$, define $g^j\mapsto j/o$ on its cyclic subgroup. To extend an
additive character from a subgroup $H$ of any finite abelian group to
$H+\mathbb Zx$, choose the least $m>0$ with $mx\in H$, choose any
$\beta\in\mathbb R/\mathbb Z$ with $m\beta=\chi(mx)$, and set
$\chi'(h+jx)=\chi(h)+j\beta$. If $h+jx=h'+j'x$, then
$j-j'=qm$ and $h-h'=-qmx$, so the two values agree. The character
equations follow by addition. Such a $\beta$ exists: if $r\in\mathbb R$
represents $\chi(mx)$, then $r/m$ is a choice. Adjoining elements
until the finite group is exhausted extends the prescribed nonzero
value $1/o$. This proves separation without an unproved extension
assumption.

\begin{lemma}[The original labelled packet as an exact character projection]
\label{lem:ns-labelled-character-packet}
Let $\mathbf X=(X_1,\ldots,X_s)$ be independent formal variables. For
$t\in G$ define
\[
 \begin{split}
 P_a(\chi;\mathbf X)
   &=\prod_{j=1}^s\left(\sum_{e_j=0}^{2v_j}
           X_j^{e_j}\langle\chi,[\ell_j]\rangle^{e_j}\right),\\
 C_{a,t}(\mathbf X)
   &=\frac1{|\widehat G|}\sum_{\chi\in\widehat G}
          \langle\chi,t\rangle^{-1}P_a(\chi;\mathbf X).
 \end{split}
\]
Then the complete coefficient identity is
\[
 C_{a,t}(\mathbf X)=
 \sum_{\substack{0\leq e_j\leq2v_j\ (1\leq j\leq s)\\
                   \prod_j\ell_j^{e_j}\equiv t\pmod R}}
       \prod_{j=1}^s X_j^{e_j}.
\]
In particular every retained coefficient is $1$, and the exponent of
$X_j$ is the original valuation of the corresponding positive divisor
$u=\prod_j\ell_j^{e_j}$. The cardinalities of the original exterior
and middle fibres are respectively
\[
 C_{a,-4^{-1}}(\mathbf1)=|E_a|,\qquad
 C_{a,-a}(\mathbf1)=|M_a|.
\]
\end{lemma}
\begin{proof}
For $g=1$, every character value is $1$. For $g\ne1$, choose
$\eta$ with $\langle\eta,g\rangle\ne1$ by separation. Translation
$\chi\mapsto\chi+\eta$ permutes the finite set $\widehat G$; hence
its character sum equals itself multiplied by
$\langle\eta,g\rangle$ and is zero. Expand the finite product in
$P_a$ without changing any exponent interval. For the term with
the specified original vector $e$, character averaging is therefore
$1$ exactly when $(\prod_j[\ell_j]^{e_j})t^{-1}=1$, and is otherwise
$0$. Distinct original vectors have distinct formal monomials and
distinct positive integers by unique prime factorization. This proves
the identity and both inverse coordinate assertions. Since $R$ is
odd and $\gcd(a,R)=1$, both targets are units. Their respective
equations are exactly $R\mid4u+1$ and $R\mid u+a$.
\end{proof}

The exact comparison with the source's two-dimensional auxiliary torus
is, for each retained unit $g\in G$, the homomorphism
\[
 E_g:\widehat G^2\longrightarrow(\mathbb R/\mathbb Z)^2,
 \qquad E_g(\alpha,\beta)=(\alpha(g),\beta(g)).
\]
The family $(E_g)_{g\in G}$ is injective: equality on all $g$ is
equality of the two functions. Its image consists precisely of pairs
of tables satisfying the character identity and product equations
above. Its inverse reads those two tables. Thus all these coordinates,
not a selected single evaluation, retain the full arithmetic datum.
With the literal source matrix $J_g=\left(\begin{smallmatrix}3&1\\1&5\end{smallmatrix}\right)$,
define
\[
 \mathcal J(\alpha,\beta)=(3\alpha+\beta,\alpha+5\beta).
\]
For every $g$ we have $E_g\mathcal J=J_gE_g$ in the torus, by
evaluating the two displayed sums. This proves the coordinate
morphism to the source operator on every original unit, including
each original prime residue and both targets.

\begin{theorem}[Complete finite image, fibres and correction cosets]
\label{thm:ns-character-coset-correction}
Write $14\widehat G=\{14\chi:\chi\in\widehat G\}$ and
$\widehat G[14]=\{\chi:14\chi=0\}$, and set
$\delta=|\widehat G[14]|$. Then
\[
 \begin{split}
 \ker\mathcal J&=\{(\tau,-3\tau):\tau\in\widehat G[14]\},\\
 \operatorname{im}\mathcal J
   &=\{(\gamma,\eta):\eta-5\gamma\in14\widehat G\}.
 \end{split}
\]
There is an exact sequence of finite abelian groups
\[
 0\longrightarrow\ker\mathcal J\longrightarrow\widehat G^2
   \xrightarrow{\mathcal J}\widehat G^2
   \xrightarrow{q}\widehat G/14\widehat G\longrightarrow0,
 \qquad q(\gamma,\eta)=[\eta-5\gamma].
\]
For every image point, its full inverse fibre is
\[
 \{(\alpha,\gamma-3\alpha):14\alpha=5\gamma-\eta\}.
\]
The displayed root equation is solved by enumerating the already
defined finite set $\widehat G$; all of its solutions are one
solution plus $\widehat G[14]$.

Choose a complete set $\mathcal R\subset\widehat G$ of coset
representatives modulo $14\widehat G$, retaining the representative
table. There are exactly $\delta$ representatives. Define
\[
 T:\widehat G^2\times\mathcal R\longrightarrow\widehat G^2,
 \qquad T(\alpha,\beta,\rho)
       =(3\alpha+\beta,\alpha+5\beta+\rho).
\]
This map is surjective, and every fibre has exactly $\delta$
points. For an arbitrary complex-valued function $F$ on
$\widehat G^2$, or a function with values in a complex polynomial
algebra, the exact identity is
\[
 \frac1{\delta|\widehat G|^2}
 \sum_{\rho\in\mathcal R}\sum_{\alpha,\beta\in\widehat G}
       F(T(\alpha,\beta,\rho))
 =\frac1{|\widehat G|^2}\sum_{\gamma,\eta\in\widehat G}F(\gamma,\eta).
\]
\end{theorem}
\begin{proof}
The first output equation forces $\beta=\gamma-3\alpha$.
Substituting in the second gives $14\alpha=5\gamma-\eta$.
This proves the image, the full fibre and the kernel, including
both directions. The quotient map $q$ is surjective because
$q(0,\eta)=[\eta]$; its kernel is precisely that image. Each fibre
of multiplication by $14$ on $\widehat G$ is a coset of
$\widehat G[14]$, so $|14\widehat G|=|\widehat G|/\delta$ and
$|\widehat G/14\widehat G|=\delta$.

For an arbitrary output $(\gamma,\eta)$ of $T$, choose the unique
recorded representative $\rho$ of $[\eta-5\gamma]$. Then
$14\alpha=5\gamma-\eta+\rho$ has exactly $\delta$ solutions,
and for each of them $\beta=\gamma-3\alpha$. This supplies every
preimage and proves both uniqueness of the coset coordinate and
the full fibre formula. Counting each output once for each of its
$\delta$ preimages proves the averaging identity, coefficient by
coefficient for polynomial-valued $F$.
\end{proof}

\begin{theorem}[Exact alias terms and their removal on original valuations]
\label{thm:ns-sieve-alias-removal}
For $t_1,t_2\in G$ and independent variable tuples $\mathbf X,
\mathbf Z$, put
\[
 F(\gamma,\eta)=
  \langle\gamma,t_1\rangle^{-1}P_a(\gamma;\mathbf X)
  \langle\eta,t_2\rangle^{-1}P_a(\eta;\mathbf Z),
 \qquad G[14]=\{g\in G:g^{14}=1\}.
\]
The uncorrected source matrix gives exactly
\[
 \frac1{|\widehat G|^2}\sum_{\alpha,\beta}
          F(\mathcal J(\alpha,\beta))
   =\sum_{g\in G[14]}
       C_{a,t_1g}(\mathbf X)C_{a,t_2g^{-3}}(\mathbf Z).
\]
The complete correction is
\[
 \frac1{\delta|\widehat G|^2}
     \sum_{\rho\in\mathcal R}\sum_{\alpha,\beta}
          F(T(\alpha,\beta,\rho))
       = C_{a,t_1}(\mathbf X)C_{a,t_2}(\mathbf Z).
\]
Every monomial on the right retains both original exponent vectors.
The corrected identity therefore preserves their complete joint
fibre, not only its cardinality.
\end{theorem}
\begin{proof}
For fixed original exponent vectors $e,f$, write
$u=\prod_j[\ell_j]^{e_j}$, $v=\prod_j[\ell_j]^{f_j}$,
$g=ut_1^{-1}$ and $k=vt_2^{-1}$. Its complex factor under the
uncorrected map is
\[
 \langle\alpha,g^3k\rangle\langle\beta,gk^5\rangle.
\]
The two independent character sums retain the monomial exactly when
$g^3k=1$ and $gk^5=1$. The first equality gives $k=g^{-3}$;
the second then gives $g^{-14}=1$. Conversely those conditions
give both equalities. This is precisely the first displayed sum,
including every original exponent bound and variable.

In the corrected map the same term has the additional factor
$\langle\rho,k\rangle$. For $k^{14}=1$, this value depends only
on the coset of $\rho$ modulo $14\widehat G$, since
$\langle14\chi,k\rangle=\langle\chi,k^{14}\rangle=1$.
If $k\ne1$, separation supplies a character whose value at $k$ is not $1$;
translation by its coset permutes all quotient cosets and forces
the sum over $\mathcal R$ to be zero. If $k=1$, every value is
$1$. In the retained uncorrected terms, $k=g^{-3}$ and $g^{14}=1$.
The two equations $g^3=1$ and $g^{14}=1$ imply $g=1$, since
$1=5\cdot3-14$. Thus precisely $g=k=1$ remains. These are the
original two target equations, proving the second identity with
their full monomials. It also follows directly by applying the
preceding theorem and then the independent character projections.
\end{proof}

\begin{corollary}[All finite covering histories preserve the full ES test]
\label{cor:ns-finite-history-sieve}
Fix an integer $k\geq0$, the original group, and the representative
table $\mathcal R$. For $z_0\in\widehat G^2$ and a retained
sequence $(\rho_0,\ldots,\rho_{k-1})\in\mathcal R^k$, define
\[
 z_{j+1}=\mathcal J z_j+(0,\rho_j)\quad(0\leq j<k),\qquad
 z_k=\mathcal J^k z_0+
       \sum_{j=0}^{k-1}\mathcal J^{k-1-j}(0,\rho_j).
\]
The history-to-output map has precisely $\delta^k$ preimages at
each output, and
\[
 \frac1{\delta^k|\widehat G|^2}
   \sum_{z_0\in\widehat G^2}\sum_{\rho_0,\ldots,\rho_{k-1}\in\mathcal R}
      F(z_k)
       =\frac1{|\widehat G|^2}\sum_zF(z)
\]
for every polynomial-valued $F$. In particular, taking
$F(\gamma,\eta)=\langle\gamma,t\rangle^{-1}P_a(\gamma;\mathbf X)$
recovers exactly $C_{a,t}(\mathbf X)$ for both original targets
after every finite history length.
\end{corollary}
\begin{proof}
The stated expansion of $z_k$ follows by substitution, with the
empty sum and identity map for $k=0$. For a fixed terminal state,
the preceding theorem supplies its unique last coset coordinate
and every one of its $\delta$ predecessor states. Repeating the
same inverse construction at each predecessor gives all histories
and precisely $\delta^k$ of them, with no choices identified.
The averaging formula follows by counting this exact fibre.
For the specified function, the average over $\eta$ contributes
exactly the factor $|\widehat G|$, and the other average is the
proved labelled coefficient projection.
\end{proof}

Returning to integers uses each retained monomial's unique divisor
$u=\prod_j\ell_j^{e_j}$, and then the complete ordered inverse of
Theorem~\ref{thm:shellwise-no-hit}. Its exterior orientation coordinate
and both middle divisors remain part of that return. In particular,
the corrected averages give the same $|E_a|$, $|M_a|$, the same
$q_{a,p}=|E_a|+|M_a|/2$, and the same $Q_p$ and $\mathcal O_p$
in Theorem~\ref{thm:global-no-hit}. This supplies an exact arithmetic
consequence of the source cover, including a proved obstruction to
using its uncorrected finite sampling.

The source's smooth fields use continuous auxiliary coordinates.
Here $E_g$ is the explicit finite-coordinate map to that torus, and
the preceding exact sequences and character products describe its
arithmetic restriction. The operator correspondence by itself does
not construct a Navier--Stokes velocity or force from a prime, or
establish nonemptiness of any previously empty ES fibre. No such
construction or prime-coverage conclusion is a premise of a result
above.


\paragraph{Return to the continuous source calculation.}
The next two subsections address the actual continuous operators in
\cite[Lemma 6.2, (6.17)--(6.20), (8.4)--(8.10), (8.14), and
Lemma 8.6, (8.19)--(8.23)]{OpenAINS}. The sharp constants, full inverse
fibres and strengthened derivative estimates are proved here. The original
cutoff remainder and slow derivatives remain in those identities.
\subsection{Returning the full-fibre calculus to the new proof}
\label{sec:ns-reverse-proof}

We now apply the full-fibre maps to the source's continuous auxiliary
calculus, in the direction requested by the user. The relevant source
statements are Lemma~6.2, equations (6.17)--(6.20), and Lemma~8.6,
equations (8.19)--(8.23), on printed pages 67 and 95--96 of the
preserved 166-page manuscript. We retain its unit-period torus,
$J$, $v_r$, $v_t$, $\Lambda=4-\sqrt2$, $T_g=4+\sqrt2$,
and the Fourier character $e_n(Y)=\exp(2\pi i n\cdot Y)$.
Write $\mathscr C=C^\infty(\mathbb T^2;\mathbb C)$,
$\Pi f=\int_{\mathbb T^2}f$, and $\mathscr C_0=\ker\Pi$.
The smooth inverse used below is proved, with its exact norm and
original constants, in the next subsection. The conclusions here
also hold directly on the previously defined finite polynomials.

\begin{theorem}[Continuous transfer, descent, and every operator fibre]
\label{thm:ns-continuous-transfer}
For $m\geq0$, retain $K_m=\ker(J^m:\mathbb T^2\to\mathbb T^2)$
and the full inverse parametrization (T8). Define
\[
\begin{aligned}
 P_m:\mathscr C&\longrightarrow\mathscr C,
 &P_mf(Y)&=f(J^mY),\\
 L_m:\mathscr C&\longrightarrow\mathscr C,
 &L_mg(X)&=14^{-m}\sum_{J^mY=X}g(Y),\\
 Q_mg(Y)&=14^{-m}\sum_{a\in K_m}g(Y+a).
\end{aligned}
\]
Then $L_mP_m=1$, $P_mL_m=Q_m$, and $Q_m^2=Q_m$.
The image of $P_m$ is exactly the $K_m$-invariant functions,
and $P_m^{-1}=L_m$ on that image. The full fibres of $L_m$ are
$P_mf+\ker L_m$, where its kernel consists of functions whose
sum on every full covering fibre is zero. The fibre of $P_m$ at
$g$ is $\{L_mg\}$ if $Q_mg=g$ and is empty otherwise.
Also $\ker Q_m=\ker L_m$. The complete fibre of $Q_m$ at
$h$ is $h+\ker L_m$ if $Q_mh=h$, and is empty otherwise.
For every $n\in\mathbb Z^2$,
\[
 L_me_n=
 \begin{cases}
 e_k,&n=(J^m)^Tk\text{ for }k\in\mathbb Z^2,\\
 0,&n\notin(J^m)^T\mathbb Z^2.
 \end{cases}
\]
The integral and product identities are
\[
 \Pi L_mg=\Pi g,\quad \Pi P_mf=\Pi f,\quad
 \langle P_mf,g\rangle=\langle f,L_mg\rangle,\quad
 L_m((P_mf)g)=fL_mg.
\]
All inner products use the original probability Haar measure.
In particular $\Pi((P_mf)(P_mg))=\Pi(fg)$.
For $v=v_r,v_t$ and its corresponding eigenvalue $\lambda$,
\[
 D_vL_m=\lambda^{-m}L_mD_v,
 \qquad D_v^{-1}L_m=\lambda^mL_mD_v^{-1}
 \quad\hbox{on }\mathscr C_0.
\]
The maps $Q_m$ commute with $D_v$ and its zero-mean inverse.
\end{theorem}
\begin{proof}
Every fibre is $Y+K_m$ and has exactly $14^m$ elements by
(T8). On it, $P_mf$ is the constant $f(X)$. Fibre summation
therefore proves both compositions, the projection, its image,
the stated fibres of $L_m$ and $P_m$, and the product identity.
Since $Q_m=P_mL_m$ and $P_m$ is injective, their kernels satisfy
$\ker Q_m=\ker L_m$. A value of the projection $Q_m$ is fixed
by it; if $Q_mh=h$, subtraction gives exactly
$Q_m^{-1}(h)=h+\ker Q_m$, proving its full fibre as well.
Locally the covering has $14^m$ affine inverse branches with
linear part $J^{-m}$;
their finite sum proves smoothness of $L_mg$ independently of
the choice of local branches.

For a character $e_n$, the fibre sum has factor
$14^{-m}\sum_{a\in K_m}e_n(a)$. Translation by any element
of $K_m$ leaves this sum unchanged. It is one if the character
is trivial on $K_m$, and otherwise is zero. Triviality is
equivalent to $(J^{-m})^Tn\in\mathbb Z^2$: evaluating on
$J^{-m}z\pmod{\mathbb Z^2}$ for every $z\in\mathbb Z^2$
proves both directions. In the trivial case write
$n=(J^m)^Tk$; then $e_n(Y)=e_k(X)$ on the fibre.
This proves the displayed multiplier without dropping its
nonimage frequencies.

The Fourier computation proves $\Pi L_mg=\Pi g$ first for
polynomials. The Fejer approximation already proved after
(T7) and the supremum bound $\|L_mg\|_\infty\le\|g\|_\infty$
pass this equality to continuous $g$. The same approximation
proves Haar preservation for $P_m$. Applying it to
$L_m((P_mf)\overline g)=f\overline{L_mg}$ proves the inner
product formula. The last product integral uses
$(P_mf)(P_mg)=P_m(fg)$, which holds pointwise.

On a surviving frequency $n=(J^m)^Tk$ one has
$v\cdot n=\lambda^m(v\cdot k)$. On nonsurviving frequencies
both sides of each asserted intertwining are zero. This proves
the identities on polynomials, with the precise factors
$2\pi i(v\cdot n)$ and their reciprocals. Smooth Fourier
convergence and the finite derivative loss proved below pass
them to $\mathscr C$ and $\mathscr C_0$. Finally $Q_m$ is
averaging of translations, which commute with all these
constant-direction multipliers; alternatively use its exact
frequency projection.
\end{proof}

\begin{theorem}[The mean-correction map on every common band]
\label{thm:ns-reverse-mean-correction}
Let $U$ be an open set of retained slow coordinates. For a
fixed source band scale $Q>0$, retain the source exponent
$h_{\rm NS}$, $A_{\rm NS}=1/2+h_{\rm NS}$,
$\varepsilon=Q^{h_{\rm NS}}$, and
$c_i=T_g^iQ^{1+h_{\rm NS}}$. The subscripts distinguish the
source's physical scaling constants from ES chart coordinates;
their values are unchanged. For an extended band residual
$E_i\in C^\infty(U\times\mathbb T^2)$, put
\[
 E_i^\circ=E_i-\Pi E_i,\qquad
 \mathcal C_i(E_i)=-c_i^{-1}D_{v_t}^{-1}E_i^\circ.
\]
This is a smooth, zero-mean coefficient with
$c_iD_{v_t}\mathcal C_i(E_i)=-E_i^\circ$.
Its full kernel consists of the coefficients independent of the
torus variable, and the full inverse fibre at any smooth
zero-mean $V$ is
\[
 \{\,-c_iD_{v_t}V+b:\ b\in C^\infty(U;\mathbb C)\,\}.
\]
For $i\ge j\ge0$ put $\Delta=i-j$, preserving $Q$. Then
\[
 \mathcal C_j(P_\Delta E_i)=P_\Delta\mathcal C_i(E_i).
\]
Here $P_\Delta$ acts only on the torus coordinate and
$c_j=T_g^{-\Delta}c_i$. The common representative and its
full descent test are $P_\Delta E_i$ and
$Q_\Delta E_j=E_j$, with inverse $E_i=L_\Delta E_j$.

For a finite set of band indices with $i_0=\min i$, the common
sum is $\sum_iP_{i-i_0}E_i$. Its absolute representative is
exactly $\sum_iP_iE_i=P_{i_0}\sum_iP_{i-i_0}E_i$.
On overlaps, equality of absolute representatives implies
equality after pullback to the smaller common index. The same
assertion holds for their zero-mean inverse corrections and
for sums and products of the original coefficients.
\end{theorem}
\begin{proof}
Character integration shows that $E_i^\circ$ has precisely its
original nonzero Fourier coefficients. The exact inverse of
$D_{v_t}$ proves the cancellation and smoothness. The kernel
and inverse fibre follow from
$E_i-\Pi E_i=-c_iD_{v_t}V$, with no condition on the freely
retained slow function $b=\Pi E_i$.

Haar preservation gives $(P_\Delta E_i)^\circ=P_\Delta E_i^\circ$.
Using the already proved inverse covariance gives
\[
\begin{aligned}
 \mathcal C_j(P_\Delta E_i)
 &=-c_j^{-1}T_g^{-\Delta}P_\Delta D_{v_t}^{-1}E_i^\circ\\
 &=-c_i^{-1}P_\Delta D_{v_t}^{-1}E_i^\circ
 =P_\Delta\mathcal C_i(E_i).
\end{aligned}
\]
The transfer theorem proves the descent test and its inverse.
The identity $P_aP_b=P_{a+b}$ proves the sum formula. If
$P_jF=P_kG$ and $j\le k$, injectivity of $P_j$ gives
$F=P_{k-j}G$; its inverse and full compatibility test are
again the transfer formulas. The calculation just made shows
that applying the correction preserves this equality.
Pullback preserves sums and products pointwise, proving the
last assertion for the original compatible coefficients.

The source's physical residual and velocity scaling gives
exactly
\[
 Q^{A_{\rm NS}}Q^{-(2A_{\rm NS}+1/2)}T_g^{-i}
 =Q^{-1-h_{\rm NS}}T_g^{-i}=c_i^{-1}.
\]
Thus the minus sign and coefficient above are the actual
chart representation of $-D_{v_t}^{-1}$ acting on the
physical residual in the absolute auxiliary variable, as
in source (8.20)--(8.23). A distinct $Q$ must retain its own
displayed scaling; the common-band identity holds at the
fixed physical normalization specified here.
\end{proof}

\begin{theorem}[Physical evaluation, chain rules, and its complete fibre]
\label{thm:ns-physical-evaluation-fibre}
Retain the source $d_r>0$, $D_{\rm NS}=1/2-h_{\rm NS}$,
$\varphi(r,t)=v_rr^{d_r}+v_tt\pmod{\mathbb Z^2}$, $r>0$.
On any open physical slow domain, define
\[
 \mathcal E_iF(r,z,t)=F(r,z,t,J^i\varphi(r,t))
\]
for $F\in C^\infty(U\times\mathbb T^2)$. This map is
surjective: its right inverse sends $f$ to the coefficient
$F(r,z,t,Y)=f(r,z,t)$ independent of $Y$. Its exact kernel
is the smooth functions vanishing on the displayed graph,
and its full fibre at $f$ is this right inverse plus that
kernel. For $i\ge j$, $\mathcal E_i=\mathcal E_jP_{i-j}$.
The derivatives are exactly
\[
\begin{aligned}
 \partial_t\mathcal E_iF
 &=\mathcal E_i(\partial_tF+T_g^iD_{v_t}F),\\
 \partial_r\mathcal E_iF
 &=\mathcal E_i(\partial_rF+d_rr^{d_r-1}\Lambda^iD_{v_r}F),\\
 \partial_z\mathcal E_iF&=\mathcal E_i\partial_zF.
\end{aligned}
\]
In the original chart $R=r/\sqrt Q$, $Z=z/Q^{D_{\rm NS}}$,
$T=(1-t)/Q$, these identities become
\[
\begin{aligned}
 Q^{1+h_{\rm NS}}\partial_t&\longleftrightarrow
       -\varepsilon\partial_T+c_iD_{v_t},\\
 \sqrt Q\partial_r&\longleftrightarrow
       \partial_R+\Lambda^iQ^{d_r/2}d_rR^{d_r-1}D_{v_r},\\
 \sqrt Q\partial_z&\longleftrightarrow\varepsilon\partial_Z.
\end{aligned}
\]
There is no linear map from evaluated physical coefficients
to slow functions whose composition with $\mathcal E_i$
equals the torus mean $\Pi$ on every extended coefficient.
Nevertheless the joint map $(\mathcal E_i,\Pi)$ is explicitly
surjective onto pairs of smooth slow functions, with the
full inverse fibres described in the proof.
\end{theorem}
\begin{proof}
Evaluation along a smooth graph for $r>0$ is smooth. The
displayed right inverse proves surjectivity; subtraction
proves the kernel and fibre. The band identity is direct
composition of the powers of $J$. The chain rule and
$J^iv_r=\Lambda^iv_r$, $J^iv_t=T_g^iv_t$ prove every
derivative identity. In chart coordinates, respectively
$\partial_tT=-1/Q$, $\partial_rR=Q^{-1/2}$ and
$\partial_zZ=Q^{-D_{\rm NS}}$. Substitution gives the
three stated coefficients, including the negative slow
time term. There is no assertion of regularity of
$r^{d_r}$ at the axis; the specified domain is $r>0$.

Fix $n\in\mathbb Z^2\setminus\{0\}$ and write
$a(r,t)=e_n(J^i\varphi(r,t))$, a nowhere zero smooth
function. The extended coefficient $H=e_n(Y)-a(r,t)$
has $\mathcal E_iH=0$ but $\Pi H=-a\ne0$.
This proves the claimed failure of a linear descent of
the mean through physical evaluation. It does not discard
the relation between the two data: for arbitrary desired
physical value $f$ and mean $b$, define
\[
 F_{f,b}(r,z,t,Y)=b(r,z,t)
       +(f(r,z,t)-b(r,z,t))\frac{e_n(Y)}{a(r,t)}.
\]
It satisfies $\mathcal E_iF_{f,b}=f$ and $\Pi F_{f,b}=b$.
The complete inverse fibre is
$F_{f,b}+(\ker\mathcal E_i\cap\ker\Pi)$, by subtraction.
For real $f,b$, the real part of this formula provides a
real right inverse as well. These exact fibres give a
structural justification for retaining auxiliary means
before evaluation: one evaluated coefficient leaves the
auxiliary mean free, with the complete correspondence
given by $F_{f,b}$ and its displayed fibre.
\end{proof}

\begin{theorem}[Sampling quotient, full original-frequency fibres, and inverse obstruction]
\label{thm:ns-reverse-sampling-fibres}
For $N\ge1$ let $D_N=\mathbb T[N]^2$ and let
$R_N:\mathcal P\to\mathbb C^{D_N}$ be restriction.
For $f=\sum_na_ne_n$ its discrete Fourier coefficient
at $r\in(\mathbb Z/N\mathbb Z)^2$ is exactly
\[
 \widehat {R_Nf}(r)=
 N^{-2}\sum_{x\in D_N}f(x)e^{-2\pi i r\cdot x}
 =\sum_{n\equiv r\pmod N}a_n.
\]
Thus $R_N$ is onto and its kernel consists of precisely
the finite coefficient families having zero sum in every
residue class. It satisfies $R_NP_m=J_N^{m*}R_N$,
where $J_N^{m*}q(x)=q(J^mx)$.

For any retained finite original spectrum $\Omega\subset\mathbb Z^2$,
let $\Omega_r=\{n\in\Omega:n\equiv r\pmod N\}$.
On $\mathcal P_\Omega=\operatorname{span}\{e_n:n\in\Omega\}$,
the image consists of grid functions whose coefficients
vanish at every empty $\Omega_r$. Above each such grid
function $q$, choose and retain one $n_r\in\Omega_r$
for each occupied class. All inverse fibres are
\[
 a_n=b_n\quad(n\ne n_r),\qquad
 a_{n_r}=\widehat q(r)-\sum_{n\in\Omega_r\setminus\{n_r\}}b_n,
\]
with every $b_n\in\mathbb C$ freely retained.
In particular restriction is injective on this domain
exactly when every occupied $\Omega_r$ is a singleton.

For either original $v=v_r,v_t$, a linear map on
$\operatorname{im}(R_N|_{\mathcal P_\Omega})$ intertwines
$R_ND_v$ with $R_N$ exactly when the elements of $\Omega$
have distinct residues modulo $N$. On the original zero-Haar-mean
subspace $\mathcal P_{\Omega,0}
=\operatorname{span}\{e_n:n\in\Omega\setminus\{0\}\}$,
an intertwining inverse map exists exactly when the elements
of $\Omega\setminus\{0\}$ have distinct residues. In each
case the map on the stated grid image is unique. The original
frequency zero may collide with a nonzero frequency without
violating the latter criterion, since the former is absent
from the inverse domain.

For either original $v=v_r,v_t$, no linear operator
$B_N$ on grid functions satisfies $R_ND_v=B_NR_N$
on all $\mathcal P$. No linear operator $A_N$ satisfies
$R_ND_v^{-1}=A_NR_N$ on all $\mathcal P_0$ either.
\end{theorem}
\begin{proof}
The product of the two finite geometric sums is one
exactly when both integer frequency differences are
divisible by $N$, and zero otherwise. Expanding $f$
proves the coefficient formula. Choosing one integer
representative of every residue class supplies an
interpolating polynomial for any grid function: the
$N^2$ discrete characters are orthonormal and hence a
basis of the $N^2$-dimensional function space. This
proves surjectivity and the kernel. Pointwise composition
proves the matrix intertwining. The same coefficient
equations restricted to $\Omega$ give the image and the
entire affine inverse fibre displayed above; the free
coordinates prove both directions of the injectivity test.

For the sharp restricted derivative criterion, suppose distinct
$n,n'\in\Omega$ have the same residue. Then
$R_N(e_n-e_{n'})=0$ but
\[
 R_ND_v(e_n-e_{n'})
   =2\pi i\,v\cdot(n-n')\,R_Ne_n\ne0.
\]
Indeed $n-n'$ is a nonzero integer vector and the nonvanishing
symbol statement applies; $R_Ne_n$ is nowhere zero. This
prevents descent of the derivative through restriction. If both
original frequencies are nonzero, the inverse instead gives
\[
 R_ND_v^{-1}(e_n-e_{n'})
  =\left(\frac1{2\pi i\,v\cdot n}
             -\frac1{2\pi i\,v\cdot n'}\right)R_Ne_n\ne0.
\]
The two denominators are nonzero and their difference is
$2\pi i\,v\cdot(n-n')\ne0$, so their reciprocals differ.
This proves necessity on precisely $\mathcal P_{\Omega,0}$.
Conversely, when the relevant residues are distinct,
restriction is a linear bijection onto its stated image.
Conjugation of the actual derivative or zero-mean inverse by
that bijection defines the required operator; both preserve
their original spectral domains. Surjectivity onto the image
forces uniqueness. These arguments include the empty domain,
whose grid image is the zero vector space.

Choose $k\in\mathbb Z^2\setminus\{0\}$. The polynomial
$e_{Nk}-1$ restricts to zero, whereas
\[
 R_ND_v(e_{Nk}-1)=2\pi iN(v\cdot k)\,1\ne0.
\]
The nonzero multiplier was proved for both original
directions. A linear map must send zero to zero, which
contradicts the proposed $B_N$. For the inverse take
$e_{Nk}-e_{2Nk}\in\mathcal P_0$. It again restricts to
zero, but its inverse restricts to
\[
 \left(\frac1{2\pi iN(v\cdot k)}
       -\frac1{4\pi iN(v\cdot k)}\right)1
 =\frac1{4\pi iN(v\cdot k)}\,1\ne0.
\]
This proves the second obstruction on exactly the stated
zero-Haar-mean domain. It concerns an unrestricted
sampling quotient, not the continuous inverse in (8.19).
\end{proof}

\begin{theorem}[Exact lifted computation of the source inverse and quadrature error]
\label{thm:ns-reverse-lifted-inverse}
Suppose the retained finite spectrum $\Omega$ has distinct
residues modulo $N$. Define the exact interpolation map
on $\operatorname{im}(R_N|_{\mathcal P_\Omega})$ by
\[
 I_\Omega q=\sum_{n\in\Omega}\widehat q([n])e_n.
\]
It is the inverse of $R_N$ on that domain. Thus
$R_ND_vI_\Omega$ is the exact original-frequency
derivative on the grid image; the inverse is
$R_ND_v^{-1}I_\Omega$ on the image of the original
zero-Haar-mean subspace. Zero mean means that the
coefficient of the original integer frequency $0$
vanishes; it does not mean the grid average vanishes.

For $m\ge0$ retain the full preimage
$H_{m,N}=(J^m)^{-1}(D_N)$ of size $14^mN^2$.
For $f=\sum_{n\in\Omega}a_ne_n$, every original coefficient
is recovered by
\[
 a_n=\frac1{14^mN^2}\sum_{Y\in H_{m,N}}
       (P_mf)(Y)e^{-2\pi i((J^m)^Tn)\cdot Y}.
\]
The resulting original source inverse is
\[
 D_v^{-1}P_mf(Y)
 =\sum_{\substack{n\in\Omega\\n\ne0}}
   \frac{a_n}{2\pi i\lambda^m(v\cdot n)}
      e_{(J^m)^Tn}(Y)
\]
whenever the original coefficient at zero is zero.
All covering labels are those of (T8); the inverse
frequency map is $(J^{-m})^T$ on its exact image.

For a smooth $f$ and real $s\ge0$ define, with the
original Fourier units,
\[
 \|f\|_{\mathcal A^s}
 =\sum_{n\in\mathbb Z^2}(1+4\pi^2|n|_2^2)^{s/2}|\widehat f(n)|.
\]
Then the full-preimage quadrature has the exact error
and the cover-level-independent bound
\[
\begin{aligned}
 \frac1{|H_{m,N}|}\sum_{Y\in H_{m,N}}P_mf(Y)-\Pi f
 &=\sum_{n\in N\mathbb Z^2\setminus\{0\}}\widehat f(n),\\
 \left|\frac1{|H_{m,N}|}\sum_{Y\in H_{m,N}}P_mf(Y)-\Pi f\right|
 &\le(1+4\pi^2N^2)^{-s/2}\|f\|_{\mathcal A^s}.
\end{aligned}
\]
\end{theorem}
\begin{proof}
Singleton residue fibres in the preceding theorem give
both interpolation compositions. Conjugating the actual
operators by these inverse maps proves the derivative
and inverse formulas on the precisely specified images.
For example, $\Omega=\{(N,0)\}$ consists of an original
nonzero frequency; its polynomials have zero Haar mean
although their grid values are constant. The original
frequency label makes its inverse unambiguous.

The product inside the lifted coefficient sum is
$P_m(fe_{-n})$. Each target of $D_N$ has $14^m$
preimages, so that sum equals its exact original-grid
average. The preceding discrete Fourier identity and
the distinct residue condition give $a_n$. Applying
the original multiplier to the recovered polynomial
proves the inverse formula, including the factor
$\lambda^m$ and the zero-mode exclusion. This computation
does not replace the original spectrum by its residue set.

For smooth $f$ the Fourier series is absolutely summable
with every polynomial weight. One elementary verification
integrates $(1-\partial_1^2-\partial_2^2)^kf$ against
$e_{-n}$ by parts; it gives
$|\widehat f(n)|\le C_k(1+4\pi^2|n|_2^2)^{-k}$.
Grouping the nonzero integer points by maximum coordinate
gives $8j$ points at radius $j$, so a sufficiently large
integer $k$ makes every stated weighted sum converge.
Uniform Fourier convergence follows; that the sum equals
$f$ follows also from the previously proved Fejer
approximation with the same coefficients.
We may therefore sum the exact finite-fibre identity
term by term. It retains exactly the frequencies in
$N\mathbb Z^2$, and removal of the original zero
frequency proves the error formula. Each remaining
frequency has Euclidean length at least $N$; multiplying
and dividing by the displayed weight and using the
triangle inequality proves the bound. There is no
unrecorded factor depending on $m$.
\end{proof}

The source's continuous mean and its temporal correction
therefore have an exact reverse connection to the repaired
full-fibre calculus. The restriction obstruction identifies
what a finite verification must retain; the lifted formula
provides a complete valid computation on its stated spectral
domain and a quantitative convergence statement on smooth
data. These proofs establish the specified operators and
correction maps. They do not substitute a sampled identity
for any remaining nonlinear estimate in the NS manuscript.

\subsection{A sharp inverse estimate returned to the manuscript's original directions}
\label{sec:ns-reverse-smooth}

This subsection applies the preceding coordinate-level operator calculation
back to the public manuscript's equation~(6.7), equation~(8.10), and
Lemma~8.6, equation~(8.19). In the preserved 166-page extraction these
occur on printed pages~64, 91, and 95, respectively. The manuscript's
inverse acts on the continuous torus with probability Haar measure.
We retain that domain and the original directions
\[
 \begin{gathered}
 v_r=(1,1-\sqrt2),\qquad v_t=(\sqrt2-1,1),\\
 D_v=v_1\partial_{Y_1}+v_2\partial_{Y_2},\qquad
 Y\in\mathbb R^2/\mathbb Z^2.
 \end{gathered}
\]
The following estimates strengthen the constants and derivative count
in those inverse steps. They require no PDE existence assertion.

\begin{theorem}[Sharp denominator constant in the original coordinates]
\label{thm:ns-sharp-denominator}
Put
\[
 \begin{gathered}
 c=\sqrt{4+2\sqrt2},\qquad d=\sqrt{4-2\sqrt2},\\
 w_r=(1,1+\sqrt2),\qquad w_t=(-1-\sqrt2,1).
 \end{gathered}
\]
For either paired choice $(v,w)=(v_r,w_r)$ or $(v_t,w_t)$,
$v\cdot w=0$, $\|v\|_2=d$, $\|w\|_2=c$, and for every
$k=(m,n)\in\mathbb Z^2\setminus\{0\}$,
\[
 |v\cdot k|>\frac1{c\|k\|_2},\qquad
 \inf_{k\ne0}|v\cdot k|\|k\|_2=\frac1c.             \tag{S1}
\]
In particular the constant $1/c$ cannot be increased.
\end{theorem}
\begin{proof}
Direct multiplication gives all dot products and squared lengths.
The products of the original symbols with their conjugates are
\[
 \begin{split}
 (v_r\cdot k)(w_r\cdot k)&=m^2+2mn-n^2=(m+n)^2-2n^2,\\
 (v_t\cdot k)(w_t\cdot k)&=n^2-2mn-m^2=(n-m)^2-2m^2.
 \end{split}                                                     \tag{S2}
\]
Each is a nonzero integer: its vanishing with a nonzero second
coordinate would make $\sqrt2$ rational, and when that coordinate
vanishes the first one must also vanish. Irrationality follows from
the parity argument given in Theorem~\ref{thm:ns-directional-intertwining}.
Their absolute values are therefore at least one. Cauchy--Schwarz
gives $|w\cdot k|\leq c\|k\|_2$. Equality would make the nonzero
integer vector $k$ parallel to $w$, whose coordinate ratio is
irrational, and is impossible. Hence
$1\leq|v\cdot k||w\cdot k|<c|v\cdot k|\|k\|_2$.

To prove optimality, retain the explicit Pell sequence
\[
 \begin{gathered}
 \alpha=1+\sqrt2,\qquad\beta=1-\sqrt2=-\alpha^{-1},\\
 p_0=1,\quad q_0=0,\quad
 p_{j+1}=p_j+2q_j,\quad q_{j+1}=p_j+q_j.
 \end{gathered}                                                   \tag{S3}
\]
Induction proves $p_j,q_j$ are integers and
$p_j+q_j\sqrt2=\alpha^j$, $p_j-q_j\sqrt2=\beta^j$.
In particular $p_j^2-2q_j^2=(-1)^j$. For $j\geq1$ set
\[
 k_{r,j}=(p_j-q_j,q_j),\qquad
 k_{t,j}=(q_j,q_j-p_j).
\]
The two vectors have the same length $K_j$. Their actual symbol
and conjugate values are respectively
\[
 \begin{array}{c|cc}
 &v\cdot k_{v,j}&w\cdot k_{v,j}\\\hline
 r&\beta^j&\alpha^j\\
 t&-\beta^j&-\alpha^j
 \end{array}
\]
by substitution. Since the retained nonzero perpendicular vectors
$v,w$ span $\mathbb R^2$, the exact decomposition is
$k=(v\cdot k)v/d^2+(w\cdot k)w/c^2$. Therefore
\[
 K_j^2=\frac{\alpha^{2j}}{c^2}+\frac{\alpha^{-2j}}{d^2},
 \qquad
 |v\cdot k_{v,j}|^2K_j^2
       =\frac1{c^2}+\frac{\alpha^{-4j}}{d^2}.          \tag{S4}
\]
The number $\alpha$ is greater than one, so $K_j$ tends to infinity
and the second expression tends to $1/c^2$. This proves (S1)'s
infimum and its nonattainment with every integer coordinate retained.
\end{proof}

We specify the analytic spaces used below. For $Y\in[0,1]^2$ write
\[
 \begin{gathered}
 e_k(Y)=e^{2\pi i k\cdot Y},\qquad W(k)=1+4\pi^2\|k\|_2^2,\\
 \widehat f(k)=\int_{[0,1]^2}f(Y)e^{-2\pi i k\cdot Y}\,dY,
 \end{gathered}
\]
For real $s$, $H^s$ is the completion of finite trigonometric
polynomials with the exact norm
\[
 \|f\|_{H^s}^2=\sum_{k\in\mathbb Z^2}W(k)^s|\widehat f(k)|^2,
 \qquad H^s_0=\{f\in H^s:\widehat f(0)=0\}.           \tag{S5}
\]
Thus elements may be represented by their weighted square-summable
coefficient arrays; truncation of such arrays converges in this norm
because the defining nonnegative series has vanishing tails.
For integer $m\geq0$ the coordinate norm on $C^m$ is
\[
 \|f\|_{C^m}=\max_{a,b\geq0,\ a+b\leq m}
          \sup_{Y\in\mathbb T^2}|\partial_{Y_1}^a\partial_{Y_2}^bf(Y)|.
                                                               \tag{S6}
\]
All integrations and suprema use the original period-one coordinates.

\begin{lemma}[Fourier reconstruction and the precise analytic estimates used]
\label{lem:ns-fourier-analytic-bridge}
For every smooth periodic function $f$, its Fourier series and each
coordinatewise differentiated series converge absolutely and uniformly
to $f$ and its corresponding derivative. For an integer $r\geq0$,
\[
 \|f\|_{H^r}\leq3^{r/2}\|f\|_{C^r}.                \tag{S7}
\]
The positive lattice constant
\[
 S=\sum_{k\in\mathbb Z^2\setminus\{0\}}W(k)^{-2}
       \quad\hbox{satisfies}\quad S\leq\frac3{4\pi^4}.            \tag{S8}
\]
For every zero-mean coefficient array $b$ and integer $m\geq0$,
\[
 \sum_{k\ne0}(2\pi\|k\|_2)^m|b(k)|
    \leq\sqrt S\left(\sum_{k\ne0}W(k)^{m+2}|b(k)|^2\right)^{1/2}.
                                                               \tag{S9}
\]
Whenever the right side is finite, the corresponding Fourier series
and its derivatives through order $m$ converge absolutely and
uniformly, with $C^m$ norm at most that right side.
\end{lemma}
\begin{proof}
The functions $e_k$ are orthonormal, since integrating each coordinate
character gives one at integer frequency zero and zero elsewhere.
For a finite set $F$ of frequencies, expansion of
$\|g-\sum_{k\in F}\widehat g(k)e_k\|_{L^2}^2\geq0$
therefore gives
$\sum_{k\in F}|\widehat g(k)|^2\leq\int|g|^2$.
Increasing $F$ proves Bessel's inequality for every continuous $g$.
Periodic integration by parts gives
$\widehat{\partial_1^a\partial_2^bf}(k)
 =(2\pi i k_1)^a(2\pi i k_2)^b\widehat f(k)$;
the boundary terms cancel at the two equal endpoint values.
Expanding $W(k)^r=(1+(2\pi k_1)^2+(2\pi k_2)^2)^r$
by the multinomial formula and applying Bessel to each derivative
with orders $a,b$, $a+b\leq r$, gives
\[
 \begin{split}
 \sum_kW(k)^r|\widehat f(k)|^2
 &\leq\sum_{j+a+b=r}\frac{r!}{j!a!b!}
                   \int_{[0,1]^2}|\partial_1^a\partial_2^bf|^2\\
 &\leq3^r\|f\|_{C^r}^2.
 \end{split}
\]
This proves (S7) without requiring an unproved Parseval identity.

For a nonzero frequency choose a coordinate with
$|k_i|\geq\|k\|_2/\sqrt2$. Integrating by parts $M$ times in
that coordinate yields
\[
 |\widehat f(k)|\leq
       \left(\frac{\sqrt2}{2\pi\|k\|_2}\right)^M\|f\|_{C^M}.
                                                               \tag{S10}
\]
Exactly $8q$ integer points have $\max(|k_1|,|k_2|)=q$.
Thus $\sum_{k\ne0}\|k\|_2^{-t}\leq8\sum_{q\geq1}q^{1-t}$
is finite for $t>2$. Formula (S10) with arbitrarily large $M$
proves absolute uniform convergence of the series and every
differentiated series. Its continuous sum has Fourier coefficients
$\widehat f(k)$ by uniform convergence and orthogonality.
The difference from $f$ has all coefficients zero. Its product
Fej\'er convolutions are zero and converge uniformly to that
difference by the explicit positive-kernel argument following
Theorem~\ref{thm:ns-torus-average}; hence the difference is zero.
This proves reconstruction and differentiation. The same argument
also yields Parseval for smooth functions if needed: finite Fourier
sums converge in $L^2$, and their squared norms are the finite sums
of squared coefficients by orthogonality.

The same $8q$ count gives
\[
 S\leq\sum_{q=1}^{\infty}\frac{8q}{(1+4\pi^2q^2)^2}
 \leq\frac1{2\pi^4}\sum_{q=1}^{\infty}q^{-3}
 \leq\frac1{2\pi^4}\left(1+\int_1^{\infty}t^{-3}\,dt\right)
 =\frac3{4\pi^4}.
\]
Finally $(2\pi\|k\|_2)^m\leq W(k)^{m/2}$.
Apply Cauchy--Schwarz to
$W(k)^{-1}$ and $W(k)^{(m+2)/2}|b(k)|$ to prove (S9).
For a derivative of order $a+b\leq m$, its multiplier is bounded
by $(2\pi\|k\|_2)^{a+b}\leq(2\pi\|k\|_2)^m$ because
$\|k\|_2\geq1$. The resulting summable majorant is uniform in
$Y$, proving all stated convergence and $C^m$ assertions.
\end{proof}

\begin{theorem}[The smooth inverse, exact Sobolev norm, and sharp derivative loss]
\label{thm:ns-sharp-smooth-inverse}
For either original direction $v$ and an integer $q\geq1$, define
\[
 \widehat{T_v^q f}(0)=0,\qquad
 \widehat{T_v^q f}(k)
       =\frac{\widehat f(k)}{(2\pi i\,v\cdot k)^q}\quad(k\ne0).
                                                               \tag{S11}
\]
For every real $s$, this is a bounded map
$T_v^q:H^{s+q}_0\to H^s_0$ with exact operator norm
\[
 \|T_v^q\|_{H^{s+q}_0\to H^s_0}
       =\left(\frac c{4\pi^2}\right)^q.              \tag{S12}
\]
For every $\delta<q$, this inverse has no bounded extension
$H^{s+\delta}_0\to H^s_0$ agreeing with (S11) on polynomials.

For smooth zero-mean $f$, $T_v^qf$ is the unique smooth zero-mean
solution of $D_v^q u=f$. More generally, for every datum
$f\in H^{s+q}_0$, its coefficients give the unique zero-mean
distributional solution. Here a distribution is a continuous
complex-linear functional on $C^\infty(\mathbb T^2)$ with its
seminorms (S6), and its coefficient at $k$ is its value on $e_{-k}$.
On smooth periodic
functions the kernel of $D_v^q$ is precisely the constants, its
image is all smooth zero-mean functions, and its full fibre at $f$
is $T_v^q f+\mathbb C$.
\end{theorem}
\begin{proof}
Theorem~\ref{thm:ns-sharp-denominator} gives the exact pointwise bound
\[
 \frac1{|2\pi i\,v\cdot k|^q W(k)^{q/2}}
 <\left(\frac{c\|k\|_2}{2\pi\sqrt{1+4\pi^2\|k\|_2^2}}\right)^q
 <\left(\frac c{4\pi^2}\right)^q.                   \tag{S13}
\]
Multiplying coefficient by coefficient and summing nonnegative
squares proves the upper bound in (S12) on polynomials and all
weighted square-summable arrays. Thus truncation extends the same
formula continuously to the stated completion, with no change of
coordinates. On the single character at $k_{v,j}$ of (S3)--(S4),
the ratio of the output norm to the input norm is
\[
 \frac{\alpha^{qj}}{(2\pi)^q(1+4\pi^2K_j^2)^{q/2}}
        \longrightarrow\left(\frac c{4\pi^2}\right)^q,
\]
because $K_j/\alpha^j\to1/c$ by (S4). This proves the exact
operator norm. If the input order is $s+\delta$ instead, the ratio
is
$\alpha^{qj}/((2\pi)^q(1+4\pi^2K_j^2)^{\delta/2})$.
Its quotient by $\alpha^{(q-\delta)j}$ tends to the positive
constant $c^\delta/(2\pi)^{q+\delta}$. It is unbounded when
$\delta<q$, proving sharpness of the derivative loss.

For smooth $f$, (S10) and the multiplier growth
$|(2\pi i\,v\cdot k)^{-q}|\leq(c\|k\|_2/(2\pi))^q$
give uniform absolute convergence of every differentiated inverse
series. Lemma~\ref{lem:ns-fourier-analytic-bridge} and termwise
differentiation therefore show smoothness and $D_v^qT_v^qf=f$.
The zero coefficient is exactly its Haar mean. Every nonzero
Fourier multiplier of $D_v^q$ is nonzero, so any smooth solution
has exactly the coefficients (S11) away from zero. Its remaining
constant is arbitrary; Fourier reconstruction proves the full
kernel and fibres.

To make the distributional assertion explicit, the resulting
$u\in H^s_0$ has coefficients of at most polynomial growth:
each term of its squared norm is bounded by that full norm, so
$|\widehat u(k)|\leq\|u\|_{H^s}W(k)^{-s/2}$.
Consequently there are $B\geq0$ and an integer $r\geq0$ with
$|\widehat u(k)|\leq B\|k\|_2^r$ for $k\ne0$;
one obtains them directly from
$W(k)\leq(1+4\pi^2)\|k\|_2^2$ and $W(k)\geq1$.
Define its pairing with a smooth periodic test function by
$u(\varphi)=\sum_k\widehat u(k)\widehat\varphi(-k)$.
Choose an integer $M>r+2$. Formula (S10) and the lattice count
give the explicit continuous-functional bound
\[
 |u(\varphi)|\leq
 \left(|\widehat u(0)|+
    8B\left(\frac{\sqrt2}{2\pi}\right)^M
       \sum_{j=1}^{\infty}j^{1+r-M}\right)\|\varphi\|_{C^M}.
\]
The series converges, so this defines a distribution with exactly
the specified coefficients. The same construction realizes
$f\in H^{s+q}_0$ as a distribution. Distributional differentiation
is defined by $(D_vu)(\varphi)=-u(D_v\varphi)$; evaluating on
$e_{-k}$ therefore multiplies its coefficient by $2\pi i\,v\cdot k$.
Formula (S11) gives the coefficient equation for $D_v^qu=f$.
Finally a distribution with all Fourier coefficients zero vanishes
on every trigonometric polynomial. The Fourier partial sums of
any smooth test function converge to it in every seminorm (S6),
by Lemma~\ref{lem:ns-fourier-analytic-bridge}. Continuity of the
distribution therefore makes its value on every smooth test
function zero. This proves the coefficient equation as an equality
of distributions and proves uniqueness: any other zero-mean
solution has the same forced nonzero coefficients and the same
zero coefficient.
\end{proof}

\begin{theorem}[A strengthened coordinate derivative estimate]
\label{thm:ns-cm-improvement}
For integers $m\geq0$, $q\geq1$ and zero-mean
$f\in C^{m+q+2}(\mathbb T^2)$, the inverse series (S11) belongs
to $C^m$ and satisfies
\[
 \begin{aligned}
 \|T_v^qf\|_{C^m}
   &\leq\left(\frac c{4\pi^2}\right)^q
             \sqrt S\,\|f\|_{H^{m+q+2}}\\
   &\leq C_{m,q}\|f\|_{C^{m+q+2}}.
 \end{aligned}                                                   \tag{S14}
\]
where the completely specified constant is
\[
 C_{m,q}=\left(\frac c{4\pi^2}\right)^q
          3^{(m+q+2)/2}\sqrt S
 \leq\left(\frac c{4\pi^2}\right)^q
          \frac{3^{(m+q+3)/2}}{2\pi^2}.             \tag{S15}
\]
For $q=1$ this uses $m+3$ input derivatives and applies to the
original temporal inverse in equation~(8.19). For general $q$ it
uses $m+q+2$ and applies to the original repeated radial inverse
in equation~(8.10). No optimality assertion is made for this
particular $C^m$ derivative count.
\end{theorem}
\begin{proof}
Apply (S9) to the inverse coefficients, and then (S13) with
$s=m+2$. This gives the first inequality (S14), including uniform
convergence of all inverse derivatives through order $m$.
Formula (S7) with $r=m+q+2$ gives the second inequality, and
(S8) gives the displayed bound on the constant. The proofs of
(S7) and (S9) only require the finite number of continuous
derivatives displayed here, so the same argument applies directly
to $C^{m+q+2}$ data. In the two source equations the operators
are exactly $T_{v_t}$ and $T_{v_r}^q$, with the same Fourier
constant $2\pi i$, sign, period and zero mode. The source bounds
with $m+4$ and $m+q+3$ input derivatives follow at once from
(S14), since increasing $C^r$ order only increases the norm (S6).
\end{proof}

\begin{proposition}[Parameter derivatives, support, reality and the source mean update]
\label{prop:ns-inverse-parameters}
Let $V\subset\mathbb R^b$ be open, and let
$f:V\times\mathbb T^2\to\mathbb C$ be jointly smooth with
$\int_{\mathbb T^2}f(x,Y)\,dY=0$ for every $x\in V$.
Define $u(x,Y)=T_v^q(f(x,\cdot))(Y)$. Then $u$ is jointly
smooth and, for every parameter multi-index $\gamma$,
\[
 \begin{gathered}
 \partial_x^\gamma u=T_v^q(\partial_x^\gamma f),\\
 \sup_{x\in K}\|\partial_x^\gamma u(x,\cdot)\|_{C^m}
 \leq C_{m,q}\sup_{x\in K}
                \|\partial_x^\gamma f(x,\cdot)\|_{C^{m+q+2}}.
 \end{gathered}                                                  \tag{S16}
\]
for every compact $K\subset V$. If $f(x,Y)=0$ for all $Y$
whenever $x\notin A\subset V$, the same statement holds for $u$.
Real-valued data give a real-valued inverse.

In the manuscript's fixed local band/common-torus chart, retain
$Q>0$, $h=h_{\rm NS}$, $c_{i_0}=T_g^{i_0}Q^{1+h}>0$,
$T_g=4+\sqrt2$, and the distinct coordinates
\[
 \begin{gathered}
 Y\in\mathbb T^2_{\rm abs},\qquad y=J^{i_0}Y\in\mathbb T^2_{\rm com},
 \qquad P_{i_0}F(Y)=F(J^{i_0}Y),\\
 D_v^Y=v\cdot\partial_Y,\qquad D_v^y=v\cdot\partial_y,\qquad
 T_v^Y=(D_v^Y)^{-1},\qquad T_v^y=(D_v^y)^{-1}.
 \end{gathered}
\]
The inverse in each coordinate acts on that torus's smooth
zero-Haar-mean functions and uses its own original Fourier
frequency with the multiplier (S11). Denote the corresponding
means by $\Pi_Y$ and $\Pi_y$. On the indicated smooth domains,
\[
 \begin{gathered}
 \Pi_Y P_{i_0}=\Pi_y,\qquad
 D_{v_t}^Y P_{i_0}=T_g^{i_0}P_{i_0}D_{v_t}^y,\\
 T_{v_t}^Y P_{i_0}=T_g^{-i_0}P_{i_0}T_{v_t}^y
       \quad\hbox{on }\ker\Pi_y.
 \end{gathered}
\]
For the actual common-chart residual write
$E_\theta^\circ=E_\theta-\Pi_yE_\theta$. Its temporal update is
\[
 \begin{gathered}
 \Delta v=-c_{i_0}^{-1}T_{v_t}^yE_\theta^\circ,\qquad
 c_{i_0}D_{v_t}^y\Delta v=-E_\theta^\circ,\\
 P_{i_0}\Delta v=-Q^{-1-h}T_{v_t}^Y(P_{i_0}E_\theta^\circ),\\
 \|\Delta v\|_{C_y^m}
       \leq c_{i_0}^{-1}C_{m,1}\|E_\theta^\circ\|_{C_y^{m+3}}.
 \end{gathered}
                                                               \tag{S17}
\]
The identical typed formulas apply to the manuscript's desired axial
increment $\gamma_d=-c_{i_0}^{-1}T_{v_t}^yE_z^\circ$.
For smooth shell-supported data the source's actual axial increment is
$\Delta\gamma=\gamma_d+a$, with $a=-A_1\gamma_d$, where its
retained cutoff operator is explicitly
\[
 \begin{split}
 (A_1f)(R,Z,T,y)
  ={}&R^{-1}(\partial_R\chi_m)(R,Z,T)\\
  &{}\cdot\int_0^\infty R'f\bigl(R',Z,T,
          y+M((R')^{d_r}-R^{d_r})v_r\bigr)\,dR',
 \qquad M=\Lambda^{i_0}Q^{d_r/2}.
 \end{split}
\]
Here $R>0$, the smooth zero extension of $f$ has radial support
locally bounded in a compact subinterval of $(0,\infty)$, and
$\chi_m(R,Z,T)$ is the source's smooth torus-independent cutoff.
The actual fast-time cancellation is consequently
\[
 c_{i_0}D_{v_t}^y\Delta\gamma
       =-E_z^\circ+c_{i_0}D_{v_t}^y a.
\]
The complete displayed time operator remains
$-\varepsilon\partial_T+c_{i_0}D_{v_t}^y$, with
$\varepsilon=Q^h$ and $\partial_T$ holding $y$ fixed. Thus its
action on the updates still contains
$-\varepsilon\partial_T\Delta v$ and
$-\varepsilon\partial_T\Delta\gamma$, respectively.
No zero or flatness assertion about $a$ is part of this inverse estimate.
\end{proposition}
\begin{proof}
On a compact parameter neighborhood inside $V$, every mixed
parameter/torus derivative of $f$ is bounded. Differentiation of
its coefficient integral in a parameter is justified by the
continuous bounded derivative on that compact set times the
finite torus measure. Repetition gives
$\partial_x^\gamma\widehat f(x,k)
=\widehat{\partial_x^\gamma f}(x,k)$ and also zero mean for
each parameter derivative. Integration by parts as in (S10)
gives every polynomial decay order uniformly on this neighborhood,
for all parameter derivatives under consideration. After multiplication
by the inverse multiplier and any fixed coordinate derivative,
choose the decay order larger than the total growth exponent plus
two. The $8q$ lattice-shell count then supplies a summable majorant
independent of the parameter and torus variables. Thus every mixed
series differentiates termwise and converges locally uniformly.
One may justify each first derivative directly by integrating its
uniformly convergent derivative series along a coordinate segment;
the fundamental theorem of calculus passes to the sum. Iterating
proves joint smoothness and (S16)'s exact commutation. Applying
(S14) at each parameter and taking suprema proves its bound.

When a parameter slice of $f$ vanishes, all its Fourier coefficients
vanish, and (S11) gives the vanishing slice of $u$. This proves the
support assertion without changing the parameter set. If $f$ is
real, $\widehat f(-k)=\overline{\widehat f(k)}$ and
$(2\pi i\,v\cdot(-k))^{-q}
=\overline{(2\pi i\,v\cdot k)^{-q}}$; pairing opposite modes
shows $u$ is real. A common-torus character $e_k(y)$ pulls back
to $e_{(J^{i_0})^Tk}(Y)$. Its absolute derivative multiplier is
$2\pi i\,v_t\cdot(J^{i_0})^Tk
=T_g^{i_0}2\pi i\,v_t\cdot k$.
The lattice map is injective and preserves zero frequency exactly,
which proves the stated Haar identity and both typed intertwining
identities, including the reciprocal factor for the inverse. Smooth
Fourier convergence proved above passes them from characters to the
specified smooth spaces. In particular
$P_{i_0}T_{v_t}^y=T_g^{i_0}T_{v_t}^YP_{i_0}$ on $\ker\Pi_y$.
Multiplying this identity by
$-c_{i_0}^{-1}=-T_g^{-i_0}Q^{-1-h}$ gives the absolute
representative in (S17) with its complete coefficient.

The equation $D_{v_t}^yT_{v_t}^yE_\theta^\circ=E_\theta^\circ$
proves (S17)'s common-chart cancellation, including its minus sign.
The estimate is (S14), now in the explicitly named common
coordinates, multiplied by $c_{i_0}^{-1}$. The same calculation
with $E_z^\circ$ gives the desired $\gamma_d$ and its bound.
For the stated shell-supported inputs, differentiation under the
displayed $A_1$ integral is justified locally by a fixed compact
radial interval and bounded smooth derivatives on that interval;
$R$ stays positive. This proves smoothness of $a$ there.
Linearity of $D_{v_t}^y$ applied to
$\Delta\gamma=\gamma_d+a$ proves the actual fast cancellation
with the retained term $c_{i_0}D_{v_t}^ya$. Finally applying the
complete time operator to each update gives exactly the displayed
slow derivatives in addition to its fast part. No slow derivative
has been commuted through the radial cutoff operator.
\end{proof}

The static source dependency for this result is explicit.
\texttt{DiophantineGraph.lean} proves the algebraic norm lower bounds
at lines~94 and~104 and the reciprocal bounds with constant~6 at
lines~226 and~241. \texttt{TorusInverse.lean} retains
$\omega=2\pi i$ at line~52, defines the same original vectors at
line~195, and proves the inverse coefficient bound and smooth
inverse at lines~278 and~323. Its Haar measure at line~342 is the
product of two actual circle Haar probability measures.
\texttt{SmoothFourierData.lean}, lines~351 and~561, proves the
rapid-coefficient and reconstruction bridge for smooth periodic
functions; its line~565 returns the directional equation for
those actual functions. Thus rapid coefficients are not being
inserted here as an unproved surrogate for a smooth function.
Equations (S1), (S12), and (S14) provide sharper analytical constants
and derivative estimates for that same source chain. These are
independent mathematical proofs with a bounded static source audit;
no Lean process or complete PDE verification is asserted.


\begin{thebibliography}{99}
\raggedright
\bibitem{UmbrellaContribution}
Reddit user \texttt{u/UmbrellaCorp\_HR},
\emph{Arithmetic, positivity-code and prime-production contributions to the
Erd\H{o}s--Straus programme}, private communications relayed by the project
author, September 2026, including the continuation numbered (71)--(100).
The contributor's public handle was explicitly supplied by the project author.
No unreceived manuscript is cited as inspected.

\bibitem{PrimewiseCoordinates}
The Clankers, \emph{The complete primewise coordinate system},
mathematical working source, 8 September 2026, section
\emph{Exact unary counts and simultaneous logical constraints}.
This intermediate derivation develops the arithmetic input credited
in \cite{UmbrellaContribution}. The bounded code and witness equations
needed here are reproduced and proved in this article.

\bibitem{PositivityTransport}
The Clankers, \emph{Transport into the colleague's stronger finite-code
positivity layers}, mathematical working source, 8 September 2026,
equations (PCT1)--(PCT11); and \emph{Fixed TCHISLA rays and their exact
intersection with fixed divisor rays}, section
\emph{The original factor scale and its complete coordinate fibre}.
These intermediate calculations retain the original colleague's contribution
\cite{UmbrellaContribution}; their retained-scale and code maps used here
are completely reproved.

\bibitem{SignedMarginAntecedent}
The Clankers, \emph{Signed margins, exact witness swaps, and the
prime-output sieve}, mathematical working source, 8 September 2026,
Theorem \texttt{thm:margin-swap-bands}.
The exact antecedent SHA-256 begins \texttt{d51347ad95517c59};
the complete digest is recorded at the start of
Section~\ref{sec:boundary-four-candidates}.
That source supplies the same-tag exchange antecedent; all four
candidate maps and proofs needed here are written out in this article.

\bibitem{ElsholtzTao}
Christian Elsholtz and Terence Tao,
\emph{Counting the number of solutions to the Erd\H{o}s--Straus equation
on unit fractions}, arXiv:1107.1010 (2011), Section 2.
\href{https://arxiv.org/abs/1107.1010}{Primary preprint}.
The Type I/II parametrization lineage is relevant; this is not attribution
of the later finite-code or auxiliary-torus results to that paper.

\bibitem{CCTrace}
Alain Connes and Caterina Consani,
\emph{Hochschild homology, trace map and zeta-cycles},
in \emph{Cyclic Cohomology at 40: Achievements and Future Prospects},
Proceedings of Symposia in Pure Mathematics 105 (2023), pp.~83--102.
DOI: \href{https://doi.org/10.1090/pspum/105/01896}{10.1090/pspum/105/01896}.
The source loci used are Proposition 3.1 and equation (3.2); the finite
integral constructions in this article are proved separately.

\bibitem{CCPericyclic}
Alain Connes and Caterina Consani,
\emph{Cyclic theory and the pericyclic category},
in the same volume, pp.~103--122.
DOI: \href{https://doi.org/10.1090/pspum/105/01897}{10.1090/pspum/105/01897}.
Section 7.4, Lemma 7.7, p.~120, and Proposition 7.8, p.~121.

\bibitem{OpenAINS}
OpenAI, \emph{Finite Time Blowup for Navier--Stokes}, 8 September 2026,
\href{https://cdn.openai.com/pdf/32d9f210-8b73-45e0-91bc-82a30aef8a9a/navier-stokes.pdf}
{public manuscript}. The 166-page version used for the reverse calculations
has SHA-256
\texttt{0e779481c4da40bd28d1e642e1d8ca5744}\allowbreak
\texttt{7d129610df28dfa5a11e9af8ae228f}.
The prior 165-page version is separately identified in the source archive.
Citation of the source is not a claim of complete independent validation.

\bibitem{OpenAISource}
OpenAI, \emph{NavierStokesAndEuler}, source revision
\href{https://github.com/openai/NavierStokesAndEuler/tree/8937a8f4cbc7abaab5e9e97d1cc7f5d2319d9538}
{\texttt{8937a8f4cbc7abaab5e9e97d1cc7f5d2319d9538}}.
The exact declarations read are identified in the operator proofs.
No Lean build for the full fluid theorem was run for this contribution.

\bibitem{BuckmasterStatement}
Tristan Buckmaster, \emph{Statement}, 8 September 2026,
\href{https://cims.nyu.edu/~tristanb/statement.pdf}{author-hosted statement},
p.~1 for the forced-fluid programme's intellectual lineage and the joint
work with Levent Alp\"oge. The statement's claims retain their author's scope.

\bibitem{CMZEuler}
Diego C\'ordoba and Luis Mart\'inez-Zoroa,
\emph{Blow-up for the incompressible 3D-Euler equations with uniform
$C^{1,\frac12-\epsilon}\cap L^2$ force}, arXiv:2309.08495v1 (2023).
\href{https://arxiv.org/abs/2309.08495v1}{Primary preprint record}.
Cited for the preceding forced-fluid programme, not as a theorem dependency
of the arithmetic claims in this article.

\bibitem{CMZIPM}
Diego C\'ordoba and Luis Mart\'inez-Zoroa,
\emph{Finite time singularities of smooth solutions for the 2D incompressible
porous media (IPM) equation with a smooth source}, arXiv:2410.22920v3 (2025).
\href{https://arxiv.org/abs/2410.22920v3}{Primary preprint record}.
The full proof is not independently audited in this contribution.

\bibitem{CMZHypo}
Diego C\'ordoba, Luis Mart\'inez-Zoroa and Fan Zheng,
\emph{Finite time blow-up for the hypodissipative Navier Stokes equations
with a force in $L^1_t C_x^{1,\epsilon}\cap L^\infty_tL_x^2$},
arXiv:2407.06776v2.
\href{https://arxiv.org/abs/2407.06776v2}{Primary preprint record}.
Cited with the source's precise hypodissipative, not full-viscosity, scope.

\bibitem{ABCIPM}
Levent Alp\"oge, Tristan Buckmaster and Matei P. Coiculescu,
\emph{IPM manuscript}, September 2026,
\href{https://cims.nyu.edu/~tristanb/ipm.pdf}{author-hosted manuscript},
title page and Theorem 2.1. The listed authors are those on the inspected
title page; this is a separate fluid result.

\bibitem{ABBoussinesq}
Levent Alp\"oge and Tristan Buckmaster,
\emph{Boussinesq manuscript}, September 2026,
\href{https://cims.nyu.edu/~tristanb/boussinesq.pdf}{author-hosted manuscript},
title page, Theorem 1.1 and Lemma 3.1.
Earlier Boussinesq work by C\'ordoba, Andr\'es La\'in-Sanclemente and
Mart\'inez-Zoroa is credited by that source.

\bibitem{PriorEdition}
The Clankers, \emph{Erd\H{o}s--Straus Project Archive},
version 2026.09.08.1, the 58-statement supplement.
DOI: \href{https://doi.org/10.5281/zenodo.22666493}{10.5281/zenodo.22666493}.
This citation identifies the earlier immutable edition, not a replacement
for the human sources credited above.
\end{thebibliography}

\end{document}
